% Options for packages loaded elsewhere
\PassOptionsToPackage{unicode}{hyperref}
\PassOptionsToPackage{hyphens}{url}
\PassOptionsToPackage{dvipsnames,svgnames,x11names}{xcolor}
\documentclass[
  a4paper,
]{article}
\usepackage{xcolor}
\usepackage[margin=25mm]{geometry}
\usepackage{amsmath,amssymb}
\setcounter{secnumdepth}{-\maxdimen} % remove section numbering
\usepackage{iftex}
\ifPDFTeX
  \usepackage[T1]{fontenc}
  \usepackage[utf8]{inputenc}
  \usepackage{textcomp} % provide euro and other symbols
\else % if luatex or xetex
  \usepackage{unicode-math} % this also loads fontspec
  \defaultfontfeatures{Scale=MatchLowercase}
  \defaultfontfeatures[\rmfamily]{Ligatures=TeX,Scale=1}
\fi
\usepackage{lmodern}
\ifPDFTeX\else
  % xetex/luatex font selection
  \setmainfont[]{DejaVu Serif}
  \setsansfont[]{DejaVu Sans}
  \setmonofont[]{DejaVu Sans Mono}
\fi
% Use upquote if available, for straight quotes in verbatim environments
\IfFileExists{upquote.sty}{\usepackage{upquote}}{}
\IfFileExists{microtype.sty}{% use microtype if available
  \usepackage[]{microtype}
  \UseMicrotypeSet[protrusion]{basicmath} % disable protrusion for tt fonts
}{}
\makeatletter
\@ifundefined{KOMAClassName}{% if non-KOMA class
  \IfFileExists{parskip.sty}{%
    \usepackage{parskip}
  }{% else
    \setlength{\parindent}{0pt}
    \setlength{\parskip}{6pt plus 2pt minus 1pt}}
}{% if KOMA class
  \KOMAoptions{parskip=half}}
\makeatother
\usepackage{longtable,booktabs,array}
\newcounter{none} % for unnumbered tables
\usepackage{calc} % for calculating minipage widths
% Correct order of tables after \paragraph or \subparagraph
\usepackage{etoolbox}
\makeatletter
\patchcmd\longtable{\par}{\if@noskipsec\mbox{}\fi\par}{}{}
\makeatother
% Allow footnotes in longtable head/foot
\IfFileExists{footnotehyper.sty}{\usepackage{footnotehyper}}{\usepackage{footnote}}
\makesavenoteenv{longtable}
\usepackage{graphicx}
\makeatletter
\newsavebox\pandoc@box
\newcommand*\pandocbounded[1]{% scales image to fit in text height/width
  \sbox\pandoc@box{#1}%
  \Gscale@div\@tempa{\textheight}{\dimexpr\ht\pandoc@box+\dp\pandoc@box\relax}%
  \Gscale@div\@tempb{\linewidth}{\wd\pandoc@box}%
  \ifdim\@tempb\p@<\@tempa\p@\let\@tempa\@tempb\fi% select the smaller of both
  \ifdim\@tempa\p@<\p@\scalebox{\@tempa}{\usebox\pandoc@box}%
  \else\usebox{\pandoc@box}%
  \fi%
}
% Set default figure placement to htbp
\def\fps@figure{htbp}
\makeatother
\setlength{\emergencystretch}{3em} % prevent overfull lines
\providecommand{\tightlist}{%
  \setlength{\itemsep}{0pt}\setlength{\parskip}{0pt}}
\usepackage{luatexja}
\usepackage{xurl}
\usepackage{bookmark}
\IfFileExists{xurl.sty}{\usepackage{xurl}}{} % add URL line breaks if available
\urlstyle{same}
\hypersetup{
  pdftitle={Reading Out Time{,} Space{,} Mass{,} Charge{,} Translational Velocity and Spin from the Relations and Interference of Complex Waves without a Background Spacetime},
  pdfauthor={Noriaki Kihara (WF System Co., Ltd.) ORCID 0009-0004-6753-4020},
  colorlinks=true,
  linkcolor={Maroon},
  filecolor={Maroon},
  citecolor={Blue},
  urlcolor={Blue},
  pdfcreator={LaTeX via pandoc}}

\title{Reading Out Time, Space, Mass, Charge, Translational Velocity and
Spin from the Relations and Interference of Complex Waves without a
Background Spacetime}
\usepackage{etoolbox}
\makeatletter
\providecommand{\subtitle}[1]{% add subtitle to \maketitle
  \apptocmd{\@title}{\par {\large #1 \par}}{}{}
}
\makeatother
\subtitle{― An Integrated Hypothesis in the Order of the Thought
Experiment, from Residents of a Closed Circle to Residents of a
Four-Dimensional Spherical Shell, and on to α, the Four-Dimensional
Lattice and E₈ ―}
\author{Noriaki Kihara (WF System Co., Ltd.) ORCID 0009-0004-6753-4020}
\date{October 15, 2026 Version DOI 10.5281/zenodo.23384929}

\ltjsetparameter{jacharrange={-2,-3}}
\begin{document}
\maketitle

\textbf{Author:} Noriaki Kihara\\
\textbf{ORCID:} 0009-0004-6753-4020\\
\textbf{Version DOI:} 10.5281/zenodo.23384929\\
\textbf{Concept DOI:} 10.5281/zenodo.23384928\\
\textbf{Version:} v1.0 (published version. v0.1, which made α the
subject, was rejected by the author; v0.2 was rewritten in full along
the order of the author's questions; v1.0 tidied the quotations and
fixed the figures)\\
\textbf{Date:} 2026-10-15\\
\textbf{Related papers:} Wave model of closed-sphere residents {[}H1{]},
Fourteenth Thought Experiment {[}A1{]}, Thirteenth Thought Experiment
{[}D1{]}, Eleventh Thought Experiment {[}D2{]}, Twelfth Thought
Experiment {[}F1{]}, finite-order resonance {[}B1{]}\\
\textbf{Classification:} Hypothesis paper following the order of a
thought experiment / an alternative interpretive hypothesis of the
standard theory / presentation of tasks that can be taken over and
verified\\
\textbf{Method of record:} The primary sources are the full texts of two
dialogues between the author and Claude (``Intrinsic geometry on a
circle'', 2026-10-13 to 14; ``The dialogue after the Fourteenth Thought
Experiment'', 2026-10-14 to 15). The order of the author's questions and
corrections is used as the order of the chapters. Claude's errors and
the author's corrections are kept as part of the derivation. All
numerical experiments were re-run and checked.\\
\textbf{Keywords:} closed-sphere residents, imaginary amplitude, true
squares, Σxₙ² = R², light cone, gauge interval and geodesic length,
pinned phase (Higgs-like background), cross term, proper time, velocity
addition, necessity of three dimensions, spin, Hopf-fibre winding,
observation as the cross term, four-dimensional Hurwitz lattice, 23/124,
three 137s, near but not locked, complete absorbing shell at infinity,
E₈ gluing

\begin{center}\rule{0.5\linewidth}{0.5pt}\end{center}

\subsection{Abstract}\label{abstract}

This paper makes one claim. \textbf{Without placing a background
spacetime, and from nothing but the relations and interference of
complex waves that are not even travelling waves of odd and even
harmonics, time, space, mass, charge, translational velocity and spin
can all be read out.} What is placed is only a phase θ and an
imaginary-direction amplitude A on a closed circle (S¹), and the single
condition that the sum of the true squares (squares without conjugation)
of the components equals R², Σxₙ² = R². Written as Σxₙ² + (iR)² = 0 this
is the light-cone condition, and R is not a radius placed outside but a
quantity of the same kind as the components (an amplitude).

The derivation proceeds in the order of the author's questions and
corrections. A resident of the circle has no direction for a transverse
wave, so the amplitude is placed in the imaginary direction. The number
n of nodes (phases where the amplitude vanishes) counts
half-wavelengths, so odd n is allowed; odd n closes in two laps and even
n in one. Summing harmonics of equal amplitude localises them at one
point; the appearance of localisation and the fixing of the origin θ₀
are one event, and the fundamental component n = 1 or n = 2 acts as the
gauge and carries a virtual origin. In two dimensions R does not enter.
Extending the residents to the four-dimensional spherical shell (S³), R
enters and θ₀ becomes three (two phases and one mixing angle). A
numerical experiment that solved a self-consistent shell with absolute
squares and a real radius was fixed as an error by the author's
corrections ``true squares'', ``not a norm'', ``there is no outside of
the circle'', and its results are not adopted. Residents cannot read
curvature, but the geodesic length per gauge interval becomes
non-uniform (the peak interval is 1.245 times longer for all orders N =
18). In slices of S³, even components raise a peak of the same sign at
the antipode and odd components one of the opposite sign, and only with
both present is the position fixed at one point.

Then the author's eleven stages follow. (1) Mass: from the question that
a travelling wave cannot be represented, through the pinned phase (which
corresponds to the Higgs field and does not oscillate) whose beat with
the wave reads out mass, time and space at once, through ``the speed of
light is not exceeded because the gauge is widened'', to the realisation
that the way of thinking of the Lorentz factor is the light cone Σxₙ² =
0 itself; on the circle v² + (1/γ)² = 1 holds as an identity (residual
10⁻¹⁵), v = tanh ln(a/b), the velocity-addition law agrees to six
digits, and mass 1/γ = m/R is read out as the fraction of the cross term
2ab of the counter-rotating terms. (2) Three dimensions: the only sphere
satisfying the three conditions---uniform translations form a group (S¹,
S³), propagation without a tail (odd dimensions), and a mass gap from
curvature the residents can read (all but S¹)---is S³. (10) Spin: the
odd/even distinction is spin 1/2 and integer spin as it stands, the
cross term has spin 0, localisation at one point and a spin eigenstate
are incompatible, a spin eigenstate is a two-point localisation, and the
opposite spin on the far side can be observed by no one. (3) Charge: the
Hopf-fibre winding number, whose integrality gives quantisation.
Observation is the cross term between particle and observer, and being
able to read, radiating, and distributing energy are the same event. (4)
Can it be read as the electron: from the questions about the weak force
(redistribution of harmonic components) and the strong force, through
Koide's relation being a zero square sum, the electron is chosen as the
testing ground (E1). (5) When α is introduced, Claude's reasoning about
``why 23/124'' is dismissed by the author; the problem is ``why does it
lock to that window''. (9) The four-dimensional lattice: only in four
dimensions, where the diagonal is an integer, do Z₄ and Z₆ coexist in
one lattice, and this is why the background is four-dimensional and the
visible space three-dimensional (E2). (6) A detour: Planck's constant
was placed at its SI value to judge the electron mass (E3). (7) In G = c
= q₀ = 1 (q₀ = e/3), ħ = 9α⁻¹ is an output; once α is fixed the electron
mass necessarily agrees, and α alone is what must be verified. (8) α is
obtained from 23/124 through N(R) = 4π/(1−R)², and the four-dimensional
lattice makes it easy to pull in there (E4). The observed value deviates
slightly because observation is an average that includes states in
transition, and the electron itself distorts the gauge (E5). The
harmonic content, in a system of ratios with no absolute scale, is fixed
by the degree of localisation, and the minimal structure is two spin
components × three three-phase components. Localisation occurs by
exchange and returns after one cycle; the only self-consistent lock is
23/124; the dynamics is inherent in the system as the record on the
complete absorbing shell at infinity, and the answer table of the
Eleventh Thought Experiment allows the answers to be checked. The
beginning of the universe (geometric expansion of an unstable
self-consistent closure of a highly symmetric monochromatic wave), the
places of the three forces, the inflaton and the Higgs, and the
background radiation line up as stages of the same system. (11) If the
Lamb shift and the g-factor cannot be derived, the E₈ that has not yet
been used supplies the constraint. E₈ is two Hurwitz quaternions glued
by ``their difference is a multiple of (1+i)''; the point 180° opposite
enters automatically; and since the way of connection has been
identified, what remains is a derivation problem.

α, 137 and 23/124 are one stage of this development, not the subject of
the paper. The paper separates what has been verified, what has been
connected as hypothesis, and the means of verification, and fixes the
whole not as a paper that aims at completion but as an alternative
interpretive hypothesis of the standard theory in a form that can be
taken over and verified.

\begin{center}\rule{0.5\linewidth}{0.5pt}\end{center}

\subsection{0. Position of this paper and how to read
it}\label{position-of-this-paper-and-how-to-read-it}

\begin{enumerate}
\def\labelenumi{\arabic{enumi}.}
\tightlist
\item
  \textbf{The order has meaning.} The order of the chapters is exactly
  the order in which, in the two dialogues, the author asked, Claude
  erred, and the author corrected until the road was fixed. Claude's
  errors are not erased; the author's corrections are written as part of
  the derivation. Nothing has been reordered, merged or shortened. The
  numbers of the eleven stages the author stated on 2026-10-15 are
  attached to the chapter headings. The correspondence between the
  eleven stages and the chapters is: stage 1 = Chapter 7, stage 2 =
  Chapter 9, stage 3 = Chapter 12, stage 4 = Chapter 14, stage 5 =
  Chapter 15, stage 6 = Chapter 17, stage 7 = Chapter 18, stage 8 =
  Chapter 19, stage 9 = Chapter 16, stage 10 = Chapter 11 (and, on the
  four-dimensional lattice, §20.3), stage 11 = Chapter 24. The two
  stages whose position in the dialogue differs from the enumeration
  (spin, the four-dimensional lattice) are placed at their position in
  the dialogue, following the author's decision that the dialogue order
  is the basis and readability may decide. Spin is read out directly
  from the geometry of S³ (left and right multiplication in SU(2)), so
  it reads best right after three dimensions; the four-dimensional
  lattice is the answer to ``why does it lock to that window'' and E2
  fixes the method of E3 onwards, so it reads best right after the
  introduction of α.
\item
  \textbf{Three distinctions.} The paper separates what has been
  verified (what has numbers from papers and numerical experiments),
  what has been connected as hypothesis (what was placed in the dialogue
  in answer to the author's questions), and the means of verification
  (the answer table of the Eleventh Thought Experiment and measured
  values). Chapter 25 lists them.
\item
  \textbf{Units.} G = c = q₀ = 1 (q₀ = e/3, Gaussian). Planck's constant
  is not placed as an input. From α = 9q₀²/(ħc), ħ = 9α⁻¹ is an output
  (Chapter 18).
\item
  \textbf{The meaning of locking.} Dynamic locking---a state being
  pulled in to a nearby value and locking there, with quasi-oscillation
  and sudden change on top---must occur. What would break the hypothesis
  is that the values chosen by several symmetries (137, 23/124, the
  integers of E₈, \ldots) coincided exactly from the start, not the
  discovery of a dynamic lock {[}A1 §7.6{]}. ``Near but not identical''
  is not a hole to be filled; it is the content of the hypothesis.
\item
  \textbf{Wording.} This is a hypothesis paper, not a record of a
  dialogue, and its claims are stated outright. At the head of each
  section, however, the author's question is placed, because it is the
  primary source that fixed the order of the derivation. The quotations
  are the dialogue text with typographical errors, colloquial
  interjections and expressions of reproach removed; the meaning has not
  been changed. The originals are in \texttt{chat\_\allowbreak{}logs/\allowbreak{}}.
\item
  \textbf{Symbols and experiment names of the earlier papers are kept.}
  The section numbers of the sphere-resident paper {[}H1{]}, the
  equation numbers and H1--H5 of the Fourteenth Thought Experiment
  {[}A1{]}, U\_R, N\_A and ε\_state of the Thirteenth {[}D1{]}, the
  comparison table of the Eleventh {[}D2{]}, N(R) and R\_\{n,m\} of the
  K4 paper {[}B1{]}, and the names of the numerical experiments E0--E5
  are unchanged.
\end{enumerate}

\subsection{1. Motivation: from the answer table to this thought
experiment}\label{motivation-from-the-answer-table-to-this-thought-experiment}

This research series stands under the design document {[}G1{]}. Its
starting point is two questions. Why is it difficult to derive the
Standard Model and general relativity from one and the same underlying
structure, although each succeeds with high precision? Why does the
Standard Model, with its high symmetry and many conserved structures,
need several experimental parameters? The candidate first principle is
``there exist states that have no names but are mutually distinguishable
and countable'', and the states are given no attributes such as
electron, photon, mass, charge, spatial position or time coordinate in
advance. Among the prohibitions are fixing spatial or temporal axes from
outside, naming particle species, and adding normalisation merely to
obtain a desired result.

The First Thought Experiment {[}G2{]} derived, from two nameless values
(a, b) and the principle that ``a state alone cannot be read; only
products can'', that the readout is uniquely an area form, that the
first readable dynamical quantity is the central acceleration (the
operational content of the equivalence principle), that the sum of
squares Σz² = C² is induced rather than input, and that C is the radius
of curvature. The Second Thought Experiment {[}G3{]} derived position
too as a product readout, x = a² − b², y = 2ab, r = a² + b², with x² +
y² − r² = 0 emerging as a zero closure, and showed that the readout
point traces a conic obeying an inverse-square force; at the same time
it concluded that ``the mass ratio and the charge ratio cannot be
written with two values''. The Third Thought Experiment {[}G4{]}
reversed the direction and confirmed that, from two-body orbits computed
with the known Newtonian equations and post-Newtonian approximations,
the mass ratio, relative distance and relative time can be read out with
zero knowledge. The Fourth to Sixth Thought Experiments {[}G5--G7{]}
generated the orbits from a state-dependent action Sₙ, made explicit
through self-audit the information remaining outside the state, and
placed the two interactions, gravity and Coulomb, in eight persistent
states and one map. The supplement to the Sixth {[}G8{]} changed the
normalisation from G = M = c = 1 to G = q₀ = c = 1 (q₀ = e/3).

The Eleventh Thought Experiment {[}D2{]} changed course here. The
radiation update rule of the Sixth had been given from outside as known
multipole formulas; ``closed in the state alone'' was the form of the
ledger, not a derivation of the interaction. The attempts to derive the
dynamics including radiation from the distribution-array operation
itself (the Seventh to Tenth Thought Experiments, unpublished), and to
expand the Einstein--Maxwell equations around the initial configuration
and place them in a discrete variational scheme, neither reached a
two-body orbit with gravity, and the run with an absorber did not reach
a lock but collapsed. So the Eleventh did not aim at derivation; it
computed the known standard theory (Dirac--Coulomb bound states with
recoil, 1PN two-body gravity with quadrupole gravitational waves,
Einstein's A coefficients) without holding any coordinate grid or metric
of a background spacetime in the state, with the rest frame of the
complete absorber at infinity as the only reference frame, in units G =
c = q₀ = 1, with inputs α, μ\_p, μ and g\_p, reading out only from the
record of emission and absorption (frequency, rate, branching ratio,
time difference, direction, polarisation), for four systems (hydrogen
atom, electron--electron, proton--proton, hydrogen molecule), and
produced the ``answer table'' that a future generating rule must
reproduce. The absorber is the set of sections where the future light
cone of each emission event cuts infinity, and ``no background
spacetime'' and ``a complete absorber at infinity'' do not contradict
each other. Recorded exchanges (emission and absorption) give the
frequency itself; unrecorded exchanges give its shift (the Lamb shift).
For ``the ground state does not radiate'', both the time evolution of
the field and a discrete lock are needed.

The Twelfth Thought Experiment {[}F1{]} asked where the α used as an
input in the answer table had come from. In the hydrogen ground state α
= v/c, and the Lorentz factor γ = 1/√(1−α²) is nothing but a component
ratio in the light-cone equation with one proper-time axis added, r² +
τ² + (it)² = 0. Treating the axes as nameless directions, when an
observer chooses the subspace that can be read directly, the remaining
directions appear only as a sum of squares, and symmetry can be read as
an expression of indistinguishability. The Thirteenth Thought Experiment
{[}D1{]} asked why a particle that has acquired mass does not dissipate
and vanish, showed that localisation shuttles reversibly between two
waves, that U\^{}n = I is a result and not a cause, and that in an
exchange system with the finite-order root R\_\{124,23\} the state
recovers at the 124th collision, but reserved: ``R is an input value,
and this is not a confirmation that the periodic closure arose from the
interaction; to determine the mixing angle from the interaction a
dynamics is needed.'' The Fourteenth Thought Experiment {[}A1{]} traced
the question back behind Planck's constant as the end point of
explanation, set up the hypothesis that several integer-selecting
symmetries cannot be fully reconciled and pull each other in, the
incomplete lock producing metastability and sudden change, and proposed
``to search for pairs that look theoretically unrelated yet lie close,
that pull in but never lock completely''. Its trigger was that the
four-dimensional cell count 137 = 1 + 8 + 128 and the finite-order root
of two-wave exchange R\_\{124,23\} = cos²(23π/124) each separately give
a value near α⁻¹.

The series up to this point had built an answer table with the known
equations, asked after the origin of the constant left in its input (α,
the charge coupling), found the structure of exchange and incomplete
locking, and stopped there. What remained was to show, in one system and
one stage at a time, where mass, charge, translation and spin are read
out from, without placing either time or space. The thought experiment
of this paper begins there: with a resident of a closed circle.

\begin{center}\rule{0.5\linewidth}{0.5pt}\end{center}

\subsection{2. Residents of the circle
(S¹)}\label{residents-of-the-circle-suxb9}

From this chapter to Chapter 6 the development of the dialogue
``Intrinsic geometry on a circle'' (2026-10-13 to 14) {[}H1{]} is
written in the order of the author's questions and corrections. The
author's question is placed at the head of each section. All formulas
are written with nothing but the phase θ on the closed circle and the
imaginary-direction amplitude A; no background time or space is used.

\subsubsection{2.1 Setting: no background time or space is
assumed}\label{setting-no-background-time-or-space-is-assumed}

\begin{quote}
Consider a resident living on a circle of radius R in two dimensions. No
background time or space is assumed. (10-13 17:59)
\end{quote}

The resident's world is the closed circle alone. The centre, and the
plane inside and outside the circle, do not exist for the resident.
``Radius R'' is a word seen from outside and is not something the
resident can measure from inside. A one-dimensional manifold has no
intrinsic curvature, so the only clue left to the resident is the
closure, ``one lap brings you back''. There is no background time
either, so if time is to be spoken of it must be assembled from the
relations among changes within the resident's world.

\subsubsection{2.2 No direction for a transverse wave: the amplitude is
in the imaginary
direction}\label{no-direction-for-a-transverse-wave-the-amplitude-is-in-the-imaginary-direction}

\begin{quote}
Now consider a cosine or sine wave. We take the wave to be transverse
only. There is no direction for a transverse wave; it should be
orthogonal to the tangent of the circle, but no such direction exists,
so we assume the amplitude may also lie in the imaginary direction.
(18:03)
\end{quote}

The direction orthogonal to the tangent is the direction of a background
space, which the resident does not have. The only direction in which the
amplitude of a transverse wave can be placed is an internal degree of
freedom attached to the resident's world, not a spatial direction. That
is taken to be the imaginary direction. Because the circle is closed,
the allowed waves are discrete.

\subsubsection{2.3 The number of nodes n: the phases where the amplitude
vanishes}\label{the-number-of-nodes-n-the-phases-where-the-amplitude-vanishes}

\begin{quote}
Here n is the point: think of it as the number of nodes (vertices) on
the circle. (18:05) No.~Nodes (vertices) are the phases where the
amplitude becomes zero. (18:06) Can a complex amplitude not become zero?
(18:07) Superfluous assumptions have crept in. All there is are the
phase θ and the imaginary amplitude A. (18:09)
\end{quote}

Here Claude erred three times. It read the points where the phase aligns
as ``nodes''; then restricted to the circularly polarised e\^{}\{iks\}
and said ``a complex amplitude cannot vanish''; then introduced C·cos(mθ
+ δ) with a complex constant C, a phase δ and an integer m. Through the
author's three corrections the elements of the system were fixed as two:
θ and the imaginary amplitude A. A node is a phase where the amplitude
vanishes, and there are n such θ in one lap.

\begin{quote}
Do you see the meaning of n here? 2n gives the ordinary 360 degrees; n
gives 180 degrees. (18:10)
\end{quote}

From node to node is 180 degrees of phase. With 2n nodes the wave
advances n×360 degrees per lap, an ordinary wave; with n nodes it
advances n×180 degrees. n counts half-wavelengths. Waves with odd n are
therefore allowed, and for them one lap returns the amplitude A as −A.
Since A may take values in the imaginary direction, −A is A rotated by
180 degrees in the complex plane and causes no contradiction.

\subsubsection{2.4 n = 0 is ±iA}\label{n-0-is-ia}

\begin{quote}
Then what waveform is n = 0? (18:11) ±iA, surely? (18:13)
\end{quote}

Claude answered that n = 0 is ``a uniform A with nothing appearing'',
and the author corrected it. At n = 0 the phase does not advance, so the
amplitude remains constant, and that amplitude can only be placed in the
imaginary direction. n = 0 is not empty: it is the state in which the
whole circle uniformly carries a purely imaginary amplitude, and there
are two states distinguished by sign, +iA and −iA.

\subsubsection{2.5 n = 1, n = 2, odd and
even}\label{n-1-n-2-odd-and-even}

\begin{quote}
Then n = 1? (18:14) Can you draw it? (18:14) Make it 3D. It should be
obvious that taking the imaginary axis in the radial direction is wrong.
(18:16) Now draw n = 2. (18:17) Is there a difference here between odd
and even? (18:19)
\end{quote}

n = 1 is the wave with one node. From +iA at some point it decreases
along the circle, passes through zero at the single point on the
opposite side (the node), and after one lap it has become −iA. n = 1 is
the wave that joins the two states +iA and −iA of n = 0 across a single
node; it does not return after one lap but reverses sign, and returns
only after two laps.

In the figure, the imaginary axis cannot be taken in the radial
direction, because the radial direction is the direction of a background
space. The imaginary amplitude is drawn perpendicular to the plane of
the circle.

n = 2 has two nodes and advances 360 degrees per lap; starting from +iA
it crosses the plane twice at the two nodes and returns to +iA after one
lap.

For odd n the phase advances an odd multiple of 180 degrees per lap, so
after one lap +iA has become −iA and two laps are needed to return. Even
n closes in one lap. Under the rotation that goes once round the circle,
odd waves return after 720 degrees and even waves after 360 degrees.
This difference in how they return has the same structure as the
difference between a spinor and a vector.

\subsubsection{2.6 Summing harmonics and
localisation}\label{summing-harmonics-and-localisation}

\begin{quote}
Can you draw n = 1, 3, 5 in one figure? (18:20) Add the linear sum as
well. (18:21) Likewise show the linear sum of n = 2, 4, 6. (18:23) Now
show only the linear sum of the odd orders up to n = 17. (18:24)
Likewise only the linear sum of the even orders up to n = 18. (18:26)
Now only the linear sum of odd and even together up to n = 18. (18:27)
\end{quote}

Each component is summed with the same amplitude A.

\begin{itemize}
\tightlist
\item
  Odd orders n = 1, 3, \ldots, 17 (9 components): a single sharp peak of
  +9iA stands at θ = 0; at the end of the first lap the same point is a
  trough of −9iA, and it returns to +9iA on the second lap. There are 17
  nodes per lap, equally spaced at θ = π/9. The closed form is
  A·sin(9θ)/(2 sin(θ/2)).
\item
  Even orders n = 2, 4, \ldots, 18 (9 components): +9iA at θ = 0,
  returning to +9iA at the same point after one lap. At θ = π the value
  is −iA, not a node. There are 18 nodes per lap, with two interleaved
  equally spaced sequences of different spacing.
\item
  Both, n = 1, \ldots, 18: +18iA at θ = 0. At the end of the first lap
  the odd sum −9iA and the even sum +9iA cancel and the amplitude is
  exactly zero; but at this point the wave only touches the plane of the
  circle without crossing it. At the end of the second lap it returns to
  +18iA. The closed form is A·sin(9θ/2)·cos(19θ/4)/sin(θ/4).
\end{itemize}

\begin{quote}
Do you see the intention? (18:30)
\end{quote}

Claude's reading was as follows, and the author said ``roughly OK''.
Since there is no background space, the direction of the amplitude is i,
and it cannot be taken radially because that is the direction of a
background space. From the closure of the circle alone the allowed waves
become discrete, and n is the only quantity the resident can read
out---the resolution itself. n = 0 is the two values ±iA. Odd n closes
in two laps and even n in one, so the distinction between fermionic and
bosonic waves comes out of the closed circle and the parity of the
number of nodes alone. Summing equal-amplitude harmonics localises a
sharp peak at one point, and a structure that can be called ``position''
is born from interference alone although there is no background space.
Summing odd and even together, the two cancel at the end of the first
lap and the localised peak appears only once in two laps. From the
closed circle and the imaginary-direction amplitude alone, quantisation
(n), the fermion/boson distinction (parity) and a localised particle
(harmonic summation) come out without bringing anything in from outside.

\subsubsection{2.7 The peak phase θ₀ and the three
formulas}\label{the-peak-phase-ux3b8ux2080-and-the-three-formulas}

\begin{quote}
Add to this figure a peak phase θ0 that can be moved with a slider.
(18:44) Show the formulas for the linear sums fodd(n,θ0), feven(n,θ0),
fall(n,θ0) when the maximum n and θ0 are specified, for odd only, even
only and both. (18:48)
\end{quote}

Put φ = θ − θ₀ and let the amplitude be iA. N is the largest n included
in the sum.

\[
f_{\mathrm{odd}}(N,\theta_0)=iA\sum_{k=0}^{(N-1)/2}\cos\frac{(2k+1)\varphi}{2}=iA\,\frac{\sin\frac{(N+1)\varphi}{2}}{2\sin\frac{\varphi}{2}},
\]

\[
f_{\mathrm{even}}(N,\theta_0)=iA\sum_{k=1}^{N/2}\cos k\varphi=iA\,\frac{\sin\frac{N\varphi}{4}\cos\frac{(N+2)\varphi}{4}}{\sin\frac{\varphi}{2}},
\]

\[
f_{\mathrm{all}}(N,\theta_0)=iA\sum_{n=1}^{N}\cos\frac{n\varphi}{2}=iA\,\frac{\sin\frac{N\varphi}{4}\cos\frac{(N+1)\varphi}{4}}{\sin\frac{\varphi}{4}}.
\]

The peak values are i(N+1)A/2, iNA/2 and iNA respectively; the periods
are 4π (two laps) for f\_odd and f\_all and 2π (one lap) for f\_even.
The θ₀ slider runs from 0 to 720°; θ₀ and θ₀ + 360° put the peak at the
same point of the circle but differ in whether the peak appears on the
first or the second lap, and are different states.

\subsubsection{2.8 Exponential representation and the E
function}\label{exponential-representation-and-the-e-function}

\begin{quote}
Write this in exponential form, including the amplitude. (18:49) Put the
amplitude in exponential form too. (18:51) You have ignored that the
amplitude is imaginary. (18:53) Do not skip the expansion of Σ. (18:54)
Now show, as Eeven and so on, the functions that give only the exponent.
(18:56) Do you see what is being done? (18:58)
\end{quote}

After four corrections, the amplitude iA = e\^{}\{α+iπ/2\} (α = ln A)
was kept as an independent factor and the sums were closed as geometric
series.

\[
f_{\mathrm{odd}}(N,\theta_0)=e^{\alpha+i\frac{\pi}{2}}\;\frac{e^{i\frac{(N+1)\varphi}{2}}-e^{-i\frac{(N+1)\varphi}{2}}}{2\left(e^{i\frac{\varphi}{2}}-e^{-i\frac{\varphi}{2}}\right)},
\]

\[
f_{\mathrm{all}}(N,\theta_0)=e^{\alpha+i\frac{\pi}{2}}\;\frac{\left(e^{i\frac{N\varphi}{4}}-e^{-i\frac{N\varphi}{4}}\right)\left(e^{i\frac{(N+1)\varphi}{4}}+e^{-i\frac{(N+1)\varphi}{4}}\right)}{2\left(e^{i\frac{\varphi}{4}}-e^{-i\frac{\varphi}{4}}\right)}.
\]

The differences in numerator and denominator are both purely imaginary,
so their ratio is real and f always lies on the imaginary axis. Only
e\^{}\{iπ/2\} fixes the direction of the amplitude; its size is
e\^{}\{α\}.

Let E\_odd, E\_even, E\_all be the exponents with f = e\^{}\{E\}.
Products become sums and quotients become differences.

\[
E_{\mathrm{odd}}(N,\theta_0)=\alpha-\ln2+i\frac{\pi}{2}+\ln\!\left(e^{i\frac{(N+1)\varphi}{2}}-e^{-i\frac{(N+1)\varphi}{2}}\right)-\ln\!\left(e^{i\frac{\varphi}{2}}-e^{-i\frac{\varphi}{2}}\right).
\]

The real part of E is the logarithm of the size of the amplitude,
ln\textbar f\textbar, and diverges to −∞ at the nodes where
\textbar f\textbar{} = 0. The imaginary part of E takes only the two
values π/2 and 3π/2 (mod 2π): the direction of the amplitude is either
+i or −i and jumps by π at each node.

To the author's ``do you see what is being done?'', Claude's reading was
as follows, and the author said ``roughly OK''. Not the wave itself but
its exponent E = ln f is being taken as the descriptive quantity of the
system. The amplitude α = ln A is only a constant added to E and
disappears on differentiation with respect to θ or θ₀, so the quantities
read out as gradients do not depend on A. ``Only ratios are physics''
holds automatically in the formula. The remaining parameters are N and
θ₀ only. The readout is the gradient of E, and since φ = θ − θ₀, ∂E/∂θ₀
= −∂E/∂θ: moving the peak and reading at that place are tied by the same
quantity.

\subsubsection{2.9 In two dimensions R does not
enter}\label{in-two-dimensions-r-does-not-enter}

\begin{quote}
From here the bar goes up. In the two-dimensional case R does not enter
this calculation, does it? (19:00)
\end{quote}

E\_odd, E\_even and E\_all are all written in φ alone, and R appears
nowhere. R appears only when the phase θ is converted to arc length s =
Rθ. The resident has no background ruler; the only unit of length
available is the wave itself (the node spacing), and measured in that
unit 1/R disappears and only the ratio of node numbers N remains. The
circle is one-dimensional and has no intrinsic curvature, so whatever
the resident measures locally, R cannot be detected.

\subsubsection{2.10 θ₀ is a virtual origin: the base n = 1 or n = 2 is
the
gauge}\label{ux3b8ux2080-is-a-virtual-origin-the-base-n-1-or-n-2-is-the-gauge}

\begin{quote}
Right. There is one more strange thing: θ0. This system has no
background origin, and yet θ0 is somehow defined. (19:06) The formula
holds without any such assumption. (19:07) Exactly. The base n = 1 or n
= 2 becomes the gauge and carries a virtual origin. (19:08)
\end{quote}

Claude answered that ``θ₀ has meaning only when there are two or more
waves, and only the phase difference can be read'', and the author
corrected it. No second wave is needed. In the exponential
representation,

\[
e^{i\frac{n\varphi}{2}}=e^{i\frac{n\theta}{2}}\,e^{-i\frac{n\theta_0}{2}},
\]

so θ₀ is not an origin placed from outside but the arrangement of the
phases nθ₀/2 of the components n themselves. The point where the phases
of all components align is the peak, and that point is θ₀. Without a
background origin, the phase relations among the components fix a single
point on the circle by themselves. This holds because they are summed:
with a single component the points of aligned phase are several and none
can be singled out as θ₀. The appearance of a localised peak and the
fixing of an origin are one and the same event.

The fundamental component has only one peak within its period: once in
two laps (4π) for n = 1, once in one lap (2π) for n = 2. So the phase of
the fundamental alone fixes a single point on the circle. That is the
virtual origin. A higher component n has n peaks and cannot fix an
origin by itself, but in the composite wave the phase nθ₀/2 of each
component is fixed as an integer multiple of the phase of the
fundamental, and the higher components merely line up according to the
origin the fundamental has fixed. The phase of the fundamental itself
can be chosen freely; shifting θ₀ to θ₀ + δ shifts all component phases
by nδ/2 in step and the waveform moves without changing shape. The
choice of the fundamental's phase is the gauge freedom, and what is
invariant under that choice is the ratio of the component phases (the
integer-multiple relation). In the odd and the full systems n = 1 plays
this gauge role; in the even system n = 2 does.

\subsection{3. To residents of the four-dimensional spherical
shell}\label{to-residents-of-the-four-dimensional-spherical-shell}

\subsubsection{3.1 R enters and θ₀ becomes
three}\label{r-enters-and-ux3b8ux2080-becomes-three}

\begin{quote}
Next. Extend this to residents of the surface of a four-dimensional
spherical shell. Suddenly R enters and θ0 becomes three, does it not?
(19:12)
\end{quote}

The surface S³ of a four-dimensional ball is a closed three-dimensional
space. Three angles are needed to specify a point of the resident's
world, so the peak position of the composite wave, the virtual origin,
is fixed by three angles (χ₀, θ₀, φ₀), and the gauge freedom becomes
three.

The circle was one-dimensional with no intrinsic curvature and R was
invisible to the resident, but S³ has intrinsic curvature 1/R² and the
resident can read R from local measurements without looking outside. On
the side of the wave, where on the circle the eigenvalue of component n
was proportional to n², for harmonics on S³ it is n(n+2)/R². This can be
written (n+1)² − 1, and the ``−1'' is the term due to curvature. The
ratio between the ruler measured by the node spacing of the wave and the
bending of the space remains as a quantity the resident can read. S³ has
the same shape as SU(2) and possesses, as a space, the spinor property
of reversing sign in one lap and returning in two. The ``wave that
closes in two laps'' seen in the odd system of the circle appears on S³
as a property of the space itself.

\subsubsection{3.2 The same exponential form: N and R are
common}\label{the-same-exponential-form-n-and-r-are-common}

\begin{quote}
Can this be expressed in the same exponential form as before? Take N and
R to be common. (19:14)
\end{quote}

Let s be the geodesic distance between a point and the peak and put ψ =
s/R.

\[
\cos\psi=\cos\chi\cos\chi_0+\sin\chi\sin\chi_0\left[\cos\theta\cos\theta_0+\sin\theta\sin\theta_0\cos(\phi-\phi_0)\right].
\]

Regarding S³ as the unit quaternions, this can also be written x =
x₀·e\^{}\{ψu\} (u a pure imaginary unit quaternion), and ψ is the angle
of the exponent as seen from the virtual origin. Replacing φ/2 by ψ in
the formulas for the circle,

\[
f_{\mathrm{odd}}=e^{\alpha+i\frac{\pi}{2}}\,\frac{e^{i(N+1)\psi}-e^{-i(N+1)\psi}}{2\left(e^{i\psi}-e^{-i\psi}\right)},
\]

\[
f_{\mathrm{even}}=e^{\alpha+i\frac{\pi}{2}}\,\frac{1}{2}\left(\frac{e^{i(N+1)\psi}-e^{-i(N+1)\psi}}{e^{i\psi}-e^{-i\psi}}-1\right),
\]

\[
f_{\mathrm{all}}=e^{\alpha+i\frac{\pi}{2}}\,\frac{1}{2}\left(\frac{e^{i\frac{(2N+1)\psi}{2}}-e^{-i\frac{(2N+1)\psi}{2}}}{e^{i\frac{\psi}{2}}-e^{-i\frac{\psi}{2}}}-1\right).
\]

Three things change relative to the circle. The denominator e\^{}\{iψ\}
− e\^{}\{−iψ\} was, on the circle, merely the result of taking the sum;
on S³ it is the geometric factor corresponding to the size R sin(s/R) of
the sphere at distance s from the virtual origin, and its deviation from
the s that flat space would give is the curvature 1/R². R enters the
resident's physics here. f\_odd is itself a single harmonic on S³ with
eigenvalue N(N+2)/R²: the ``composite wave'' on the circle is one wave
on S³. One lap of the circle (φ = 2π) corresponds on S³ to the antipode
(ψ = π); f\_odd reverses sign at the antipode, f\_even returns, and
f\_all with even N is exactly zero there.

\subsubsection{3.3 The three θ₀: two phases and one mixing
angle}\label{the-three-ux3b8ux2080-two-phases-and-one-mixing-angle}

\begin{quote}
Now put in the three θ0. (19:16)
\end{quote}

Representing a point of S³ by two complex numbers (z₁, z₂) = (cos
η·e\^{}\{iξ₁\}, sin η·e\^{}\{iξ₂\}) and the peak by (η₀, ξ₁₀, ξ₂₀),

\[
\cos\psi=\mathrm{Re}\left[\cos\eta\cos\eta_0\,e^{i(\xi_1-\xi_{10})}+\sin\eta\sin\eta_0\,e^{i(\xi_2-\xi_{20})}\right].
\]

ξ₁₀ and ξ₂₀ appear only as phase differences, like e\^{}\{−inθ₀/2\} on
the circle, and are gauges of the same kind as on the circle. η₀ appears
as cos η₀ and sin η₀, fixing the ratio of the two phase terms: a mixing
angle, not a phase. Whatever coordinates are chosen, the three θ₀ cannot
all be made pure phases at once, because among the rotational symmetries
of S³ at most two rotations commute. Setting η = η₀ = 0 recovers the
circle. The single θ₀ of the circle becomes on S³ ``two phases (gauges)
and one mixing angle''.

\subsubsection{3.4 One equation localises at any position of
three-dimensional
space}\label{one-equation-localises-at-any-position-of-three-dimensional-space}

\begin{quote}
Is this not strange? With (A,N,θn), n = 1,2,3, the same equation can
express movement to and localisation at any position of the
three-dimensional space of the universe, can it not? (19:23)
\end{quote}

Every point of S³ is equivalent (homogeneous), so however the three θ₀
are chosen, the waveform around the peak keeps the same shape and can be
localised at any point of the space. One equation with five parameters
(A, N, θ₁, θ₂, θ₃) can place a localised object anywhere in
three-dimensional space. This movement is not movement through a
background space; because S³ is SU(2) itself, it is realised by the
symmetry of the space, as multiplication by an element of the group. A
is the strength of the peak, N its width (about πR/(N+1), the ratio of
the particle's size to the size of the universe), and the three θ₀ the
position. What the equation expresses at this point is the position and
movement of a localised object; the odd and even systems distinguish
fermionic from bosonic. Where a quantity such as charge is read out from
has not yet been shown.

\subsubsection{3.5 Why does the wave appear to translate: it is at
rest}\label{why-does-the-wave-appear-to-translate-it-is-at-rest}

\begin{quote}
And why does the wave appear to translate? That cannot be explained yet,
can it? It is at rest. (19:25)
\end{quote}

Time enters nowhere in the formulas so far. Every f is a stationary
waveform, θ₀ is only a parameter saying where it is localised, and
waveforms with different θ₀ are different stationary states. For
translation to be seen, the states with different θ₀ would have to be
ordered and one would have to pass into another; that is the same as
bringing in a direction of time, and this system, which assumes no
background time, has no such mechanism yet. What can be said from the
formula is that the phase of each component is e\^{}\{−inθ₀/2\}, so for
the peak to translate means that the phases of all components advance in
step in proportion to n.~The problem narrows to ``does a common
parameter that advances the phases of all components in proportion to n
come out of the system itself, without being brought in from outside?''
Chapter 7 answers this question.

\subsubsection{3.6 R becomes R(θn, φn): fixed by Σxₙ² =
R²}\label{r-becomes-rux3b8n-ux3c6n-fixed-by-ux3c3xux2099uxb2-ruxb2}

\begin{quote}
Here it suddenly becomes difficult when R becomes R(θn). In reality a
direction of the centroid is needed, so it is R(θn,φn). (19:28) Properly
it is fixed by Σxn\^{}2 = R\^{}2. (19:32)
\end{quote}

While R is constant every point of S³ is equivalent and the three θ₀ are
pure gauges. When R varies from place to place, θ₀ ceases to be a gauge,
the closed formula ceases to be an eigenfunction, the node spacing
stretches and shrinks in proportion to the local R and the wave is
refracted along the gradient of R. To define R(θ) a centre ``from which
the radius is measured'' is needed; if that is fixed by the centroid, a
circle arises in which the centroid is fixed by the distribution of the
wave, R by the centroid, and the wave by R.

R is fixed neither from outside nor from \textbar f\textbar² by some
separate rule, but by the sum of squares of the components themselves:

\[
\sum_n x_n(\theta_n,\phi_n)^2=R(\theta_n,\phi_n)^2 .
\]

The circle of wave and geometry is not an assumption added later; it is
contained in this formula from the start. That the wave fixes the shape
of the shell and the shape of the shell fixes the wave on it are two
faces of one formula. Since the square of a component in the imaginary
direction is negative, moving R² to the left and writing it as (iR)²
gives Σxₙ² + (iR)² = 0, and R is not a special quantity but enters the
same sum as one of the components.

\subsection{4. The error of the self-consistent shell, and the true
squares}\label{the-error-of-the-self-consistent-shell-and-the-true-squares}

\subsubsection{4.1 Claude's numerical experiment: absolute squares and a
real
radius}\label{claudes-numerical-experiment-absolute-squares-and-a-real-radius}

\begin{quote}
Can you examine it? (21:01)　Go ahead. (21:04)　Can this be extended to
the four-dimensional spherical shell? (21:32)　Do both and compare.
(21:33)
\end{quote}

Claude computed a self-consistent shell numerically with the following
setting: components xₙ = iA·cos(n(φ−φ₀)/2); the shape of the shell fixed
in each direction as R(θ)² = A²·Σcos²(n(φ−φ₀)/2); the resident's phase
measured by arc length on the deformed shell; the centroid of the shell
recomputed and R re-measured from it; iterated until the shape stops
changing. On the circle, in systems containing odd components the
centroid and the centre from which R was measured did not coincide, and
for both N = 18 the centroid was displaced towards the peak by d = 0.682
(2.86\% of the circumference); for even only d ≈ 0; the peak became a
thin radial spike. On S³, with components (a) cos(nψ) the symmetry ψ →
π−ψ gave d = 0, and with (b) Uₙ(cos ψ)/(n+1) the cases both N = 2, 4
converged to asymmetric solutions. Claude went on to propose ``fixing R
from the square of the composite wave \textbar Σxₙ\textbar²'' {[}H2:
\texttt{円周住民\_\allowbreak{}自己無撞着殻\_\allowbreak{}20261013},
\texttt{四次元球殻住民\_\allowbreak{}自己無撞着殻\_\allowbreak{}20261013}{]}.

\subsubsection{4.2 The author's corrections: true squares, not a norm,
there is no outside of the
circle}\label{the-authors-corrections-true-squares-not-a-norm-there-is-no-outside-of-the-circle}

\begin{quote}
The absolute square? Who decided that definition? It is the true square.
(21:42) Σxn\^{}2 = 0 is not a norm. (21:43) There is no outside of the
circle. (21:44)
\end{quote}

Σxₙ² is the sum of true squares without conjugation, not a
positive-definite norm. With complex components the sum can vanish even
for a non-zero vector; that is the meaning of Σxₙ² = 0. Claude had set
the shell equation as R² = A²·Σcos², adding \textbar xₙ\textbar² of each
component. With true squares, xₙ = iA·cos(\ldots) gives xₙ² = −A²·cos²,
so Σxₙ² = R² gives R² = −A²·Σcos² and R is purely imaginary. Treating R
as a positive real number, Claude had drawn the shell as a figure in the
real plane or real four-space and computed the centroid in it. That is
precisely the background space that had been excluded at the start. The
centroid displacement d, the radial spikes, the d = 0 by symmetry on S³
and the asymmetric solutions all arose from this embedding by a norm and
are not results that follow from the axiom. This paper does not adopt
them as results {[}H1 §5{]}.

Claude further wrote that ``R points in the imaginary direction, that
is, in a direction outside the system'', and the author corrected:
``there is no outside of the circle''. A resident of the circle has only
θ and the imaginary amplitude; the imaginary direction is not an outside
space but the direction of the amplitude the resident has had from the
start. Σxₙ² = R² says: each xₙ is an imaginary amplitude and xₙ² a
negative real number; the R² equal to their sum is also negative, so R
too is imaginary; that is, R is a quantity of the same kind as xₙ, an
amplitude, and not a radius placed outside. R(θ) = iA√(Σcos²(n(θ−θ₀)/2))
is a quantity fixed by the sum of squares of the components at each
phase, closed in θ and the imaginary amplitude alone. This correction is
the premise of every readout that follows.

\subsection{5. The non-uniform geodesic length per gauge
interval}\label{the-non-uniform-geodesic-length-per-gauge-interval}

\subsubsection{5.1 Curvature cannot be read, but the same tick has
different
lengths}\label{curvature-cannot-be-read-but-the-same-tick-has-different-lengths}

\begin{quote}
First we think on a two-dimensional arc, surely? Curvature cannot be
read, but the geodesic length per gauge interval becomes non-uniform.
(21:48) Go ahead and draw it too. (10-14 07:58)
\end{quote}

The resident cannot read curvature. What can be read is the relation
between the gauge intervals of the phase and the corresponding geodesic
lengths. The gauge is cut in the phase θ: the fundamental fixes the
origin and the nodes of the higher components cut θ at equal intervals.
On the other hand R(θ)² = −A²·Σcos²(n(θ−θ₀)/2) varies with θ, so the
geodesic length Δs = ∫R dθ = i∫\textbar R\textbar dθ corresponding to an
equal phase interval Δθ differs from place to place. Since i is common
to all intervals, the ratio Δs/⟨Δs⟩ is real. The nodes of the highest
component n = N (equally spaced in phase, 2π/N) were used as gauge
ticks. The earlier numerical experiment had, conversely, measured the
phase by arc length; here the phase is fixed as the gauge and the
non-uniformity of length is read.

{\def\LTcaptype{none} % do not increment counter
\begin{longtable}[]{@{}llll@{}}
\toprule\noalign{}
System & Peak interval & Antipodal interval & Others \\
\midrule\noalign{}
\endhead
\bottomrule\noalign{}
\endlastfoot
All orders N=18 & 1.245 & 0.988 (ordinary) & 0.97--0.99 \\
Odd orders N=17 & 1.252 & 0.835 (short) & about 1.00 \\
Even orders N=18 & 1.236 & 1.236 (same length) & 0.96--0.98 \\
All orders N=2 / N=4 & 1.159 / 1.211 & 0.841 / 0.938 & --- \\
\end{longtable}
}

The resident cannot read curvature, but at the peak one tick is about
1.25 times longer. Localisation can be read inside the resident's world
as a distortion of the ruler, not as a shape seen from outside. What
distinguishes the peak from the antipode is the odd components. With
even orders alone the peak and the antipode have the same length and
which is the peak is undecided; in the odd systems all components share
a node at the antipode and R → 0, so the peak is long and the antipode
short; in the all-orders system only the peak is long and the antipode
has ordinary length---a single localisation in which one point alone
stands out. The odd/even difference seen first (two-lap or one-lap
return) appears as the asymmetry of geodesic length {[}H2:
\texttt{円周住民\_\allowbreak{}ゲージ間隔測地線長\_\allowbreak{}20261014}{]}.

\pandocbounded{\includegraphics[keepaspectratio]{figs/fig1_R_vs_phase.png}}

\textbf{Fig. 1.} Phase dependence of the amplitude
\textbar R(θ)\textbar{} = R/i fixed by R² = Σxₙ² (true squares). From
top: all orders N=18, odd orders N=17, even orders N=18. Dotted lines
are the nodes of the highest component (gauge ticks).
\texttt{gauge\_\allowbreak{}geodesic\_\allowbreak{}s1.py}, \texttt{plot\_\allowbreak{}gauge\_\allowbreak{}geodesic\_\allowbreak{}s1.py}.

\pandocbounded{\includegraphics[keepaspectratio]{figs/fig2_gauge_interval_lengths.png}}

\textbf{Fig. 2.} Ratio of geodesic length Δs/⟨Δs⟩ for the same tick
(equally spaced in phase, 360°/N). All orders N=2, 4, 18, odd orders
N=17, even orders N=18. Only the peak interval is long (1.245 times); in
the odd orders the antipodal interval is short (0.835 times).

\pandocbounded{\includegraphics[keepaspectratio]{figs/fig3_resident_ruler.png}}

\textbf{Fig. 3.} The resident's ruler. Upper row: gauge ticks equally
spaced in phase; lower row: the positions of the same ticks in geodesic
length (normalised to the total length). The same tick has different
lengths at different places.

\subsubsection{5.2 Negative curvature, and the slope being the same
under the tensor
reading}\label{negative-curvature-and-the-slope-being-the-same-under-the-tensor-reading}

\begin{quote}
That a larger distortion widens the gauge is intuitively easy to
understand, but what happens in the case of negative curvature? (08:29)
Correct. In fact, whether negative or positive, when the distortion is
read as a tensor only the direction differs and the slope is the same.
(08:32)
\end{quote}

With true squares R² = −A²·Σcos², so R is imaginary and 1/R² negative,
and the computation above is already on the side of negative curvature.
At the peak \textbar R\textbar{} is maximal, the curvature is negative
and close to zero, and one tick is long. At the antipode of the odd
system R → 0, the curvature goes to −∞ and one tick is short. If the
components are taken as the ± rotation terms of the exponential
representation, the sum of true squares is −(A²/2)·Σcos(nφ) and changes
sign. For all orders, near the peak Σcos(nφ) ≈ N \textgreater{} 0 so R²
\textless{} 0 (R imaginary, negative curvature); away from it ≈ −1/2 so
R² \textgreater{} 0 (R real, positive curvature); at the boundary R = 0
and the length of one tick is zero. In this case Δs is imaginary in the
peak interval and real elsewhere, and the kind of length itself is
exchanged.

The metric is the single component g = R²dθ². Read as a tensor, what
matters is not g itself but the gradient of its logarithm; the
connection is Γ = ½∂θ ln g. For R² \textless{} 0, ln g =
ln\textbar R²\textbar{} + iπ; for R² \textgreater{} 0, ln g =
ln\textbar R²\textbar; the difference is only the constant iπ, which
disappears on differentiation with respect to θ. The slope ∂θ
ln\textbar R²\textbar{} is the same for positive and negative curvature,
and the sign decides only the direction (imaginary length or real
length). This is the same structure as in the E function: the real part
ln\textbar f\textbar{} carried all the gradient information and the
imaginary part only showed the direction with the two values π/2 and
3π/2. At the boundary R = 0, ln\textbar R²\textbar{} goes to −∞ and
across it the imaginary part jumps by π. The reversal of the sign of
curvature can be read as the π jump of the phase that occurs at a node
of the metric.

\subsubsection{5.3 Drawing the real curve radially, orthogonal to the
imaginary
amplitude}\label{drawing-the-real-curve-radially-orthogonal-to-the-imaginary-amplitude}

\begin{quote}
Now draw this real curve in the radial direction, orthogonal to the
imaginary-direction amplitude on the circle. (08:36)
\end{quote}

The circle is laid in a horizontal plane, the imaginary amplitude is
drawn perpendicular to the plane and the slope ∂θ
ln\textbar R²\textbar{} in the radial direction. Coral marks R²
\textless{} 0, green R² \textgreater{} 0, and black dots the boundary R
= 0; the radial curve continues as a single curve across the change of
colour. With ``components = iA cos'', R² is negative all the way round
and no black dots appear. Moving θ₀ through 360° reverses the amplitude
in systems containing odd components, but the slope, built from squares,
is unchanged {[}H2:
\texttt{circle\_\allowbreak{}imag\_\allowbreak{}amplitude\_\allowbreak{}radial\_\allowbreak{}log\_\allowbreak{}gradient.html}{]}.

\subsection{6. Slices of the four-dimensional spherical shell and the
antipode}\label{slices-of-the-four-dimensional-spherical-shell-and-the-antipode}

\subsubsection{6.1 The projection surface is the radius of the
shell}\label{the-projection-surface-is-the-radius-of-the-shell}

\begin{quote}
Now I want to extend this figure to the four-dimensional spherical
shell, but it cannot be drawn without cutting a projection surface
somewhere. Onto which surface do you project? And can the position of
the projection surface be made slidable? (08:40) Both. I want the peak
to be easy to see. Projected to three dimensions it looks like a sphere.
I want to cut it in the way that makes the size of the sphere change as
it slides, so the projection surface is precisely the radius of the
shell, is it not? (08:44) Fix it, and vary the transparency with the
shading too. (08:46)　Add an on/off for the slope lines. (08:57)　Add
odd and even check boxes, both on by default. (09:00)
\end{quote}

A point of S³ is written (sin χc·n̂, cos χc) and cut by the spherical
shell at geodesic distance χc from the pole (radius sin χc). As χc
moves, the shell swells from a point, is largest at the equator, and
shrinks again. When the cut reaches the latitude of the peak fixed by
the three θ₀, a dark spot and a radial bulge appear in the direction
(θ₀, φ₀). The imaginary amplitude is drawn by colour and opacity (the
ratio of \textbar amplitude\textbar{} to the global peak value, the
number of components), the slope ∂ψ ln\textbar R²\textbar{} by radial
relief, and the sign of R² by the boundary line R = 0. The components
can be switched between cos(nψ) (direct correspondence from the circle)
and Uₙ(cos ψ)/(n+1) (the eigenmodes of S³), and between ± rotation terms
and iA × component. With true squares, R² is −(A²/2)·Σcos(2nψ) for
cos(nψ) with ± rotation terms, −A²·Σcos²(nψ) for cos(nψ) with iA ×
component, A²·Σcos(2(n+1)ψ)/(2(n+1)²sin²ψ) for Uₙ/(n+1) with ± rotation
terms, and −A²·ΣUₙ²/(n+1)² for Uₙ/(n+1) with iA × component {[}H2:
\texttt{s3\_\allowbreak{}shell\_\allowbreak{}slice\_\allowbreak{}amplitude\_\allowbreak{}gradient.html}{]}.

\subsubsection{6.2 Even orders too raise a peak at the antipode: without
both parities the position is
ambiguous}\label{even-orders-too-raise-a-peak-at-the-antipode-without-both-parities-the-position-is-ambiguous}

\begin{quote}
In three dimensions, does a peak not form at the symmetric point even
for even orders? (09:02) Can N be specified separately for odd and even
orders? (09:07) This is easy to see. Indeed, with even components
present a peak always forms at the symmetric point. (09:25) So,
properly, without both odd and even, the position becomes ambiguous.
(09:26)
\end{quote}

At the antipode ψ = π the value of a component is (−1)ⁿ in both the
direct-correspondence form and the eigenmode form. With even orders
alone all components are +1 and a second peak of +i, equal to the first,
stands at the antipode; with odd orders alone all components are −1 and
a peak of reversed sign −i stands there; with all orders (N even) +1 and
−1 are equal in number and cancel, and the antipode is zero. The
antipode ψ = π of S³ corresponds to one lap (φ = 2π) of the circle, and
the parity difference that appeared on the circle as ``how many laps to
close'' appears on S³ as the spatial arrangement ``what stands at the
antipode''. With separate maxima N for odd and even orders, the value at
the antipode is the number of even components minus the number of odd
ones (6 odd with N = 11 and 6 even with N = 12 cancel; odd N = 13 leaves
a faint −i spot; even N = 14 a faint +i spot).

``Always, with even components present'' cannot quite be said.
Precisely: even components make a peak of the same sign at the antipode,
odd components one of the opposite sign, and when both are present they
cancel and the localisation is at one point. With even orders alone,
equal peaks stand at p and −p and the resident cannot tell which of the
two points it is at; the even components are functions on the space that
identifies p with −p (S³ divided by ±) and cannot point to a single
point in the first place. With odd orders alone, p is +i and −p is −i
and they seem distinguishable by sign, but the direction ±i of the
amplitude is a gauge that can be reversed as a whole; reversing the
overall sign exchanges the two peaks, their sizes are equal and the
square R² is the same at p and −p, so the position stays ambiguous
between the two points. With both present they cancel at the antipode
and a single peak remains; reversing the overall sign does not move the
position of the peak, and only here is the position fixed uniquely. The
even system is the bosonic component and the odd system the fermionic
one; localisation at one point does not occur with either alone, only
when both are present and balanced.

\subsubsection{6.3 The symmetric point on the far side of the universe
has not been observed. A photon, a single tone, has no position. n = 1
alone is
special}\label{the-symmetric-point-on-the-far-side-of-the-universe-has-not-been-observed.-a-photon-a-single-tone-has-no-position.-n-1-alone-is-special}

\begin{quote}
But it is strange: has anyone observed that a boson has a symmetric
point on the far side of the universe? (09:27) A photon is generally
regarded as a single tone without harmonics, so the localisation in
spacetime itself has no meaning, has it? (09:28) But even the state with
odd 1 alone shows some localisation in the figure, does it not? (09:36)
So it diverges at the boundary R = 0. But this is a typical ten-gallon
hat. (09:38) I merely thought it resembles a potential curve with
symmetry. (09:40)
\end{quote}

It has not been observed. Even if the universe is a closed S³, the
curvature is almost zero, the antipode lies beyond the horizon, and
searches for pairs of ghost images in the cosmic microwave background
have found no trace. On the side of the model there is no need to read
``a boson has a double on the far side of the universe''. A wave of even
components alone always takes the same value at p and −p; for the even
components p and −p are the same single point from the start, and there
is no place called ``the far side''. p and −p become distinguishable as
different places only when odd components are added. The even system
(bosonic) is not a particle localised at one point but a spread
existence that does not distinguish p from −p, and a particle localised
at one point appears only when both the odd and the even systems are
present.

A single component n alone has nodes and peaks lined up periodically and
no peak that points to one point; since the appearance of localisation
and the fixing of the origin are one event, a single tone has no
position defined, and all that is fixed is n (frequency, energy). The
position of a photon is not something the photon has; it is read out,
when it meets a localised system (a composite wave with odd and even
together), by the localisation on the other side.

n = 1 alone, however, is special. The single component n = 1 is cos ψ
and has only one peak and one trough over the whole of S³. This is
precisely the property that ``n = 1 becomes the gauge and, having only
one peak in its period, fixes the origin uniquely''. A single tone with
n ≥ 2 has several peaks of the same height and none can be chosen, but n
= 1 has one peak. It differs from the localisation made by summing
harmonics in two respects. The width is ψ \textless{} 90°, spread over
half of the space. Reversing the overall sign ±i exchanges peak and
trough, so what n = 1 fixes is strictly not a point but the ``axis and
orientation'' joining p and −p, and only when even components are added
to cancel the antipode is it narrowed to one point. n = 1 is the gauge
that gives the orientation of position; the harmonics (odd system)
sharpen the peak and the even system erases the antipode, and the
localisation at one point is complete.

With ± rotation terms and n = 1, R² = −(1/2)·cos 2ψ and the slope is ∂ψ
ln\textbar R²\textbar{} = −2·tan 2ψ. In the crown (ψ \textless{} 45°) R²
\textless{} 0; on the ring ψ = 45°, R = 0 and the slope jumps from −∞ to
+∞; in the brim (ψ \textgreater{} 45°) R² \textgreater{} 0. The ring is
a two-sphere surrounding the peak inside S³ and appears as a circle in
the slice. A rotationally symmetric ring around a centre, a central
rise, a wall around it and an outer brim is the shape of a symmetric
potential curve, but here that shape comes from inside, as the slope of
the logarithm of R fixed by Σxₙ² = R², not as a potential given from
outside.

\subsection{7. Stage 1: first, read out
mass}\label{stage-1-first-read-out-mass}

From this chapter we enter the order of derivation (eleven stages) the
author stated on 2026-10-15. Stage 1 is the readout of mass.

\subsubsection{7.1 A travelling wave cannot be
represented}\label{a-travelling-wave-cannot-be-represented}

\begin{quote}
Let us begin a different examination. This cannot represent a travelling
wave. What do you think? Does it travel in the direction of the slope?
(10-14 11:42)
\end{quote}

The composite waves so far are made of cosines, standing waves in which
the ± rotation terms are added with equal weight. The direction of
travel is not the direction of the slope ∂ ln\textbar R²\textbar. As the
gauge-interval calculation showed, the distortion of geodesic length is
symmetric on both sides of the peak (0.972, 1.245, 0.972); the slope ∂
ln\textbar R²\textbar{} is positive on one side of the peak and negative
on the other and points nowhere on balance. It shows where the
localisation is, not where it goes. The difference between standing and
travelling waves shows up in the imaginary part of ln R². For a standing
wave R² is real and the imaginary part of ln R² takes only the two
values 0 and π, jumping at the nodes R = 0. With only one rotation term
(xₙ = iA·e\^{}\{inφ/2\}) the true square is xₙ² = −A²·e\^{}\{inφ\}, R²
is complex,

\[
R^2=-A^2\,e^{i(N+1)\varphi/2}\,\frac{\sin(N\varphi/2)}{\sin(\varphi/2)},
\]

its modulus is localised at the peak and its argument arg R² increases
uniformly as (N+1)φ/2. The slope of the modulus (real part) carries the
shape of the localisation; the slope of the phase (imaginary part)
carries the direction and its strength. A phase slope alone is not
motion; motion needs a common parameter advancing all component phases
in proportion to n (the problem of §3.5). For a composite wave with one
rotation term only, advancing the phases in proportion to n coincides
with shifting θ₀ uniformly, and the peak translates in one direction
without changing shape.

\subsubsection{7.2 A state travelling at the speed of light and a state
travelling slowly, with no
t}\label{a-state-travelling-at-the-speed-of-light-and-a-state-travelling-slowly-with-no-t}

\begin{quote}
So, in a state with no t, can a stationary wave represent both a state
that travels at the speed of light and a state that looks as if it has
energy and travels slowly? (11:45)
\end{quote}

A background time cannot be placed, so the advance of the phase of some
reference component is used as the clock. Speed is the ratio ``how far
the peak position θ₀ shifts while the clock's phase advances one unit'',
and no t appears. When the phases of all components advance in
proportion to n, the peak shifts without changing shape, the ratio is 1,
and this is read as the speed of light. When the advance of each
component's phase departs from proportionality and goes as √(k²+m²), the
group velocity of the peak becomes k/√(k²+m²) \textless{} 1, and the
smaller the imbalance of weights the less the peak moves while only its
phase turns. That is the state ``with energy but not moving'', the rest
mass.

m comes out without being brought in. On the circle curvature cannot be
read, so the phase advance can only be proportional to n and only the
light-speed state can be represented. The eigenvalues of S³ are (n+1)² −
1 with the curvature term −1. With true squares R is imaginary, so the
space is the sphere of imaginary radius, the hyperbolic side, where the
eigenvalues are (k²+1)/\textbar R\textbar² and the gap
1/\textbar R\textbar² plays the role of m². The mass needed to travel
slowly comes from the curvature 1/\textbar R\textbar{} that the
residents can read, and since R is fixed by the wave itself through Σxₙ²
= R², the mass is made on the side of the wave. At this stage Claude
wrote that ``the assumption that the phase advance is proportional to
the square root of the eigenvalue presupposes a second-order wave
equation and has not been derived from the axiom''; this was an error,
corrected by the author in §7.8.

\subsubsection{7.3 Dividing the rotation of phase into three
steps}\label{dividing-the-rotation-of-phase-into-three-steps}

\begin{quote}
Here I want to express the rotation of phase generically: a part that
finds the rotation angle, a part that distributes it by direction, and
an array operation that actually rotates three components at once
according to the distribution. (11:48)
\end{quote}

Since S³ \(\cong\) SU(2), a rotation on S³ is written by a single unit
quaternion and corresponds directly to this division. Step 1 finds the
rotation angle (a scalar): the phase advance Δφₙ of component n is
converted to a geometric angle αₙ = Δφₙ/κₙ (κₙ the wavenumber of the
component, n/2 for cos(nφ/2) on the circle, n for cos(nψ) on S³). If the
phase advance is proportional to n all αₙ are equal and the translation
is rigid (speed of light); with a gap αₙ differs with n and the peak
lags and spreads. The difference is confined to this step. Step 2
distributes by direction (3 components): aₙ = αₙ·u with a unit vector u.
Step 3 rotates everything at once by an array operation: qₙ = cos αₙ +
sin αₙ·(u₁i + u₂j + u₃k) is multiplied from the left onto the quaternion
p of every lattice point. On S³, left multiplication by a quaternion
moves every point the same distance in the same direction with no fixed
point (a Clifford translation). A rotation of S² always has fixed
points, but on S³ this translation plays the role of translation in flat
space. Multiplication from both sides, q p q*, is a rotation that
changes orientation around the peak and is not a translation. Right
multiplication is also a fixed-point-free translation, but of a
different kind from the left one, and may correspond to the two hands of
travel (left- and right-handed).

\subsubsection{7.4 What corresponds to the Higgs field: a pinned phase
that does not
oscillate}\label{what-corresponds-to-the-higgs-field-a-pinned-phase-that-does-not-oscillate}

\begin{quote}
Here, let us suppose that what corresponds to the Higgs field is an
intrinsic internal clock. Do you see the picture? (11:50) Think
carefully. This intrinsic clock does not oscillate. It merely
distributes phase over space and is pinned there. That is exactly the
Higgs field. (11:51)
\end{quote}

Claude had treated the internal clock as ``a component of rotation
assigned to no direction'', wrote α² = α\_clock² + α\_space², said
``zero to the clock is the speed of light, all to the clock is rest'',
and added that ``the vacuum expectation value of the Higgs field in the
standard theory is a static value that does not oscillate, so Higgs
field = internal clock is a reinterpretation peculiar to this model''.
The author corrected this. The model has no t, so nothing turns. What
corresponds to the clock is a distribution of phase assigned to each
place and fixed there; it does not move and does not oscillate. The
reading of the clock is the value of the phase pinned at that place, and
differences of time appear only as differences of phase from place to
place. This is the same form as the vacuum expectation value of the
Higgs field in the standard theory, which does not oscillate and sits
statically over the whole of space; precisely because it is static it
corresponds to the Higgs field itself. When a composite wave (particle)
changes position within this background it picks up the difference ΔΦ =
∫∇Φ·dl of the pinned phase. A component that does not couple to the
background (photon-like) picks up no phase difference and passes
straight through at ratio 1; a component that couples must, each time it
changes position, rotate by the amount needed to match the pinned phase,
and by that amount its rotation in the spatial direction is reduced and
it is slower. That corresponds to mass. For a particle at rest the clock
naturally does not advance; the phase changes only when the place
changes. The origin of the pinned phase need not be placed from outside:
the arrangement of the component phases nθ₀/2, and especially the phase
of the gauge component n = 1 that fixes the origin, is a static
distribution of phase pinned to space.

\subsubsection{7.5 The beat with the static Higgs field reads out mass,
time and space at
once}\label{the-beat-with-the-static-higgs-field-reads-out-mass-time-and-space-at-once}

\begin{quote}
So the beat with the static Higgs field gives the quantity read as mass,
and the phase velocity and direction of translation, realising the
readout of mass, time and space simultaneously. (11:54)
\end{quote}

From one quantity (the difference between the phase of the static
background and the phase of the composite wave, the beat) three readouts
come out at once. Mass is the strength of the beat that remains even at
rest; as long as there is coupling to the background, the phase
difference does not vanish without a change of position. Space is the
direction and size of the slope of the beat, giving the direction of
translation and k. Time is the number of beats accumulated while
changing position; since the background clock is pinned and does not
move, the time the particle experiences (proper time) appears only as
the count of beats. The three are not independent: the square of the
strength of the beat splits as m² + k², the Pythagorean relation. A
photon-like component that does not couple to the background has no
beat, accumulates zero however many times it changes position, and its
proper time is zero.

\subsubsection{7.6 The speed of light is not exceeded: the spatial gauge
is
widened}\label{the-speed-of-light-is-not-exceeded-the-spatial-gauge-is-widened}

\begin{quote}
The speed of light is not exceeded. Have you forgotten that the spatial
gauge has widened? (11:55) Look into it. (11:56)
\end{quote}

Claude had written that ``the phase velocity of the beat is ω/k
\textgreater{} 1 and exceeds the speed of light''. That measured the
phase velocity with a flat ruler of uniform ticks---the same error as
bringing in an outside ruler. All the resident can use is the gauge of
phase, and where there is localisation one tick is longer (1.245 times
at the peak for all orders N = 18). For the same phase advance, if the
tick that measures it is widened s times, the phase velocity read with
the resident's ruler is (ω/k)/s. If this is consistent, s = ω/k = √(1 +
m²/k²) should hold, and mass m would be readable as the widening of the
gauge (the distortion of geodesic length) itself.

On verification, with the present definitions s = ω/k did not hold. s
was taken in two ways (the value of the interval containing the peak and
the value at the peak point); ω/k was taken as √(1 + 1/κ\_N²) with m =
1/\textbar R\textbar{} (the hyperbolic gap) and k =
κ\_N/\textbar R\textbar{} (the wavenumber of the highest component). For
N = 2 to 18, ω/k tends to 1 while s stays constant (converging to √2 at
the peak point) or keeps growing; the trends are opposite and do not
agree. The comparison itself was flawed: all the composite waves
examined were cosine standing waves with zero translational momentum k,
whose ω/k is properly infinite, and using the highest wavenumber in
place of k was beside the point. That the widening at the peak point
converges to √2 independently of N is a clear result in itself, because
\textbar R\textbar² = N at the peak and about N/2 on average {[}H2:
\texttt{ゲージ広がりと分散比の検証\_\allowbreak{}20261014}{]}. This failure produced
the next question.

\pandocbounded{\includegraphics[keepaspectratio]{figs/fig_gauge_stretch_vs_dispersion.png}}

\textbf{Fig. 4.} Comparison of the gauge widening s (peak interval, peak
point) with ω/k = √(1 + 1/κ\_N²). From left: circle cos(nφ/2), S³
cos(nψ), S³ Uₙ/(n+1). As N grows ω/k tends to 1 while s stays near √2 or
keeps growing; they do not agree.
\texttt{gauge\_\allowbreak{}stretch\_\allowbreak{}vs\_\allowbreak{}dispersion.py},
\texttt{plot\_\allowbreak{}gauge\_\allowbreak{}stretch\_\allowbreak{}vs\_\allowbreak{}dispersion.py}.

\subsubsection{7.7 In general relativity both time and space are
distorted: it depends on how the gauge is
fixed}\label{in-general-relativity-both-time-and-space-are-distorted-it-depends-on-how-the-gauge-is-fixed}

\begin{quote}
Let us think it over. In general relativity gravity distorts both time
and space, so it depends on how the gauge is fixed. (14:10)
\end{quote}

The previous check fixed only the spatial gauge (the nodes of the
highest component) and compared with ω/k on the premise that the time
side is undistorted. In the Schwarzschild solution rulers lengthen and
clocks slow near the source; the speed of light measured in coordinates
looks slow, but measured with the local ruler and clock it is always
c.~The resident can have two gauges: the spatial gauge, the ticks cut by
the nodes of the phase, whose geodesic length is distorted by
\textbar R\textbar; and the temporal gauge, the count of beats with the
background. What must be checked is not a one-sided equation like s =
ω/k but how the ratio of the distortions of the two gauges is fixed.
With ± rotation terms, R² changes sign from place to place; where R is
imaginary the length is imaginary, where real it is real. Two kinds of
length of different sign coming from one R² closely resembles the
distinction between timelike and spacelike intervals (the signature of
the metric). The distortions of time and space are not to be placed
separately; they are fixed together by the same Σxₙ² = R².

\subsubsection{7.8 The way of thinking of the Lorentz factor: the light
cone Σxₙ² =
0}\label{the-way-of-thinking-of-the-lorentz-factor-the-light-cone-ux3c3xux2099uxb2-0}

\begin{quote}
It depends on how you think of the Lorentz factor γ = √(1 − v²/c²). The
result is the light cone Σxn\^{}2 = 0. (14:14)
\end{quote}

Σxₙ² = 0 is the light-cone condition itself. What Claude wrote in §7.2,
that ``the dispersion √(k²+m²) is an assumption not derived from the
axiom'', was wrong: it was the axiom Σxₙ² = 0 itself. Divide the
components into three groups: R (the total budget of rotation)
corresponds to ω, the components with a direction (the share assigned to
spatial directions) to k, and the components without a direction (the
beat with the pinned background) to m. Then Σxₙ² = R² is

\[
R^2=k^2+m^2,
\]

and moving one term across as imaginary gives the light-cone form in
which the sum of squares is zero. The whole system is always on the
cone, always ``turning'' at the speed of light. The speed ratio is v/c =
k/R and 1/γ = √(1 − v²/c²) = m/R; 1/γ is the share of the rotation
budget that goes to the internal clock (the beat with the background),
and this is how proper time advances. Photon-like (m = 0): v = c, 1/γ =
0, proper time does not advance. At rest (k = 0): 1/γ = 1, the whole
budget goes to the internal clock. In between, the faster it moves the
smaller the share to the clock and the slower time runs. Time dilation
and length contraction both come from the same projection ratio 1/γ, so
the speed of light measured by the resident with the local ruler and
clock is 1 in every state. The disagreement in §7.6 arose from comparing
only the widening of the spatial gauge, without the temporal gauge, with
ω/k. Speed, mass, proper time and the constancy of the speed of light
are none of them newly placed; they are the result of reading the one
cone condition Σxₙ² = 0 by splitting it into components with and without
a direction.

\subsubsection{7.9 Checking: the light cone holds exactly on the circle
and mass is read
out}\label{checking-the-light-cone-holds-exactly-on-the-circle-and-mass-is-read-out}

\begin{quote}
Check it. (16:19)
\end{quote}

Take the components xₙ = iA·(a·e\^{}\{inφ/2\} + b·e\^{}\{−inφ/2\}) (a, b
the weights of the ± rotation terms). The true square is

\[
x_n^2=-A^2\left(a^2e^{in\varphi}+b^2e^{-in\varphi}+2ab\right).
\]

The cross term 2ab does not depend on φ and is a uniform component
pinned to space. R² = Σxₙ² is Fourier-decomposed and three quantities
are read from the field: the directional component k, the sum of the
coefficients of positive rotation minus the sum for negative rotation;
the pinned component m, the φ-independent coefficient (the cross term);
and the total R, the sum of the coefficients of both directions. The
speed is v = k/R and the clock share 1/γ = m/R. The results (N = 18) are
as follows {[}H2: \texttt{円周住民\_\allowbreak{}光円錐読出し\_\allowbreak{}20261014},
\texttt{lightcone\_\allowbreak{}readout\_\allowbreak{}s1.py},
\texttt{results\_\allowbreak{}summary\_\allowbreak{}lightcone\_\allowbreak{}readout\_\allowbreak{}s1.json}{]}.

\begin{enumerate}
\def\labelenumi{\arabic{enumi}.}
\tightlist
\item
  \textbf{The light cone holds exactly.} The residual of v² + (1/γ)² = 1
  is 10⁻¹⁵. The coefficients are Na², Nb² and 2abN, so (a² + b²)² = (a²
  − b²)² + (2ab)² holds identically and N cancels.
\item
  \textbf{The speed is fixed by the logarithm of the weight ratio.} v =
  tanh(ln(a/b)), 1/γ = sech(ln(a/b)). ln(a/b) is the rapidity.
\item
  \textbf{The two ends.} With one rotation only (b = 0), v = 1 and 1/γ =
  0: no pinned component and proper time does not advance (photon-like).
  With a = b, v = 0 and 1/γ = 1: everything is pinned component (rest).
\item
  \textbf{Velocity addition.} Composing by multiplying the weight
  ratios, the speed read from the field agrees with the relativistic
  addition law (v₁ + v₂)/(1 + v₁v₂) to six decimal places. The simple
  sum gives values exceeding 1 (1.33, 1.73).
\end{enumerate}

\pandocbounded{\includegraphics[keepaspectratio]{figs/fig_lightcone_readout_s1.png}}

\textbf{Fig. 5.} Reading Σxₙ² = R² (true squares) as a light cone on the
circle. Left: v = k/R and 1/γ = m/R read from the field lie on
tanh(ln(a/b)) and sech(ln(a/b)). Centre: all points lie on the circle v²
+ (1/γ)² = 1, with the photon-like (b = 0) and rest (a = b) states at
the ends. Right: the speed read from the field for the product of weight
ratios agrees with the relativistic addition law, while the simple sum
exceeds 1. \texttt{lightcone\_\allowbreak{}readout\_\allowbreak{}s1.py},
\texttt{plot\_\allowbreak{}lightcone\_\allowbreak{}readout\_\allowbreak{}s1.py}.

The background pinned to space that corresponds to the Higgs field did
not need to be placed from outside. It arises naturally from the true
squares as the cross term of the two counter-rotating terms.
\textbf{Mass (1/γ) is the fraction of the whole occupied by this pinned
component, and speed is the remaining directional component.} The
constancy of the speed of light and the velocity-addition law both come
out as the result of reading Σxₙ² = R² with true squares. This is the
end point of stage 1. Without a background spacetime, from the relations
and interference of the complex components of a standing wave that is
not even a travelling wave, mass has been read out.

\subsection{8. The action of phase rotation, the observer, prior work,
and z = a +
ib}\label{the-action-of-phase-rotation-the-observer-prior-work-and-z-a-ib}

\subsubsection{8.1 The action that rotates the phase: shifting θ₀ is
enough}\label{the-action-that-rotates-the-phase-shifting-ux3b8ux2080-is-enough}

\begin{quote}
Now add the action that rotates the phase, and add an on/off for this
action to the figure, on by default. (16:33) I do not see why changing
the phase directly is a problem. It is just shifting θ0, is it not?
(16:44) Moreover, the separate N for odd and even harmonics and their
on/off have disappeared. (16:46) The rotation speed is artificial and
unnecessary. Rather, put the translated distance on the slider. (16:48)
\end{quote}

The action is the three steps of §7.3 as they stand. The total rotation
dβ is advanced, distributed in the proportion of the light cone read
from the field as v·dβ to the spatial direction and dβ/γ to the internal
clock, and all components are translated by the same amount at once (θ₀
advanced by v·dβ). The internal clock's share is not drawn as an
oscillation but counted as the accumulation of beats with the pinned
background (proper time τ). Always β² = Δθ₀² + τ².

Claude had objected that ``a naive action advancing each rotation term's
phase in proportion to n splits the + and − parts in opposite directions
at the speed of light'' and wrote that ``the measured Δθ₀/β = v confirms
that the implementation moves as distributed and is not a physical
proof''. The author corrected this. Shifting θ₀ by δ changes the phase
of the + rotation term by −nδ/2 and of the − term by +nδ/2, in opposite
directions, so the composite wave translates rigidly keeping its shape,
and since the cross term 2ab of the true square does not contain θ₀ the
background does not move and only the directional component moves: the
part corresponding to the Higgs field stays pinned while only the
particle moves. What Claude called ``a problem'' was a different
procedure, multiplying both terms by the same-direction phase as in
standard time evolution; the model has no t, so there is no need to turn
all terms in the same direction, and the change of phase is sufficiently
defined as the shift of θ₀. What fraction of β the shift of θ₀ takes is
fixed by the field's own light-cone decomposition R² = k² + m² as k/R =
v, and both the action and the distribution close within the axiom.

The rotation speed is a tick for display and not a physical quantity.
The translated distance Δθ₀ is put on the slider; with the action on it
advances automatically, and dragging it back-computes β = Δθ₀/v and τ =
β/γ. Only the ratios of Δθ₀, β and τ have meaning. The maximum N of the
odd and even sides are specified separately and each can be switched on
and off {[}H2:
\texttt{circle\_\allowbreak{}phase\_\allowbreak{}rotation\_\allowbreak{}action\_\allowbreak{}lightcone\_\allowbreak{}v2.html}{]}.

\subsubsection{8.2 The four-dimensional shell version: the Clifford
translation}\label{the-four-dimensional-shell-version-the-clifford-translation}

\begin{quote}
Can you make the four-dimensional shell version of this? (17:32)　I will
test it first. (17:38)　It works. (17:44)
\end{quote}

The peak p₀ is multiplied from the left by the quaternion q = cos d +
u·sin d and moved a distance d.~Left multiplication is the
fixed-point-free Clifford translation and corresponds to the shift of θ₀
on the circle. The direction of travel u is chosen from i, j, k. v and
1/γ are fixed by ln(a/b) alone, as on the circle, and do not depend on N
or on the dimension. Colour and opacity show the imaginary part of the
amplitude, (a+b)·Σcos(nψ); the radial relief shows the real part (the
component in the direction of travel). Three assumptions were made. The
real part was taken as −(a−b)·Σsin(nψ)·cos α (α the angle to the
direction of travel seen from the peak; on the circle cos α = ±1 and
this agrees with the sine part, but the extension to S³ is not unique).
Only the direct-correspondence components cos(nψ) were used. Great
circles of S³ close at 360°, so the translated distance runs over ±360°
(corresponding to 720° on the circle) {[}H2:
\texttt{s3\_\allowbreak{}phase\_\allowbreak{}rotation\_\allowbreak{}action\_\allowbreak{}lightcone.html}{]}.

\subsubsection{8.3 There is no time, yet time evolution is
visible}\label{there-is-no-time-yet-time-evolution-is-visible}

\begin{quote}
Leaving the details for later, this is an interesting experiment, is it
not? There is no time, and yet time evolution is visible. (17:55)
\end{quote}

What enters this system is only phase, imaginary amplitude, Σxₙ² = R²,
and the action that shifts θ₀. Yet the following come out as readouts of
the field: speed (v = tanh(ln(a/b)) from the weight ratio of the ±
rotation terms), proper time (the count τ of beats with the pinned
background, which stays 0 in the photon-like state), and the constancy
of the speed of light and velocity addition (measured with the two
gauges, translated distance and beat count, every state lies on the
light cone and composition follows the relativistic law). ``Time
passes'' is written nowhere. What is seen is a fixed relation among
three quantities---translated distance, total rotation, beat count---and
that relation has the same form as what relativity calls time. What is
still placed is the parameter β that orders the sequence of states (what
β is advanced one step against has not been derived from inside) and the
tick of the animation for display. The reference for β is fixed in
Chapter 10 as the wave of the observer.

\subsubsection{8.4 Close prior work}\label{close-prior-work}

\begin{quote}
Are there prior papers close to this idea? (17:56)　Look into it.
(17:58)
\end{quote}

There is close prior work. Shanahan {[}X8-Shanahan{]} models a massive
particle as a standing wave of two counter-propagating waves and argues
that when the particle moves the two waves are Doppler-shifted in
frequency and become a travelling wave. The present model differs in
three points. Shanahan changes the frequencies of the two waves while
keeping the amplitudes equal; the present model keeps the frequencies
(harmonics n) and changes the amplitude weights a, b. On computation
a²/b² = (1+v)/(1−v) = ω₁/ω₂, so the ratio of squared amplitudes plays
the same role as Shanahan's Doppler frequency ratio. Shanahan states
explicitly that he assumed the Lorentz factor γ, whereas here 1/γ =
2ab/(a² + b²) comes out of the cross term of the true squares. Shanahan
places t and x from the start and argues for a preferred frame in which
ω₁ = ω₂; this corresponds to the state a = b in the present model, in
which only the pinned background remains.

Feynman's checkerboard model {[}X8-Feynman--Hibbs{]} and
Jacobson--Schulman {[}X8-Jacobson--Schulman{]} have the particle zigzag
at the speed of light in one dimension with mass proportional to the
frequency of reversals. For plane waves of the 1+1-dimensional Dirac
equation with right-moving component R and left-moving component L,
\textbar R\textbar²/\textbar L\textbar² = (E+p)/(E−p), v =
(\textbar R\textbar² − \textbar L\textbar²)/(\textbar R\textbar² +
\textbar L\textbar²), the ratio of the scalar density
2\textbar R\textbar\textbar L\textbar{} to the total density is m/E =
1/γ, and (\textbar R\textbar² + \textbar L\textbar²)² =
(\textbar R\textbar² − \textbar L\textbar²)² +
(2\textbar R\textbar\textbar L\textbar)²; with a = \textbar R\textbar{}
and b = \textbar L\textbar{} these are the same formulas as the
light-cone readout here. In the standard theory, too, the mass term ψ̄ψ
is the cross term of the left- and right-handed components and the Higgs
vacuum expectation value multiplies it. In Bondi's k-calculus v = (k² −
1)/(k² + 1), and the present v = (a² − b²)/(a² + b²) is the same formula
with k = a/b. Page--Wootters {[}X3-Page--Wootters{]} derived time
evolution as internal correlation within a globally stationary entangled
state. De Broglie's internal clock and Hestenes's zitterbewegung tie
mass to an internal frequency.

What overlaps with prior work is that mixing two light-speed components
yields speed, Lorentz factor and mass; this is established as
1+1-dimensional Dirac theory and no novelty is claimed for it. What is
not in prior work is reaching that structure from the closed circle and
the true-square condition Σxₙ² = R² alone, without spacetime, t or a
Hamiltonian; the localisation born of harmonic superposition and
narrowed to one point by the balance of odd and even; and translation
expressed on S³ as a Clifford translation.

\subsubsection{8.5 z = a + ib: a and b are both
real}\label{z-a-ib-a-and-b-are-both-real}

\begin{quote}
a and b are real: z = a + ib. Both are real. (18:00)
\end{quote}

Claude had raised the question ``if a, b are given phases, the true
square a² and the conjugate product \textbar a\textbar² no longer
agree'', but it was not a question that stands. a and b are real from
the start, and treated as a complex number they are the real and
imaginary parts of one number z = a + ib. The true square of z is

\[
z^2=(a^2-b^2)+i\cdot2ab,
\]

whose real part is the directional component k, imaginary part the
pinned component m, and modulus \textbar z²\textbar{} = a² + b² = R. The
light-cone decomposition R² = k² + m² is the relation among real part,
imaginary part and modulus of the single complex number z². Writing z =
\textbar z\textbar·e\^{}\{iφ\}, v = cos 2φ and 1/γ = sin 2φ; in the
photon-like case (b = 0) φ = 0 and z² is real, at rest (a = b) φ = 45°
and z² is purely imaginary. What fixes the speed is not a phase brought
in from outside but the argument of z itself, and what enters the
imaginary part of z² is the clock (mass) and what enters the real part
is the advance in the spatial direction. The bilinear relations of Dirac
theory are gathered in this model into the real part, imaginary part and
modulus of z².

\subsection{9. Stage 2: why three dimensions? Because three dimensions
has the
symmetry}\label{stage-2-why-three-dimensions-because-three-dimensions-has-the-symmetry}

\subsubsection{9.1 Neither time nor space is assumed. Only the
three-dimensional reading is
assumed}\label{neither-time-nor-space-is-assumed.-only-the-three-dimensional-reading-is-assumed}

\begin{quote}
This experiment is strange. Not only time but space, too, is not
assumed. R is merely interpreted as space. What I am dissatisfied with
is that only the three-dimensional reading is assumed. (18:02)
\end{quote}

In this experiment not only time but space has not been placed. There is
only the phase on the closed ring, the imaginary amplitude and the
true-square condition Σxₙ² = R²; reading R or the translated distance as
``space'' is only an interpretation. ``Three dimensions'' alone has been
placed: at the moment S³ was chosen and at the moment the directions of
translation were taken as the three quaternion units i, j, k.

This three can be narrowed down from the condition on the action. The
translation of step 3 is an action in which ``every point moves the same
distance in the same direction and no point stays fixed'', and two
translations performed in succession must compose into one (velocity
addition). This is the condition that the space itself can move
uniformly over itself as a group. Among spheres, only S¹ and S³ can
carry such a group structure. On S² any flow necessarily has a
stationary point (the hairy-ball theorem). On S⁷ something like a
uniform translation can be made with the octonions, but the octonions
are not associative, so two translations do not compose into one. Higher
spheres admit no uniform translation at all. The structure in which the
light cone closes under the true square of z = a + ib is subject to the
same restriction: the number systems in which squaring preserves the
size relation are only the reals, complexes, quaternions and octonions
(Hurwitz's theorem {[}X1-Hurwitz{]}), with 0, 1, 3, 7 directions, and of
these only 1 (complex, S¹) and 3 (quaternions, S³) are also associative.
From the condition on the action alone---``on a closed space, rotations
of phase can be composed as uniform translations''---the number of
directions is limited to 1 or 3. Three dimensions is not an assumption
but can be read as the largest dimension in which composition works.

\subsubsection{9.2 In other than three dimensions the system is not
stable: it cancels cleanly at the symmetric
point}\label{in-other-than-three-dimensions-the-system-is-not-stable-it-cancels-cleanly-at-the-symmetric-point}

\begin{quote}
Exactly. In fact, from the four-dimensional symmetry and the special
nature of the three-sphere it should be derivable that the system is not
stable in other than three dimensions. Only in three dimensions does it
conveniently vanish cleanly at the symmetric point. (18:03)
\end{quote}

The sign (−1)ⁿ at the antipode alone holds in any dimension. The
eigenfunctions on Sᵈ symmetric about the peak are Gegenbauer polynomials
Cₙ\^{}λ(cos ψ) (λ = (d−1)/2) with Cₙ\^{}λ(−1) = (−1)ⁿ·Cₙ\^{}λ(1), so
normalised, the antipodal value is ±1 in every dimension. That ``odd and
even cancel at the antipode when equal in number'' happens on S² and S⁵
too. What is peculiar to S³ is that they cancel ``cleanly''. On S³, λ =
1 and the eigenfunctions are Uₙ(cos ψ) = sin((n+1)ψ)/sin ψ---the sine
wave of the circle divided by sin ψ---and the eigenvalues have the form
``square minus one'', (n+1)² − 1. A spherical wave on S³ propagates as a
one-dimensional wave f(ψ ± β) divided by sin ψ without changing shape
(Huygens' principle holds exactly and no tail is left behind), so a wave
leaving the peak refocuses sharply at the antipode, where the odd and
even components cancel without excess or deficiency. On S² (λ = 1/2,
Legendre polynomials) Huygens' principle fails, the wave spreads with a
tail, the cancellation at the antipode is incomplete and the remainder
accumulates in the system. This is the concrete content of ``in other
than three dimensions the system is not stable''.

\subsubsection{9.3 A higher dimension with the same symmetry:
S⁷}\label{a-higher-dimension-with-the-same-symmetry-sux2077}

\begin{quote}
So, since why it becomes three dimensions can be explained, it can no
longer be said to be contrived. What is curious here is whether there is
a higher dimension that in fact has the same symmetry. (18:05)
\end{quote}

There is: S⁷. S¹, S³ and S⁷ are the only spheres on which a flow with
aligned direction can be laid over every point (Adams' theorem,
parallelisability), and left multiplication by a unit octonion is an
isometry of S⁷ with no fixed point. There is a Hopf fibration S³ → S⁷ →
S⁴ of the same form as S¹ → S³ → S². The octonions are one of Hurwitz's
four and the structure that closes the light cone from the real part,
imaginary part and modulus of z² can be built directly (7 directions).
Being odd-dimensional, Huygens' principle holds and the refocusing at
the antipode is sharp. What S⁷ lacks is associativity: (pq)r ≠ p(qr), so
two successive translations do not make one, left multiplications do not
close into a group, and velocity addition cannot be kept. The simplicity
of the S³ eigenfunctions, sin((n+1)ψ)/sin ψ, is also lost. S⁷ reaches to
within one step of building a space with uniform translation and a light
cone, and breaks at composition. As a conjecture: S³, where composition
closes, becomes the ``space'' in which particles translate; S⁷, where it
does not, cannot be a space but part of its symmetry survives as an
internal degree of freedom. This agrees in direction with work (Furey
{[}X7-Furey{]}) that uses the fact that the subgroup of the octonion
automorphism group G₂ fixing one imaginary unit is SU(3) to tie the
octonions to the colour symmetry.

\subsubsection{9.4 Only in three dimensions do conditions A, B and C all
hold}\label{only-in-three-dimensions-do-conditions-a-b-and-c-all-hold}

\begin{quote}
Conditions A and B hold for S3 alone, do they not? (18:31) Shall we
correct it to: A, B and C hold only in three dimensions? (18:31)
\end{quote}

The sphere-resident paper {[}H1 §9{]} narrows to S³ by three conditions.
Condition A (uniform translations form a group) is satisfied by S¹ and
S³. Condition B (Huygens' principle, propagation without a tail) by the
odd-dimensional spheres S¹, S³, S⁵, \ldots. Condition C (a mass gap
arises from curvature the residents can read, and left and right
translations differ so that handedness arises) by all but S¹. All three
of A, B and C hold only for S³, three-dimensional space. A localised
massive particle can exist stably only in three dimensions. In fixing
paper v1, Claude made three self-corrections: a note that the early
self-consistent-shell results are not adopted (§5); the antipodal ratio
in the table of §6, 0.972 → 0.988 (0.972 was the interval next to the
peak); and the distinction that Adams' theorem says only ``the spheres
with a uniform flow are S¹, S³, S⁷'' while ``those that compose as a
group are S¹ and S³'' is a separate fact (§9).

\subsection{10. The observer is inside the system: a resident of the
circle, not a god's-eye
view}\label{the-observer-is-inside-the-system-a-resident-of-the-circle-not-a-gods-eye-view}

\begin{quote}
The remaining issue is β. How is β being used now? (18:32) Is it
something to agonise over? β is, as its name says, the observer's
observing wave. Here is the point: where is the observer? A resident of
the circle, or a god's-eye view? (18:34) Exactly. This allows the
experiment of special relativity. When the observer moves at some speed
the visible result changes. (18:36) Does it not affect the
three-dimensional version too? (18:48)
\end{quote}

β was the total rotation, advanced one step at a time in step 1 of the
action. At a = b (rest) β = τ, and in general β = γτ; β played the role
of the time of an observer at rest with respect to the pinned
background, coordinate time. Claude was ticking it from outside and, as
a candidate for fixing it from inside, named ``the proper time of
another composite wave'', tying it to the unsolved problem of the
coexistence of several localisations. The author answered. There is
nothing to agonise over. β is the observer's observing wave. The
observer is a resident of the circle; the system has no outside and so
no god's-eye view. β is the advance of the phase of the observer's own
wave, its proper time, and the observer measures the progress of the
other wave against its own wave as a ruler.

That a = b looks like rest is because the observer had tacitly been
placed at a = b; the pinned background, too, is not a special preferred
frame but the component at rest as seen by the observer. With the
observer's rapidity η\_obs = ln(a\_obs/b\_obs) and the other's η, the
speed the observer reads is v' = tanh(η − η\_obs); changing the observer
is multiplying weight ratios, which is the relativistic addition law
already confirmed. Advancing β by the browser clock was ticking with a
god's eye; that alone had brought in an outside.

Adding one slider for the observer's weight ratio ln(a\_obs/b\_obs)
changes three things: the visible speed tanh(η\_particle − η\_obs), the
visible advance of proper time sech(η\_particle − η\_obs), and the
ticking of β (the phase of the observer's own wave). Matching the
observer to the particle's speed makes the particle appear at rest with
maximal proper time; moving the observer the other way makes the speed
look larger and the proper time shrink; two observers measuring the same
particle read different speeds but both lie on the light cone. The laws
of physics have the same form for every observer and there is no
preferred observer. This is special relativity coming out of the model
without bringing in an outside {[}H2:
\texttt{sr\_\allowbreak{}observer\_\allowbreak{}lightcone\_\allowbreak{}s1.html}{]}.

On S³ too the observer is a resident of the same sphere, exists as a
wave localised on the same sphere as the particle, and its peak position
(χ\_obs, θ\_obs, φ\_obs) can be specified independently. The speed
readout is fixed by the difference of rapidities tanh(η\_particle −
η\_obs) alone and does not depend on where in space the two peaks are.
That every resident of S³ can use the same law is inherent in the
structure of this readout. v'² + (1/γ')² stays 1.000000 however the
sliders are moved {[}H2: \texttt{sr\_\allowbreak{}s3\_\allowbreak{}observer\_\allowbreak{}lightcone.html}{]}.

\subsection{11. Stage 10: can spin be read
out?}\label{stage-10-can-spin-be-read-out}

In the author's eleven stages this is stage 10 (after the
four-dimensional lattice), but in the dialogue it was asked right after
three dimensions (10-14 18:54) and is read out directly from the
geometry of S³, so the paper places it here. The two spin components on
the four-dimensional lattice reappear in §20.3.

\subsubsection{11.1 Spin is already in the
model}\label{spin-is-already-in-the-model}

\begin{quote}
With this, in a sense, we can express the photon and a particle that has
mass but no charge. Does it not concern you where spin is read from?
(18:54)
\end{quote}

The odd/even distinction of Chapter 2 is the readout of spin as it
stands. An odd component reverses sign in one lap (360°) and returns
only after two (720°). That is the definition of spin 1/2. An even
component returns in one lap: integer spin. Since S³ \(\cong\) SU(2),
the three generators i, j, k of left quaternion multiplication are the
spin operators Sₓ, Sᵧ, S\_z, and the axis chosen as ``direction of
travel i / j / k'' in the figure is both the direction of spatial
translation and the axis along which spin is oriented. Moving the peak
from the left by q = cos d + u·sin d, at d = π (360°) q = −1 and p → −p
jumps to the antipode, where the odd-component wave has reversed sign
and is not back in its original state; only at d = 2π (720°) does q =
+1. Left multiplication (SU(2)\_L) and right multiplication (SU(2)\_R)
are independent and together give SU(2)\_L × SU(2)\_R \(\cong\) SO(4).
On S¹ there was no such distinction; on S³ two kinds of hand appear for
the first time. One may correspond to spatial rotation (spin) and the
other to a different internal degree of freedom; since in the standard
theory the weak interaction has an SU(2) symmetry distinguishing left-
from right-handed, the right SU(2) is a candidate. The procedure for
reading out spin is three steps: whether odd components are included
decides spin 1/2 or integer; which axis (i, j, k) generates the
translation decides the orientation; which of the two SU(2)s, left or
right, is used decides whether it is spin or another internal quantum
number.

\subsubsection{11.2 Mixing odd and even harmonics makes the spin
indefinite}\label{mixing-odd-and-even-harmonics-makes-the-spin-indefinite}

\begin{quote}
This is where it is strange: if the even and odd harmonics are mixed,
what happens to the spin? (18:56)
\end{quote}

Multiplying from the left by q = −1 (a 360° translation), the even
components cos(nπ) = +1 are unchanged and the odd ones −1 are reversed.
The mixed state after 360° is a different state in which the even part
is unchanged and only the odd part is reversed---neither the original
nor minus the original. Only after 720° do all components return. In the
representation theory of SU(2), integer spin (even) and half-integer
spin (odd) belong to different irreducible representations, and the
mixed state belongs to neither and has no definite spin. In Chapter 6 it
was confirmed that localisation narrows to one point only when odd and
even are balanced; but then the spin is a superposition of irreducible
representations and indefinite. Conversely, making the spin definite
(purely odd or purely even) destroys the localisation at one point. In
nature no particle exists in a superposition of half-integer and integer
spin (a superselection rule). In this model that prohibition appears
geometrically as ``fixing the position at one point makes the spin
indefinite''.

\subsubsection{11.3 Why the Higgs wave has spin
0}\label{why-the-higgs-wave-has-spin-0}

\begin{quote}
And do you see why the Higgs wave of this waveform has spin 0? (18:57)
\end{quote}

The reason fits in one line: the cross term is the product of the + and
− rotation terms. Taking the true square of xₙ = iA(a·e\^{}\{inφ/2\} +
b·e\^{}\{−inφ/2\}), the cross term is a·e\^{}\{inφ/2\} ×
b·e\^{}\{−inφ/2\} = ab. The + rotation term carries angular momentum
+n/2 and the − term −n/2, and their product carries 0. The same holds
for every n, so the cross term 2ab·N is a constant depending on neither
φ nor ψ. A constant function on S³ is unchanged by every element of
SU(2)\_L × SU(2)\_R: the trivial representation, the definition of a
spin-0 scalar field. The reason the Higgs field has spin 0 in the
standard theory is essentially the same: the left- and right-handed
components are combined into a scalar whose angular momenta cancel. In
this model it comes out automatically from the structure that ``the +
and − rotation terms live together on the same circle''. The particle
wave (odd + even mixed) is localised but spin-indefinite; the Higgs-like
background (cross term) is pinned to the whole of space with spin 0
definite; when the particle moves through the background the beat with
it gives proper time and produces mass. Because the partner of the beat
has spin 0, the particle's spin is unchanged, in agreement with the
property of the standard theory that acquiring mass through the Higgs
mechanism does not change spin.

\subsubsection{11.4 Odd harmonics alone are also localised to some
extent: two-point
localisation}\label{odd-harmonics-alone-are-also-localised-to-some-extent-two-point-localisation}

\begin{quote}
You misread. In fact, the odd harmonics alone also look localised to
some extent. (19:00)
\end{quote}

Claude's wording was inaccurate. A composite wave of odd components
alone has sharp peaks in two places, the peak (+i) and the antipode
(−i). This is not ``not localised'' but ``localised at two points''.
That the two peaks stand in a sign-reversed relation is precisely the
geometric form of spin 1/2. A spinor reverses sign under a 360°
rotation, and in space too this appears as ``the peak and the antipode
form a ±i pair''. Odd only: sharply localised at two points with
reversed signs (spin 1/2 definite). Even only: localised at two points
with the same sign (integer spin definite). Odd and even balanced: the
antipodal peak cancelled, localisation at one point (spin indefinite).
The correct form of the ``position--spin uncertainty'' is that
localisation at one point and a spin eigenstate are incompatible. A spin
eigenstate is localised at two points; to localise at one point odd and
even must be mixed, and then the spin is indefinite.

\subsubsection{11.5 The opposite spin on the far side of the world. But
no one can observe
it}\label{the-opposite-spin-on-the-far-side-of-the-world.-but-no-one-can-observe-it}

\begin{quote}
But this does not contradict the localisation of the particle, does it?
Strangely, there is an opposite spin on the far side of the world?
Strangely, no one can observe it. (19:01)
\end{quote}

It does not contradict. The particle is localised at θ₀. The antipodal
peak exists, but from the resident's point of view reaching it needs a
geodesic distance πR, and if S³ is the whole universe the antipode lies
beyond the observable horizon. The observer, too, is a resident of the
same S³ with a wave localised at θ₀, and the beat readout has meaning
only where two waves share the same phase, so as long as the observer's
wave is localised at θ₀ the information at the antipode cannot be read.
In a spin-1/2 eigenstate of odd components alone, a +i peak stands at θ₀
and a −i peak at the antipode. This is the structure ``whenever a
particle is created, a pair with opposite spin is created on the far
side of S³'', and it can be read as the geometric origin of the
antiparticle: the positive- and negative-energy solutions of the Dirac
equation realised as the peak and the antipodal peak on S³. Adding even
components to localise at one point cancels the antipodal peak, and this
erasing operation is the same as the operation that gives mass (produces
the Higgs-like background). To have mass and to erase the shadow on the
far side are one and the same event. The photon-like state (spin
definite, a shadow at the antipode, zero mass) and the massive particle
(a spin-indefinite superposition, the antipodal shadow erased, mass
present) can be drawn in one figure as this contrast.

\subsection{12. Stage 3: can charge be
read?}\label{stage-3-can-charge-be-read}

\begin{quote}
Now, you think charge cannot be read, do you not? Is that really so?
(19:03) Probably that is charge. (19:05)
\end{quote}

It can. Moreover it is already in the model. S³ carries the natural Hopf
fibration S¹ → S³ → S². In the representation (z₁, z₂) = (cos
η·e\^{}\{iξ₁\}, sin η·e\^{}\{iξ₂\}) the sum (or difference) of the
phases ξ₁ and ξ₂ goes round this S¹ fibre. The winding number---how many
times the wave xₙ goes round this fibre---is the charge. The component
e\^{}\{inφ/2\} winds the fibre +n/2 times and e\^{}\{−inφ/2\} −n/2
times.

{\def\LTcaptype{none} % do not increment counter
\begin{longtable}[]{@{}lll@{}}
\toprule\noalign{}
State & Fibre winding number & Charge \\
\midrule\noalign{}
\endhead
\bottomrule\noalign{}
\endlastfoot
Purely + rotation terms (b = 0) & +Σn & Positive \\
Purely − rotation terms (a = 0) & −Σn & Negative \\
Balanced standing wave (a = b) & 0 & Neutral \\
Cross term (Higgs-like background) & 0 (does not wind the fibre) & 0 \\
\end{longtable}
}

That the cross term has charge 0 is consistent with its spin 0 in §11.3.
The quantisation of charge comes from the structure that was in the E
function (§2.8): when Im(E) jumps by π at each node, the total winding
of the phase per lap is an integer. That integer is the origin of charge
quantisation, not something placed from outside but a topological
invariant, the winding number of phase on a closed space. When a ≠ b and
the cross term (Higgs-like background) appears, the symmetry SU(2)\_L ×
SU(2)\_R is broken at the same time and the remaining U(1) is the
electromagnetic U(1); the eigenvalue of the generator of this remaining
U(1) is the fibre winding number, the charge. To establish the readout
of charge, the fibre winding number must be distinguished from the
rapidity imbalance k = P₊ − P₋. k is the information ``which way it
moves'' (translation) and the winding number the topological information
``how many times it winds the fibre''; at rest (a = b) k = 0 but the
winding number may be non-zero, and that is the key to the distinction.
The author's verdict was ``probably that is charge''.

\subsection{13. Observation is the cross term: being able to read,
radiating, and distributing energy. Uncertainty and
quantisation}\label{observation-is-the-cross-term-being-able-to-read-radiating-and-distributing-energy.-uncertainty-and-quantisation}

\subsubsection{13.1 How could the measuring apparatus
accelerate?}\label{how-could-the-measuring-apparatus-accelerate}

\begin{quote}
Another curious thing: you placed the measuring apparatus C for readout,
but how is it able to accelerate? Neither gravitational nor
electromagnetic waves can be observed without radiation. Have you taken
account of the interaction with the apparatus? The apparatus, too,
changes R. (19:08)
\end{quote}

The implementation of Chapter 10 had placed three things from outside.
The observer's rapidity was moved by a slider (acceleration by the hand
of god). The observer's wave and the particle's wave were drawn
separately and R computed separately. Observation was done as a
computation of the rapidity difference, with no physical process that
makes that readout possible. The axiom Σxₙ² = R² is a condition on the
whole system and the observer must enter the same sum. Taking the true
square of the whole with the particle's components x\^{}(p) and the
observer's x\^{}(o),

\[
\sum\left(x^{(p)}+x^{(o)}\right)^2=\sum\left(x^{(p)}\right)^2+\sum\left(x^{(o)}\right)^2+2\sum x^{(p)}x^{(o)},
\]

and the cross term 2Σx\textsuperscript{(p)·x}(o) is the interaction, and
is observation itself. It contains the terms a\_p·b\_o and b\_p·a\_o
whose ratio is e\^{}\{η\_p − η\_o\}; the relative rapidity the observer
reads is the quantity that should properly be read from this cross term.
Changing the observer's a\_o/b\_o changes the weights of the ± rotation
terms; in a closed system that can only happen through the cross term
with the particle or the background, and when the weights change the
shared R² changes, and that change is radiation. Without radiation there
is no acceleration and without the cross term no observation. Adding the
observer's components changes R²; when ``the observer measures the
particle'', seen from the particle ``the observer's wave changed my R''.
The interaction is symmetric and observation is not one-sided.

\subsubsection{13.2 Being able to read means energy has been
distributed}\label{being-able-to-read-means-energy-has-been-distributed}

\begin{quote}
No need to experiment on this yet, but in the end, being able to read
means that an electromagnetic or gravitational wave must have been
emitted; in a sense energy is distributed. (19:09) This is an
interesting next task: charge and spin, and further the emission and
absorption of electromagnetic and gravitational waves, and from them
recoil and even attraction, may all become drawable. (19:16)
\end{quote}

R² is constant by the axiom, but its distribution changes. The observer
can read the particle because the cross term is non-zero, and by that
amount R² is subtracted from the self-terms of the particle and the
observer. The moment a readout happens, the energy is distributed three
ways: the particle's self-term (its mass and motion), the observer's
self-term, and the cross term (interaction = radiation). If the cross
term is zero nothing can be read; if something can be read the cross
term is necessarily present and that part has been emitted as radiation.
To read and to radiate are one and the same event. Expanding the cross
term gives two kinds of terms. The position-dependent terms (products of
same-direction rotations such as a\_p·a\_o·e\^{}\{in(2φ−φ\_o)/2\}) ride
on the spatial phase and propagate---the radiated wave (the
electromagnetic part). The position-independent terms (products of
opposite rotations such as a\_p·b\_o, b\_p·a\_o) are beat terms
depending on the relative rapidity η\_p − η\_o---the very speed the
observer reads. Since R was the reading of space itself, when the
observer reads the particle the distribution of R changes and so does
the metric of space: this is the gravitational aspect of the
interaction. The part in which the Hopf-fibre winding (charge) is
exchanged through the cross term is the electromagnetic aspect. Both are
inside the same cross term, and which part is read decides whether it
looks like a gravitational wave or an electromagnetic one. In standard
quantum theory the disturbance of the object by measurement is a
separately placed principle; in this model, once Σxₙ² = R² is imposed on
the whole system, reading cannot help being a redistribution of R, and
the back-action is a consequence of the axiom.

The next tasks are five: emission and absorption (the movement of the
position-dependent part of the cross term), the distinction between
gravitational and electromagnetic waves (redistribution of R versus
exchange of fibre winding), recoil (the change in the ratio of k to m
after emission), attraction (the distance dependence of the cross term's
distribution), and the contribution of spin and charge to the
interaction. All of them are expected to be reachable with no new
assumption, by imposing Σxₙ² = R² on the whole system and reading the
cross term. The numerical experiments are left to the next stage. In
Chapter 21 it turns out that this task is inherent in the system as the
complete absorbing shell at infinity of the Eleventh Thought Experiment.

\subsubsection{13.3 Uncertainty and quantisation are consequences of the
axiom}\label{uncertainty-and-quantisation-are-consequences-of-the-axiom}

\begin{quote}
And do you not already see the possibility of uncertainty and of the
discretisation and quantisation of phase? (19:20)
\end{quote}

Both are already visible. Three kinds of uncertainty have appeared.
Position and harmonic number n (Chapter 2: a single n has uniform
amplitude and no position; many n localise to width πR/(N+1); Δθ·Δn ≈
constant; reading n as momentum, this is Heisenberg's position--momentum
uncertainty itself, coming automatically from Fourier conjugacy).
Position and spin (Chapter 11: one-point localisation needs odd and even
mixed, and then the spin is a superposition of irreducible
representations; a spin eigenstate is a two-point localisation; a form
not found in standard quantum theory). Readout and back-action (§13.1:
if something can be read there is a cross term, and if there is a cross
term R has been redistributed; there is in principle no way to read
without disturbing).

Five kinds of quantisation have appeared. n ∈ ℤ (closure on the circle).
Winding number ∈ ℤ (the Hopf-fibre winding, the charge, is a topological
integer). Discretisation of phase (in a standing wave Im(E) takes only
the two values π/2 and 3π/2). The closure condition (component n
accumulates phase 2π·(n/2) per lap and closes in one lap if even, two if
odd; \(\oint\)(phase) = 2π × integer is the Bohr--Sommerfeld
quantisation condition itself). The discreteness of readout (position
can only be read as the gauge interval, proper time only as the number
of beats; a continuous θ₀ or τ is invisible to the resident, who reads
only which interval and how many beats). As a mere suggestion: the
rapidity a/b is continuous at this stage, but once the observer is put
into the same Σxₙ² = R², the cross terms e\^{}\{i(n\_p ± n\_o)φ/2\} must
also satisfy the closure condition, from which a discrete constraint on
a/b may follow. If so, relative velocity would be quantised---a
prediction absent from the standard theory. Uncertainty follows from
``readout can only be a ratio'', quantisation from ``on a closed space
the phase must close'', both as consequences of the axiom, with nothing
placed as a separate principle.

\subsubsection{13.4 Fixing what has been reached so
far}\label{fixing-what-has-been-reached-so-far}

\begin{quote}
Tidy it up. (19:22)　Correct it and save. (19:24)
\end{quote}

The development to this point was organised as the sphere-resident paper
v1 {[}H1{]} with an abstract and 14 sections and saved to Drive
(2026-10-14 19:54). The sections are: the setting of the problem / the
circle-resident model / the exponential representation and the E
function / the extension to S³ / the self-consistent shell and Σxₙ² = R²
/ the gauge interval and the non-uniform geodesic length / the
light-cone readout / the action of phase rotation and time evolution /
why three dimensions / the readout of spin and the spin 0 of the Higgs
field / the readout of charge / uncertainty and quantisation as
consequences of the axiom / comparison with prior work / conclusion and
future tasks. Chapters 2--13 of this paper rewrite the content of that
paper, giving priority to the order of the dialogue and the author's
corrections.

\subsection{14. Stage 4: can it be read as the
electron?}\label{stage-4-can-it-be-read-as-the-electron}

From this chapter to Chapter 24 we follow the development of the
dialogue ``The dialogue after the Fourteenth Thought Experiment''
(2026-10-14 20:29 to 10-15 07:32) {[}A2{]}. By Chapter 13 the photon, a
massive particle without charge, and a candidate for charge (the
Hopf-fibre winding number) had been read out. The next question is how
this system reads the neutron, baryons, quarks, and the electron.

\subsubsection{14.1 The weak force appears as a redistribution of
harmonic components. The lifetime of the
neutron}\label{the-weak-force-appears-as-a-redistribution-of-harmonic-components.-the-lifetime-of-the-neutron}

\begin{quote}
The weak force appears when the flavour changes, so in this system it
should look like a redistribution of harmonic components. If the strong
force is thought of plainly, the neutron is a heavy particle without
charge, and it has a finite lifetime, does it not? In the case of a
single particle? (10-14 20:29)
\end{quote}

A free neutron has a finite lifetime, a mean of about 880 seconds,
decaying as n → p + e⁻ + ν̄ₑ. What causes this decay is the weak force:
one d quark turns into a u quark, a change of flavour, the
representative example of a phenomenon that should appear in this system
as ``a redistribution of harmonic components''. A neutron in a nucleus
does not decay because decay would raise the total energy. Claude here
placed the strong force as ``the condition that locks the components
into one closed bundle (confinement)'', the weak force as ``the process
in which the distribution of harmonic components inside the bundle
shifts to lower energy'', and the lifetime as ``whether the destination
of the shift is energetically open''. This placement is a connected
part, a hypothesis.

\subsubsection{14.2 Baryons and the stability of the proton, the
harmonics of quarks, the 1/3 charge, the energy
hierarchy}\label{baryons-and-the-stability-of-the-proton-the-harmonics-of-quarks-the-13-charge-the-energy-hierarchy}

\begin{quote}
So it will do if the baryons, the composite particles of quarks, can be
explained. The point is why the proton is stable. Quarks are indivisible
elementary particles, but in this system they clearly have harmonics.
The energy hierarchy and the 1/3 charge are also keys. (20:32) Are the
order and amplitude integers? Besides, oscillating between integers
looks like a continuous value. (20:34) Probably there is a relation that
is near an integer but cannot lock, which makes a potential wall that
prevents falling to the ground state. (20:36) So E8 symmetry may come
out. (20:38) What does it mean that three generations of fermions cannot
be embedded chirally? (20:40) Probably it oscillates in the irrational
relations among several pull-ins---the pull-in of the E8 symmetry,
three-phase current, probably also the pull-in of 23/124, the crystal
structure of four-dimensional spacetime---so that the energy hierarchy
looks irrational. (20:43)
\end{quote}

In this system ``indivisible'' and ``having harmonics'' are compatible.
Harmonics are an internal structure that cannot be taken out singly,
consistent with confinement. The order of a harmonic is an integer but
its amplitude need not be, and if the state oscillates among integer
states the observed value looks continuous as a time average. That quark
masses change with the scale of measurement (running masses), that
flavour states are superpositions of mass states (CKM mixing), and
neutrino oscillation are not in conflict with this picture. The task is
therefore not ``integer or not'' but ``what fixes the weights of the
oscillation''.

A relation that is near an integer ratio but cannot lock makes a
potential wall that prevents falling to the ground state. Claude placed
this in correspondence with the theory of synchronisation (the Adler
equation dφ/dt = Δω − K sin φ, a tilted washboard potential) and noted
that the exponential dependence of the escape rate on the wall height
(Kramers) could correspond to the spread of particle lifetimes over many
orders (proton \textgreater{} 10³⁴ years, neutron 880 s, muon 2.2 μs).
On E₈ it cited Zamolodchikov's result that the mass ratios of
excitations in the one-dimensional Ising model near its critical point
follow the E₈ structure (m₂/m₁ = the golden ratio), confirmed
experimentally by Coldea et al.; that projecting the 240 roots of E₈ to
four dimensions yields the H₄ symmetry containing the golden ratio; and
that Lisi's E₈ theory was shown by Distler--Garibaldi to be unable to
embed even one generation chirally (the representations of E₈ are real,
so every left-handed particle comes with a right-handed partner). The
author's view was that the energy hierarchy looks irrational because the
system oscillates in the irrational relations among several
pull-ins---the pull-in of E₈, three-phase current, the pull-in of
23/124, the crystal structure of four-dimensional spacetime. That an
integer structure in a higher dimension looks irrational when read in a
lower one has precedents in quasicrystals (cut-and-project from a
higher-dimensional integer lattice) and in the projection E₈ → H₄.

\begin{quote}
Read this (the note on the Fourteenth Thought Experiment) and the
original paper and the meaning will be clear. The arrangement of the
electron's mass is certainly involved. (20:45)
\end{quote}

After reading the Fourteenth Thought Experiment {[}A1{]}, the mass
arrangement of the charged leptons was checked (E0). Koide's relation
(mₑ + m\_μ + m\_τ)/(√mₑ + √m\_μ + √m\_τ)² = 2/3 is, for v = (√mₑ, √m\_μ,
√m\_τ) split into its (1,1,1) component and the perpendicular component,
equivalent to \textbar v⊥\textbar² − \textbar v∥\textbar² = 0: the
light-cone condition of a zero square sum Σx² = 0 with (1,1,1) as the
imaginary axis. The central projection onto the perpendicular plane
gives three phases 120° apart, and the amplitude √2 is fixed by this
zero condition. Imposing Q = 2/3 exactly and solving for the τ mass
gives 1776.969 MeV (measured 1776.93 ± 0.09), and the three-phase angle
is δ = 2/9 − 1.75×10⁻⁷ rad. The uncertainty of δ from the electron and
muon masses is about 4×10⁻¹⁰, so this offset is about 450 times the
error. It is extremely close to 2/9 yet not equal, and 2/9 rad is not a
rational multiple of π, so it never recurs finitely. This is precisely a
``pulls in but never locks completely'' pair of the Fourteenth Thought
Experiment. Nambu's m\_μ/mₑ ≈ 3/(2α) (1 + 3/(2α) = 206.55, measured
206.77) is also a near-but-not-equal pair {[}A2 §3,
\texttt{te14\_\allowbreak{}E0\_\allowbreak{}koide\_\allowbreak{}phase\_\allowbreak{}check.py}{]}.

\subsubsection{14.3 Try it with the electron:
E1}\label{try-it-with-the-electron-e1}

\begin{quote}
Why not try it with the electron, which is free of the strong force and
little affected by the weak force? (20:49) Go ahead. (20:49)
\end{quote}

The electron is the best testing ground. It feels no strong force and
can be taken out singly and its mass measured directly. Being the
lightest charged lepton it has nowhere to decay, with a lifetime above
10²⁸ years, ``the lowest valley of its series'' like the proton among
baryons. The weak force acts only in the decays of the muon and τ to the
electron. Its mass is measured to about 3×10⁻¹⁰ and its g-factor to
10⁻¹³ relative precision. With the electron series, the
``zero-square-sum framework plus a near-but-unlocked phase'' of the
Fourteenth Thought Experiment can be isolated with the strong and weak
forces almost removed.

The design of E1 uses the method of the Fourteenth Thought Experiment as
it stands. The state is the three-phase vector (its amplitude fixed by
the zero square sum, not an input). The two constraints are the rotation
orders allowed by the four-dimensional lattice, \(\mathcal{C}_A\) =
\{1,2,3,4,5,6,8,10,12\}, and those allowed by integer phases,
\(\mathcal{C}_B\) = \{1,2,3,4,6\} (Z₄, Z₆). The update is of circle-map
form with no external time. The readouts are the dwell-time distribution
of the phase δ and the frequency of phase slips. The falsification
conditions fixed in advance: if δ settles on a rational multiple of π
(locks), the hypothesis fails; whether it comes near 2/9 is looked at
only after the rule and parameters are fixed (adjusting the rule after
seeing the result is no verification).

Results of E1. In the control (compatible pairs m = 3, 6) all 860 orbits
locked completely with period m and the conserved quantity (norm) was
constant. In the test (incompatible pairs m = 5, 8, 10) the rotation was
mostly pulled in to 1/m, but the conserved quantity kept fluctuating
within the period on almost every periodic orbit, and some orbits moved
to other rationals near 1/m (11/54, 13/66 for m = 5). The ``pulls in but
never locks completely'' state did appear. U\^{}n = I came out as a
result. What did not appear was metastability in the strong sense (H4),
long dwelling with rare sudden changes; the intervals of constant
conserved quantity were short, 1--3 updates. m = 4, 12 were excluded
because of ties in rounding. An orbit of rotation number 2/9 appeared
for m = 5, but its unit is rotations per update, unrelated to Koide's δ.
The cause is that this model has no knob to vary ``how near''
continuously and cannot be brought to the edge of an Arnold tongue {[}A2
§3, §23,
\texttt{数値実験E1\_\allowbreak{}レプトン三相位相の格子引き込み\_\allowbreak{}20261014}{]}.

\pandocbounded{\includegraphics[keepaspectratio]{figs/te14_E1_fig.png}}

\textbf{Fig. 6.} E1: lattice pull-in of the three-phase lepton angle.
Compatible pairs (m = 3, 6) lock completely on every orbit; incompatible
pairs (m = 5, 8, 10) have their rotation pulled in to 1/m while the
conserved quantity fluctuates, and some orbits move to other rationals.
\texttt{te14\_\allowbreak{}E1\_\allowbreak{}lepton\_\allowbreak{}phase\_\allowbreak{}lattice.py},
\texttt{te14\_\allowbreak{}E1\_\allowbreak{}plot.py}.

\subsection{15. Stage 5: when α is
introduced}\label{stage-5-when-ux3b1-is-introduced}

\begin{quote}
Then conversely, if we do it with 23/124, this agrees from the outset if
we use the electron mass and α. The ground state, that is. (21:17)
\end{quote}

Claude here set up two arguments. The electron, being a non-decaying
ground state, should be on the completely locked side; for R\_\{124,23\}
U\^{}124 = I holds exactly, so this corresponds to the electron, while
the decaying muon and τ come on the ``near but unlocked'' side---an
assignment. And an argument for ``why there is no metastability without
23/124'': in the circle map a low-order rational (1/6) has a wide
locking interval so that all nearby states lock deeply, whereas a
high-order root such as 23/124 has an extremely narrow interval so that
states just outside it sit at the edge and behave intermittently.

\begin{quote}
No ``why'' is needed. Because it is so, the electron mass can be
predicted---that is all. Is a reason needed? Rather, the problem is why
it locks to that window. The answer is written in the note I gave you.
(21:23)
\end{quote}

The author dismissed it. The Fourteenth Thought Experiment states
explicitly ``the reason for the nearness is not explained; nearness is
the input of the search'', and the argument about the width of the
locking interval was unnecessary. The assignment ``electron = exact lock
U\^{}124 = I'' also conflicts with the paper. The answer written in the
paper is this. Why does it lock to that window? Because the
integer-selecting structures lie near each other there. The
four-dimensional cells on the space side, 137 = 1 + 8 + 128; the 23/124
on the evolution side; the 248 = 8 + 112 + 128 and 124 = 4 × 31 of
E₈---these share integers in this neighbourhood while differing in form.
So the state is pulled into this window but, agreeing with none of them,
never locks completely and keeps moving among the three 137s. The
observed α is the value read by integrating that history along the
direction of evolution {[}A1, eqs. (3)(4){]}. Numerically: the integer
137; the observed α⁻¹ = 137.035999177; x\_g = 137.036010988;
x\_\{23/124\} = 137.036042915. The observed value is smaller than both
non-integer values, so moving only between those two can never reach it,
but it lies within the range of the three 137s including the integer.

\subsection{16. Stage 9: the symmetry of the four-dimensional
lattice---why the background is
four-dimensional}\label{stage-9-the-symmetry-of-the-four-dimensional-latticewhy-the-background-is-four-dimensional}

In the author's eleven stages this is stage 9, but in the dialogue it
was brought in immediately (10-14 21:25) as the answer to ``why does it
lock to that window'' of Chapter 15, and E2 designed here fixes the
method of E3 onwards, so the paper places it here.

\subsubsection{16.1 The length of the diagonal is an integer
ratio}\label{the-length-of-the-diagonal-is-an-integer-ratio}

\begin{quote}
Have you forgotten the special symmetry of the four-dimensional lattice?
Have you forgotten that the length of the diagonal becomes an integer
ratio? (21:25)
\end{quote}

The diagonal of the four-dimensional unit cube is √4 = 2, an exact
integer multiple of the edge. This happens only in four dimensions. Half
the diagonal is 1, so the half-integer point (½,½,½,½) at the centre of
the cube has the same length 1 as the axis point (1,0,0,0); 8 integer
points and 16 half-integer points lie on the same sphere, 24 in all.
This is the 24-cell, the 24 units of the Hurwitz integers in quaternion
terms. In four dimensions the cell centre (integer point) and the cell
corner (half-integer point) are equivalent under the symmetry.
Recounting 137 by shells of the farthest corner (half-integer point) of
each cell gives 137 = 1 + 8 + 24 + 40 + 64 (corner norms 1, 3, 5, 7, 9),
and the 128 of 1 + 8 + 128 splits as 24 + 40 + 64. The 128 is a quantity
counted at half-integer points, on the same side, in the manner of
counting, as the 128 half-integer roots of E₈.

The discrepancy of E1 resolves in four dimensions. In two dimensions the
90° Z₄ (Gaussian integers) and the 60° Z₆ (Eisenstein integers) are
separate lattices and incompatible; the rounding ties that excluded m =
4 on the hexagonal lattice in E1 came from this incompatibility. Among
the 24 Hurwitz units in four dimensions, the order-4 elements i, j, k
and the order-6 element (1+i+j+k)/2 are present together. Z₄ and Z₆
coexist in one lattice in four dimensions. The evolution side (23/124)
is λₐ³¹ = i in the K4 paper {[}B1 §6.5{]}, and the ``4'' of 124 = 31 × 4
is the four phases of 90°, the Z₄ side; the three phases (leptons) are
the 120° and 60° phases, the Z₆ side.

\subsubsection{16.2 Why the background is four-dimensional. This is the
answer}\label{why-the-background-is-four-dimensional.-this-is-the-answer}

\begin{quote}
Yes. You seem to have noticed the reason for three dimensions, but you
have forgotten why the background is four-dimensional. This is the
answer. (21:28)
\end{quote}

The background is four-dimensional because four is the smallest
dimension in which all the needed integer symmetries fit simultaneously
into one lattice. The three-phase leptons of E0 took the (1,1,1)
direction of three dimensions as the invisible axis; but the diagonal of
the three-dimensional cube is √3, not an integer, so the length of the
invisible axis is irrational and the three phases and the invisible axis
cannot be closed on an integer lattice. The three-dimensional cubic
lattice has 90° (Z₄) and 120° (Z₃) rotations but no pure 60° (Z₆)
rotation. Adding one axis makes the diagonal (1,1,1,1) of length 2 with
half-length 1, so centre and corner become equivalent and the 24 Hurwitz
units arise, with the order-4 i (the Z₄ of the evolution side) and the
order-6 (1+i+j+k)/2 (the Z₆ of the three phases) coexisting, and a point
of the four-dimensional light cone arising automatically from a pair of
complex numbers {[}A1 eq. (17c){]}. Choosing one axis as the projection
axis in central projection drops the dimension by one {[}A1 §11.4{]}.
Since the background is four-dimensional, the visible space is
three-dimensional and the direction of the projection axis becomes
invisible. The three phases (three generations) and three-dimensional
space are the same ``3'': both are ``what remains of a four-dimensional
space with an integer diagonal when one axis is dropped by projection''.
The ``why three dimensions'' of Chapter 9 (the three conditions on S³)
and this ``why the background is four-dimensional'' (the integer
diagonal) are two faces of one answer.

\subsubsection{16.3 E2: the four-dimensional Hurwitz lattice and
23/124}\label{e2-the-four-dimensional-hurwitz-lattice-and-23124}

\begin{quote}
Go ahead. (21:29)
\end{quote}

The state is placed on the Hurwitz lattice (four dimensions); the
evolution side identifies λₐ³¹ = i with the quaternion i and updates
with U\_\{R₁₂₄,₂₃\}; the lepton three-phase is represented by the
Hurwitz units of order 3 and 6; the 23/124 rotation and the rounding
onto the lattice are imposed on the same state at once; no external
clock is placed. Results: the controls (the 90° Z₄ axis and the 60° Z₆
axis, compatible with the lattice) lock completely on every orbit with
constant conserved quantity and zero rounding ties. At radius 3 (the
137-cell ball) the 23/124 rotation is fully captured by the lattice; on
the Z₆ axis (the three-phase axis) all 936 points lock to 1/6, and on
the Z₄ axis they fall to orbits of period 4, 6, 16. As the radius grows,
the effective rotation number is captured step by step by the
convergents of the continued fraction {[}0; 5, 2, 1, 1, 4{]} of 23/124
and approaches 23/124, in the order 1/6 → 3/16 → 5/27 → 13/70, 18/97 →
23/124; from radius 16 orbits of period 124, 248, 372 appear. This
staircase was not input; it emerged from the combination of lattice and
rotation. The readout of α varies from orbit to orbit and has not
converged to a single value {[}A2 §5,
\texttt{数値実験E2\_\allowbreak{}四次元フルヴィッツ格子と23\_\allowbreak{}124\_\allowbreak{}20261014}. Note: E2
measured the rotation number as the mean of the per-step angle, which is
not additive in four dimensions; the α values of E2 come from that
measurement and are not used in this paper. From E3 on it is computed
from winding numbers{]}.

\pandocbounded{\includegraphics[keepaspectratio]{figs/te14_E2_fig.png}}

\textbf{Fig. 7.} E2: the 23/124 rotation and rounding onto the
four-dimensional Hurwitz lattice imposed together. As the radius grows
the orbit is captured along the staircase of convergents 1/6 → 3/16 →
5/27 → 13/70, 18/97 → 23/124. The α readout in panel (b) uses the
withdrawn measurement (angle averaging) and is not used in this paper.
\texttt{te14\_\allowbreak{}E2\_\allowbreak{}hurwitz\_\allowbreak{}124\_\allowbreak{}23.py},
\texttt{te14\_\allowbreak{}E2\_\allowbreak{}radius\_\allowbreak{}scan.py}, \texttt{te14\_\allowbreak{}E2\_\allowbreak{}plot.py}.

\subsection{17. Stage 6: a detour when mass cannot be
read}\label{stage-6-a-detour-when-mass-cannot-be-read}

\subsubsection{17.1 If α is known the electron mass is fixed (the
standard
theory)}\label{if-ux3b1-is-known-the-electron-mass-is-fixed-the-standard-theory}

\begin{quote}
Have you forgotten that the standard theory explains that the electron
mass is fixed once α is known? (21:35)
\end{quote}

The electron mass follows directly from α: mₑ = 2hR∞/(α²c), where the
Rydberg constant R∞ is measured separately by hydrogen spectroscopy to
about 10⁻¹², and CODATA indeed fixes the electron mass by this formula.
Here Claude fixed h and c at their SI defined values and computed: the
integer 137 → 0.5107305 MeV, x\_g → 0.5109990388 MeV, x\_\{23/124\} →
0.5109992769 MeV, observed 0.5109989500 MeV. Since the mass is
proportional to the square of α⁻¹, δmₑ/mₑ = 2δ(α⁻¹)/α⁻¹.

\subsubsection{17.2 The ground state: the electron has no internal
structure, but its energy levels are not
one}\label{the-ground-state-the-electron-has-no-internal-structure-but-its-energy-levels-are-not-one}

\begin{quote}
So have you forgotten that it is the ground state? The electron has no
internal structure, but is its energy level only one? (21:37) Go ahead.
(21:38)
\end{quote}

The electron itself has no internal structure, but bound to a nucleus it
has the ladder of energy levels numbered by the integer n, Eₙ =
−α²mₑc²/(2n²), and the energy of its ground state n = 1 is the Rydberg
constant itself, hcR∞ = ½α²mₑc² = 13.6057 eV. ``Once α is known the
electron mass is fixed'' means ``fixed through the ground-state
energy'', and to read the electron mass from α one must read it from the
ground state. The ground state has nowhere lower to fall and does not
decay: a completely locked state. The excited states are metastable with
finite lifetimes (2p about 1.6 ns, 2s about 0.12 s) and fall by emitting
one photon, a single transition without fraction. This connects to the
starting question of the Fourteenth Thought Experiment, ``why does the
electron not fall into the nucleus''. The answer to reading the
orbit-to-orbit scatter of α in E2 as one value is not to average
everything but to read from the orbit corresponding to the ground state.

\subsubsection{17.3 E3, and the verdict in SI units (the
detour)}\label{e3-and-the-verdict-in-si-units-the-detour}

In E3 the rotation number was computed instead from the integer winding
numbers of the two complex components of the isoclinic rotation; then
every orbit of period 124 is exactly 23/124. Results: orbits with
rotation number 23/124 and constant conserved quantity (norm): 0 of 480.
Orbits whose rotation number is exactly 23/124 are many, reaching 85\%
on the Z₄ axis at radius 96, but all of them have a norm that fluctuates
within the period---the rotation is fully engaged but the conserved
quantity is not. In the level series on the Z₄ axis 23/124 is most
occupied, then 3/16, 29/156, 5/27, 13/70; on the three-phase axis (Z₆)
5/27 is most occupied (141 of 240). The transients before entering a
periodic orbit were short, a few hundred updates, and no stepwise
descent through the levels was seen.

Here Claude took α⁻¹ = 137.036042915 from the state locked at 23/124,
derived an electron mass of 0.5109992769 MeV, and judged it ``larger
than the observed value by 6.4×10⁻⁷ relative, about 2000 times the
measurement precision; locking on the 23/124 side alone does not give
the electron mass''. This is the detour: with Planck's constant left as
an input, it concluded that mass could not be read {[}A2 §7, §23,
\texttt{数値実験E3\_\allowbreak{}基底状態と準位系列\_\allowbreak{}20261014}{]}.

\pandocbounded{\includegraphics[keepaspectratio]{figs/te14_E3_fig.png}}

\textbf{Fig. 8.} E3: ground state and level series. With the rotation
number taken from winding numbers, orbits locked exactly to 23/124 reach
85\% on the Z₄ axis at radius 96, but no orbit has a constant conserved
quantity (norm). The occupied levels are 23/124, 3/16, 29/156, 5/27,
13/70. \texttt{te14\_\allowbreak{}E3\_\allowbreak{}ground\_\allowbreak{}and\_\allowbreak{}levels.py},
\texttt{te14\_\allowbreak{}E3\_\allowbreak{}plot.py}.

\subsection{18. Stage 7: reading α plainly, and deriving the electron
mass}\label{stage-7-reading-ux3b1-plainly-and-deriving-the-electron-mass}

\subsubsection{18.1 The natural units G = c = q₀ = 1: Planck's constant
is an
output}\label{the-natural-units-g-c-qux2080-1-plancks-constant-is-an-output}

\begin{quote}
It is the other way round. You have forgotten to compute in the natural
units G = c = q0 = 1, q0 = e/3. (21:45)
\end{quote}

In computing the electron mass as 2hR∞/(α²c), h had been fixed at its SI
defined value. That is placing Planck's constant as an input, the very
procedure the Fourteenth Thought Experiment rejected at the outset.
Recomputing in G = c = q₀ = 1 (q₀ = e/3, Gaussian): α = e²/(ħc) =
9q₀²/(ħc), so

\[
\hbar=9\,\alpha^{-1},
\]

and Planck's constant is not an input but an output fixed by α. The
direction is reversed: α comes first and ħ comes afterwards. With the
23/124 readout ħ = 9 × 137.036043 = 1233.32439; with the observed value
1233.32399. The unit of mass is m₀ = q₀/√G = 6.1974×10⁻¹⁰ kg and the
electron mass is mₑ = 1.46988×10⁻²¹ m₀. The relative uncertainty of the
electron mass in this unit system is fixed by the measurement precision
of G, about 1.1×10⁻⁵. The difference in α between the 23/124 readout and
the observed value is 3.2×10⁻⁷ relative, far smaller than the
uncertainty of the electron mass, so in this unit system it cannot be
distinguished and is not rejected. The conclusion written in E3, that
``the ground state locked at 23/124 does not give the electron mass'',
was an appearance created by a unit system that fixed ħ, and was
withdrawn {[}A2 §6, §23{]}.

\subsubsection{18.2 If α is 137, the electron mass necessarily
agrees}\label{if-ux3b1-is-137-the-electron-mass-necessarily-agrees}

\begin{quote}
Or rather, if α is 137 the electron mass necessarily agrees, does it
not? (21:48)
\end{quote}

The electron mass is not an independent test from α. mₑ = 4πħR∞/(α²c)
with R∞ measured by hydrogen spectroscopy, and in G = c = q₀ = 1 ħ =
9α⁻¹ is also fixed by α. Fixing one α, combined with the measured R∞,
fixes the electron mass automatically and no contradiction arises.
Checking separately whether the electron mass agrees adds no
information. What must be verified is α alone. α is dimensionless and
its value is the same in every unit system. The difference between the
observed 137.035999177(21) and the 23/124 readout 137.036042915 is
4.4×10⁻⁵, and as α itself they do not agree. This does not contradict
the position of the Fourteenth Thought Experiment, which presupposes
that none of the three 137s agrees exactly with the observed value and
reads the observed value as the state that never locks completely among
them, read along the direction of evolution.

\subsection{19. Stage 8: why α, why 23/124, is it not a circular
derivation?}\label{stage-8-why-ux3b1-why-23124-is-it-not-a-circular-derivation}

\subsubsection{19.1 α is obtained from 23/124, and the four-dimensional
lattice condition makes it easy to lock
there}\label{ux3b1-is-obtained-from-23124-and-the-four-dimensional-lattice-condition-makes-it-easy-to-lock-there}

\begin{quote}
Rather, the point is that α is obtained from 23/124, and that the
four-dimensional lattice condition makes it easy to lock to that α.
(21:50)
\end{quote}

There are two points. α is obtained from 23/124: from the finite-order
root R\_\{124,23\} of the K4 paper {[}B1{]} through N(R) = 4π/(1−R)².
And the four-dimensional lattice condition makes the state easy to lock
to that 23/124. The E3 results show the second point: on the Z₄ axis
(the four-phase side of λₐ³¹ = i) the fraction of orbits locking exactly
to 23/124 grows with radius to 85\% at radius 96, gathering through the
convergents 3/16 → 5/27 → 13/70 → 18/97. The conserved quantity still
fluctuates, so only the rotation number locks and the rest
fluctuates---``pulls in but never locks completely''. To say ``the
four-dimensional lattice selects 23/124'' a comparison is missing: as
the radius grows the effect of rounding shrinks, and any rotation number
may become easier to lock to its own value. Hence E4 compares the
locking rate with neighbouring roots and measures the pull-in width of
23/124. The answer to the suspicion of circularity---that the agreement
is an input if 23/124 was chosen using α---is to show, without looking
at α, that 23/124 is easily selected by the lattice condition.

\subsubsection{19.2 Why the observed value deviates slightly:
observation is an average over many
samples}\label{why-the-observed-value-deviates-slightly-observation-is-an-average-over-many-samples}

\begin{quote}
Further, why does it deviate slightly from the exact observed value? You
have not noticed that you explained it yourself: because observation is
an average over many samples. Because values in transition are observed
too. (21:52)
\end{quote}

The observed α is not a single value but an average over many samples,
which include not only the state locked at 23/124 but states in
transition towards the lock. So the observed value is not 137.036043
itself but slightly off. Checking the direction with the E3 data: of the
states in transition (Z₄ axis, 122 unlocked orbits) 86 lie below
137.036043 and 36 above, with mean 135.39 (3/16 gives 131.90, 13/70
136.44, 18/97 136.82, 29/156 135.96, 5/27 137.82). The states in
transition lie mainly below, so mixing them lowers the average, and the
observed 137.035999 lies below---the direction agrees. With this mean, a
fraction of about 2.65×10⁻⁵ of the observation time in transition would
suffice, but that breakdown was taken from an ensemble of initial
conditions, not weighted by time. The running of α (137.036 at low
energy, 128.9 at the Z-boson scale) admits a reading of the same
direction: measuring at higher energy samples more states in transition
and pulls α⁻¹ down. This is a verification task.

\subsubsection{19.3 E4: selection, pull-in,
recovery}\label{e4-selection-pull-in-recovery}

\begin{quote}
Try it. (21:53)
\end{quote}

\begin{enumerate}
\def\labelenumi{(\Alph{enumi})}
\item
  Does the four-dimensional lattice make 23/124 easy to lock? Replacing
  the rotation by 23/124 and by nearby roots, the fraction locking to
  its own ratio was compared at radii 32, 64, 96. 23/124: 40\% → 70\% →
  86\%. Neighbouring roots of similar denominator (23/122, 23/126,
  22/119, 24/131): 0--27\%. 17/92: 53\% → 93\% → 100\%. A further check
  showed that roots whose denominator is a multiple of 4 lock easily
  (40--76\% at radius 64 against 10--40\% for the others): with a
  denominator divisible by 4 the rotation is exactly i after n/4 steps,
  and i is a symmetry of the lattice, so the rounding error can be
  closed. 124 = 4 × 31 belongs to this group, and the ``4'' in λₐ³¹ = i
  of the K4 paper is precisely the reason it locks easily on the
  four-dimensional lattice. Since 17/92 and others of the same group
  lock just as easily, the four-dimensional lattice alone does not
  select 23/124 uniquely.
\item
  Pull-in to 23/124. With the rotation displaced slightly from 23/124,
  83--93\% were pulled in to 23/124 for displacements within 1.4×10⁻⁵
  and 24--31\% even at 1.4×10⁻⁴. The comparison 23/122 reached at most
  43\% and was not pulled in at 1.4×10⁻⁴. The pull-in width of 23/124 is
  wide and strong.
\item
  Leaving the lock and returning. Disturbing a state locked at 23/124 by
  one lattice unit, 18 of 20 returned to 23/124, after 274 updates on
  average. The time-weighted mean of the readout on the way back was α⁻¹
  = 136.74, below the 23/124 value, agreeing in direction with the
  downward offset of the observed value. With this mean, a fraction of
  about 1.46×10⁻⁴ of the observation time in transition would give the
  observed value, and the lock would have to break about once in 1.9
  million updates. The mechanism producing that frequency is not in this
  deterministic system {[}A2 §5, §8,
  \texttt{数値実験E4\_\allowbreak{}選択・引き込み・回復\_\allowbreak{}20261014}{]}.
\end{enumerate}

\pandocbounded{\includegraphics[keepaspectratio]{figs/te14_E4_fig.png}}

\textbf{Fig. 9.} E4: selection, pull-in, recovery. (A) Radius dependence
of the locking rate of 23/124 and neighbouring roots (roots with
denominator divisible by 4 lock easily); (B) pull-in width to 23/124;
(C) recovery after a one-lattice-unit kick, with the readout during the
return lying below. \texttt{te14\_\allowbreak{}E4\_\allowbreak{}selection\_\allowbreak{}pullin\_\allowbreak{}recovery.py},
\texttt{te14\_\allowbreak{}E4\_\allowbreak{}mod4\_\allowbreak{}check.py}, \texttt{te14\_\allowbreak{}E4\_\allowbreak{}plot.py}.

\subsubsection{19.4 The electron itself distorts the gauge:
E5}\label{the-electron-itself-distorts-the-gauge-e5}

\begin{quote}
You seem to have forgotten that the mass-energy of the electron itself
slightly distorts the four-dimensional gauge, so that it can fall out of
the gauge. (22:09) Save only the programs and the documents. Then check
the gauge deviation. (22:16)
\end{quote}

The lattice used in E4 was fixed and the electron merely moved on it, so
once locked at 23/124 it never left, and the frequency of leaving had to
be assumed from outside. In reality the electron's own mass-energy
slightly distorts the four-dimensional lattice (gauge); on the distorted
lattice the locking condition itself shifts slightly, so the state
leaves the lock from time to time and returns. The cause of leaving need
not be brought in from outside. In the language of the standard theory
it is the same kind of effect as the self-energy of the electron
interacting with its own field, which produces the anomalous magnetic
moment (g − 2 ≈ α/π).

In E5 the lattice used for rounding was stretched at each update in
proportion to the electron's Coulomb-like self-energy α/r (r the radius
in lattice units), in both the stretching (s = +1) and the shrinking (s
= −1) direction, decided before looking at results; α was updated from
the readout and iterated three times, with no externally tuned
parameter. Results: the distortion does change the orbits (15 of the
same 24 initial points fell to different periodic orbits). But it does
not act to break the lock: the orbits locked at 23/124 were 14 without
distortion, 19 with stretching and 17 with shrinking---more, not fewer.
No leaving-and-returning appeared: within the locked orbits the
intervals where the local rotation number departs from 23/124 were zero;
every orbit is either ``always locked'' or ``always on another level''.
The α⁻¹ averaged over all orbits was 137.0059 for stretching and
137.0940 for shrinking, but these are ensemble averages, not time
averages, and cannot be compared with the observed value. Updating α
changed it only by 2×10⁻⁴ relative and it was a fixed point after the
first iteration. With a distortion fixed instantaneously by the
electron's position, the map on a finite lattice is deterministic and
must enter a periodic orbit, so leaving the gauge and returning cannot
in principle occur. Claude wrote here that ``the gauge side must have a
state of its own and follow with a delay'', but in Chapter 21 it turns
out that the partner that leaves and returns is the partner wave of the
exchange itself {[}A2 §9,
\texttt{数値実験E5\_\allowbreak{}ゲージの自己歪み\_\allowbreak{}20261014}; the plotting program
\texttt{te14\_\allowbreak{}E5\_\allowbreak{}plot.py} for E5 was added by the author{]}.

\pandocbounded{\includegraphics[keepaspectratio]{figs/te14_E5_fig.png}}

\textbf{Fig. 10.} E5: self-distortion of the gauge. (a) Number of orbits
locked at 23/124 (none 14, stretched 19, shrunk 17); (b) difference
between each orbit's rotation number and 23/124; (c) number of orbits
whose periodic orbit changed under distortion (15 of 24). The
self-distortion does not break the lock but increases it, and intervals
of leaving and returning are zero.
\texttt{te14\_\allowbreak{}E5\_\allowbreak{}gauge\_\allowbreak{}backreaction.py},
\texttt{te14\_\allowbreak{}E5\_\allowbreak{}control.py}, \texttt{te14\_\allowbreak{}E5\_\allowbreak{}plot.py}.

\subsubsection{19.5 Single electrons have been observed. Time average
and ensemble
average}\label{single-electrons-have-been-observed.-time-average-and-ensemble-average}

\begin{quote}
Fine for now. But is the behaviour of a single electron not, in fact,
little observed? (22:20)
\end{quote}

Single electrons are well observed. In Penning-trap experiments that
confine one electron for months, the g-factor of a single electron was
measured in 2023 to 10⁻¹³, giving α⁻¹ = 137.035999166(15) {[}X5-Fan{]}.
This is not an average over many electrons but the time average of one
electron measured for a long time. In 1999 the quantum jumps of a single
trapped electron between the lowest cyclotron level and the next were
observed directly: the electron dwells in one level for seconds and
suddenly jumps one step {[}X5-Peil{]}. The ``stays almost engaged for a
long time and occasionally slips one step'' of the Fourteenth Thought
Experiment is seen directly in one electron. Precise values of α come in
two families: from the g-factor of a single electron (time average), and
from atomic recoil (ensembles of many atoms: rubidium 137.035999206(11)
{[}X5-Morel{]}, caesium 137.035999046(27) {[}X5-Parker{]}), which
disagree by more than 5σ for reasons unknown. In this system's view,
measuring one electron for a long time and measuring many particles at
once pick up states in transition in different proportions and need not
give the same α.

\subsubsection{19.6 Expectation values do not behave like classical
probability}\label{expectation-values-do-not-behave-like-classical-probability}

\begin{quote}
As expected. In the standard theory, energy shifts merely because
particles are crowded together, so it shifts easily. Likewise a
deviation appears even in the time average. Do you see? The behaviour of
expectation values differs from classical probability. (22:23) The
experiment can wait; write down the insights. There are many other
things to do. (22:25)
\end{quote}

The average taken in E4 and E5, ⟨α⁻¹⟩ = (1−f)·x\_\{23/124\} +
f·x\_transition, weighting the locked and transitional values by their
fractions, is the average of a classical mixture, being in one state or
the other with some probability, and cannot be used as it stands. For a
superposition ψ = a·ψ\_lock + b·ψ\_transition the expectation value is
⟨A⟩ = \textbar a\textbar²A₁₁ + \textbar b\textbar²A₂₂ +
2Re(a*b·A₁₂·e\^{}\{iΔωt\}), and the last interference term is
proportional to the amplitude b, not the probability
\textbar b\textbar². Even a slight admixture of the transitional state
shifts the value by its square root (for \textbar b\textbar² = 10⁻¹⁰ the
interference term can be 10⁻⁵). This is as §5 of the Fourteenth Thought
Experiment says: ``the readout after integrating complex amplitudes with
their phases differs in general from the average of the constant
itself''. When particles are crowded, the interaction energy enters as
an overlap term and the levels themselves move (collision and density
shifts in atomic clocks). The interference term oscillates as
e\^{}\{iΔωt\}, but when the frequency ratio is near but not locked
(irrational), the beats do not cancel completely in a finite time and a
deviation remains; even one electron measured for a long time gives
slightly different values depending on how the measurement time is
taken. The transitional fraction 1.46×10⁻⁴ estimated in E4 assumed a
classical mixture; with amplitudes superposed, the admixture needed is
orders of magnitude smaller. The readout should not average α over
orbits but integrate the exchange amplitude of the K4 paper (G\_R = 1 −
R with a phase attached) along the history as a complex number, and read
α afterwards {[}A2 §10, §11{]}.

\subsection{20. The harmonic content of the electron, the absence of an
absolute scale, and the minimal structure of the
wave}\label{the-harmonic-content-of-the-electron-the-absence-of-an-absolute-scale-and-the-minimal-structure-of-the-wave}

\subsubsection{20.1 The arrangement of harmonics can be predicted: the
quarter-period
symmetry}\label{the-arrangement-of-harmonics-can-be-predicted-the-quarter-period-symmetry}

\begin{quote}
The problem is: can the harmonic content of the electron be predicted?
(22:27)
\end{quote}

Part of it can. Of the 33 orbits of period 124 locked at 23/124 (radius
64), in 29 the state multiplied by i coincided with the state exactly 93
updates later (3 × 31). Since i is a symmetry of the four-dimensional
lattice (a 90° rotation), these orbits coincide with themselves when
shifted by a quarter period. With this symmetry, when the orbit is
decomposed in frequency the components that appear are limited to (23 +
4k)/124 (k an integer): the fundamental 23/124 and sidebands at spacing
1/31 on either side. In all 29 orbits 100\% of the energy fell in this
series. In the remaining 4 without the symmetry the state multiplied by
i belonged to another orbit, and 0.01--0.07\% fell outside the series.
This is a rule of the same form as for a square wave: a square wave has
the symmetry of reversing sign under a half-period shift, so only odd
harmonics appear; on the four-dimensional lattice it becomes a
quarter-period symmetry and only every fourth component appears. The
``4'' and the ``31'' of the period 124 = 4 × 31 are directly the
selection rule and the spacing of the harmonics. At radius 64 the
fundamental holds 99.93\%, the adjacent sidebands 19/124 and 27/124
1.8×10⁻⁴ and 3.2×10⁻⁴, decreasing with order. These strengths, however,
change with the ratio of the lattice cell to the orbit radius {[}A2
§12.1, §23, \texttt{te14\_\allowbreak{}harmonic\_\allowbreak{}selection\_\allowbreak{}check.py}{]}.

\subsubsection{20.2 This system has no absolute scale; everything is a
ratio}\label{this-system-has-no-absolute-scale-everything-is-a-ratio}

\begin{quote}
Before that: this system has no absolute scale; everything is a ratio.
Is the electron always localised? If unobserved, is its position not
undetermined, like the photon? Depending on that, the harmonic content
differs by 10\^{}60 orders. (22:33)
\end{quote}

Claude had written that ``the radius at which the electron lives must be
decided'', which presupposed an absolute scale. This system has only
ratios. An unobserved electron, like the photon, has no determined
position. In a spread state few harmonics are needed, close to the
fundamental alone; observed and localised, many harmonics are needed.
The number of harmonics needed is roughly the ratio ``total extent /
width of localisation'': the ratio of the size of the universe to the
electron's Compton wavelength is about 10³⁸, and to the Planck length
about 10⁶¹. The same electron can have a number of harmonics differing
by 10⁶⁰ orders according to its degree of localisation. ``The harmonic
content of the electron'' cannot be fixed as one thing; it is a readout
quantity depending on how localised it is (how it was observed). What
can still be predicted as a ratio is the selection rule (that the
components are limited to the series (23 + 4k)/124), independent of
scale and of the degree of localisation. What was called the initial
radius in E2--E4 was not an absolute size but the ratio of orbit size to
lattice cell, the resolution of localisation; that the sidebands
weakened at larger radius expresses that a coarser (more spread)
localisation needs fewer harmonics. The question becomes not ``what is
the harmonic content of the electron'' but ``how are the harmonic
strengths fixed when the localisation ratio is given, and how does
observation (interaction) choose that ratio''.

\subsubsection{20.3 Yet it is not a single tone: the minimal structure
is two spin components × three three-phase
components}\label{yet-it-is-not-a-single-tone-the-minimal-structure-is-two-spin-components-three-three-phase-components}

\begin{quote}
Well done. But as long as interactions and masses differ, there must be
a minimal structure of the wave; it cannot be a single tone, can it? The
degrees of freedom are insufficient. (22:35)
\end{quote}

It cannot be a single tone, for two reasons. First, the four-dimensional
lattice does not permit a single tone. A single tone is a pure rotation
whose phase turns at a constant rate; a 23/124 single tone would be a
rotation of order 124, but the rotation orders the four-dimensional
lattice allows are 1, 2, 3, 4, 5, 6, 8, 10, 12 only (the
crystallographic restriction {[}A1 §10.1{]}). A pure rotation of order
124 cannot exist on the lattice, so a state locked at 23/124 on the
lattice necessarily contains components besides the fundamental. In E3,
even among orbits locked exactly to 23/124, not one had a constant
conserved quantity (a pure rotation). The sidebands (23 ± 4k)/124 weaken
as the localisation ratio changes but do not vanish; in this system the
ratio is always finite, so the sidebands never disappear. This is the
smallest part of ``the minimal structure of the wave''. Second, a single
tone has too few degrees of freedom. In a system of ratios alone, the
amplitude and phase of a single tone are relative to the observer and
the only thing it can own is one frequency ratio; one number cannot
express the differences of charge, spin, mass, generation and
interaction. The minimum needed is: spin (a state on the Hurwitz lattice
is a quaternion with two complex components z₁, z₂, and left
multiplication on the Z₄ axis rotates both at once, so it behaves as a
two-component spinor), charge (the electron's charge is 3 × q₀ with q₀ =
e/3; only with the three three-phase components together does an integer
charge arise), generation and mass (the position of the three-phase
angle, Koide's δ, distinguishes the generation, and the rotation ratio
fixes the mass), and the sidebands forced by the lattice. The candidate
minimal structure of the electron is ``two spin components × three
three-phase components, each with its fundamental and the sidebands the
lattice forces''. When observation localises it, the many harmonics
needed for localisation are added on top of this minimal structure; the
minimal structure remains even when spread, and what changes by up to
10⁶⁰ orders is the added part. These ``two spin components'' are the
spin read out in Chapter 11 as left quaternion multiplication on S³,
reappearing on the four-dimensional lattice {[}A2 §12.3{]}.

\subsection{21. Localisation and return, the self-consistent lock,
observation itself as dynamics, checking the
answers}\label{localisation-and-return-the-self-consistent-lock-observation-itself-as-dynamics-checking-the-answers}

\subsubsection{21.1 The mechanism by which observation localises the
electron, and by which it
returns}\label{the-mechanism-by-which-observation-localises-the-electron-and-by-which-it-returns}

\begin{quote}
Having examined this far, I had forgotten one thing. Do you remember the
mechanism by which observation localises the electron? And there is also
a mechanism by which it returns even after localisation. Read the note
on the Thirteenth Thought Experiment and the original paper. (10-15
03:12)
\end{quote}

The mechanism of the Thirteenth Thought Experiment {[}D1{]} is this.
Localisation occurs not because an observer shrinks a wave packet but
because harmonics gather through interaction. When a spread, nearly
single-tone wave A (effective order 1) scatters by exchange with a
sharply localised wave B (effective order 32), B's harmonics move to A
and A localises (at R = 0.70, A's effective order reaches 31.1 at the
third collision). An unobserved electron is spread and appears localised
when it exchanges with the localised wave of the measuring apparatus and
receives harmonics. The mechanism of return is that a closed exchange is
reversible. The transferred harmonics are not lost but shuttle back and
forth; there is no dissipation, the total power of each wave stays 1 and
the sum of effective orders is always N\_A + N\_B = 33. At exchange
phase kω = 90° the two waves share equally, at 180° they are completely
exchanged, at 360° they return. For R = R\_\{124,23\} the shapes are
exactly exchanged at the 62nd collision and the complex state itself
returns to the initial state at the 124th (ε\_state = 5.1×10⁻¹⁵) {[}D1
§6.2{]}. A localised electron returns to the spread state when the
exchange completes one cycle.

The connection with the present examination is fourfold. It is the same
structure of 124: the phase difference of the exchange eigenvalues is
ω/2π = 101/124 with 101 = 124 − 23; the 23/124 rotation on the
four-dimensional lattice and the phase of one exchange are the same
number with opposite direction, and the quarter marks λ³¹ = i, λ⁶² = −1,
λ⁹³ = −i coincide. That the state multiplied by i appeared 93 updates
later in the lattice orbits is this structure. The harmonic content is
not fixed by the electron alone; the effective order shuttles between
the electron and its partner and what is conserved is the sum. The
``difference of 10⁶⁰ orders'' of §20.2 is the difference of the phase at
which this shuttling is read. The mechanism of ``leaving and returning''
was here: in E5 Claude wrote that ``the gauge side needs delay or
memory'', but that is the partner wave of the exchange itself, which
holds the harmonics and gives them back. When the electron's own energy
distorts the gauge slightly and R departs a little from R\_\{124,23\},
the complete return at the 124th collision becomes a ``nearly returns''
like the 103rd at R = 0.70, and a small remainder (1.6 parts in 10⁴ at R
= 0.70) is left. This remainder that fails to return is a candidate
cause of the slight deviation of the observed α from the 23/124 value
{[}A2 §13{]}.

\subsubsection{21.2 The mechanism that pulls α in to 137 is the
mechanism that makes the reflectance R
0.7}\label{the-mechanism-that-pulls-ux3b1-in-to-137-is-the-mechanism-that-makes-the-reflectance-r-0.7}

\begin{quote}
This connects with the later paper and dynamics has nothing to do with
it. In the examination so far, the mechanism that pulls α in to 137 is
itself the answer, and it is the mechanism that makes the reflectance R
0.7. At the time of that note I simply had not noticed. (03:17)
\end{quote}

The Thirteenth Thought Experiment had reserved: ``R is a value put in by
hand and the return at the 124th collision is merely its consequence; to
fix R from the interaction a dynamical model is needed.'' The order was
reversed. R and α are tied one-to-one by N(R) = 4π/(1−R)². The mechanism
that pulls α in to 137 and the mechanism that pulls the reflectance R in
to about 0.7 are one and the same phenomenon. R\_\{124,23\} =
cos²(23π/124) = 0.69718, which in α⁻¹ is 137.036043. R was not a value
put in from outside but the value selected as the result of locking to
23/124. That 0.7 looked cleanest in the coarse sweep of July (0, 0.51,
0.55, 0.6, 0.7, 0.9, 1) and that the centre converged to R\_\{124,23\}
on fine sweeping can be read as this pull-in seen from outside (the
coarse sweep was on the two-wave exchange system and the fine sweep on
the metastable-interface system, different systems sharing the same U\_R
{[}D1 §6.2{]}). On the four-dimensional lattice, roots with denominator
divisible by 4 lock easily and 124 = 4 × 31 belongs to that group (E4);
the pull-in width of 23/124 is wider and stronger than that of its
neighbours (E4); and the phase of one exchange, 101/124, and the lattice
rotation 23/124 are the same number with coinciding marks. The ``state
returns at the 124th collision'' of the Thirteenth Thought Experiment is
then not the consequence of a period put in by hand but one continuous
story: the four-dimensional lattice selects 23/124, as a result R
becomes 0.697, and at that R the exchange closes in 124 collisions.
17/92, of the same ``denominator divisible by 4'' group, locks at least
as easily in E4 and R\_\{92,17\} = 0.69920 is also right next to 0.7,
but in α⁻¹ it is 138.885, not 137. That R is pulled in to about 0.7 can
be explained by the four-dimensional lattice alone; for 23/124 to be
selected within that, the other constraint of the Fourteenth Thought
Experiment---the 137 cells on the space side, or the 124 = 4 × 31 of
E₈---is needed {[}A2 §14{]}.

\subsubsection{21.3 The only self-consistent lock is this one.
Observation itself is
dynamics}\label{the-only-self-consistent-lock-is-this-one.-observation-itself-is-dynamics}

\begin{quote}
That is, because this lock is the only one that becomes self-consistent.
Simple. Further, the reason I thought dynamics was needed at that time
was that the mechanisms of translation and of recoil and approach by
emission and absorption had not been integrated; but this mechanism,
too, is inherent in the system. Observation itself is dynamics. To be
exact, it suffices to add a complete absorbing shell at infinity to the
present system---the note on the Eleventh Thought Experiment. (03:23)
\end{quote}

Once α is fixed, in G = c = q₀ = 1 ħ = 9/α is fixed and so is R, by N(R)
= 4π/(1−R)². That R must be a finite-order root the four-dimensional
lattice can lock to. And the α read from the locked state must equal the
starting α. This loop closes for 23/124. 17/92, although it locks easily
on the lattice, gives α⁻¹ = 138.885 and the loop does not close with the
137 of the space side; it falls out by the condition of self-consistency
alone. At present, however, the only evidence that 23/124 is the only
closing root is that 17/92 falls out; no computation has yet enumerated
the roots with denominators divisible by 4 and counted those whose N(R)
falls near 137. That is a verification task.

Observation itself is dynamics. In the framework of the Eleventh Thought
Experiment {[}D2{]} the state is only the occupation of levels and the
centre-of-mass velocity relative to the absorber, with neither
coordinates nor metric. Motion is read only from the record (time,
direction, frequency, polarisation) arriving at the complete absorbing
shell at infinity, and no readable change exists between events. That
from 29 records of 21 cm the velocity was recovered to 2.4×10⁻¹⁰ m/s,
recoil history included, and from the records of 200 atoms the
centre-of-mass velocity of 370 km/s to 0.09 m/s, shows that translation
and recoil are already in the record. An isotropic absorbing shell has
no direction and neither translation nor recoil can be defined {[}D2
§6.1{]}; direction arises only as an assignment to the anonymous
elements of the shell. The shell itself is not a background spacetime
but the map onto the section (the null surface) where the future light
cone of each emission event cuts infinity, the same operation as central
and radial projection {[}D2 §2.5{]}. The number of elements of the shell
is N = 10⁶⁰, fixing the resolution of direction (a value placed in the
Eleventh Thought Experiment, not derived {[}D2, non-claim 11{]}). The
localisation of the electron is fixed by which of these anonymous
elements it exchanged with. The discrete lock needed, in §6.5 of the
Eleventh, for ``the ground state does not radiate'' together with the
time evolution of the field, is the self-consistent lock to 23/124. In
the same §6.5, recorded exchanges (emission and absorption) give the
frequency itself and unrecorded exchanges give its shift (the Lamb
shift). The difference between the 23/124 value 137.036043 and the
observed 137.035999 can be read as a ``shift by unrecorded exchange'' of
the same kind.

Adding a complete absorbing shell at infinity to the present system (the
four-dimensional Hurwitz lattice, the exchange of the Thirteenth Thought
Experiment, the lock of the Fourteenth) brings together in one system
localisation (observation), the mechanism of return (closed exchange),
translation and recoil (the record), and the slight deviation of α
(unrecorded exchange). The five tasks at the end of Chapter 13 (emission
and absorption, the distinction of gravitational and electromagnetic
waves, recoil, attraction, the contributions of spin and charge) are
inherent in the system as this record on the absorbing shell. Both
reservations of the Thirteenth Thought Experiment (that R is an input,
that a dynamics is needed) are lifted. This is a connected part, a
hypothesis {[}A2 §15{]}.

\subsubsection{21.4 The answers can be checked against the
theory}\label{the-answers-can-be-checked-against-the-theory}

\begin{quote}
Please. Moreover, that note also prepares the numbers for checking the
answers for the hydrogen atom. So the theory can be checked against the
answers. (03:29)
\end{quote}

The Eleventh Thought Experiment has already produced the answer table of
the standard theory in the same unit system (G = c = q₀ = 1) and with
the same absorbing-shell reference. The integrated system must reproduce
it. For the hydrogen atom: the total emission from 5p₃/₂ F=1,
13.054539268 eV (equal to the level drop); the 1S ionisation energy
13.5984684 eV (its difference from measurement coincides with the
omitted QED and nuclear size, 8170.44 MHz, with a discrepancy of 0.12
MHz); E1 A-coefficients within 1.2×10⁻⁴ of NIST; the Lyman-α width 99.71
MHz (measured 99.7); the 21 cm A-coefficient 2.868×10⁻¹⁵ s⁻¹; the spin
temperature 2.72550 K at 2.7255 K; the recovery of the centre-of-mass
velocity from the record, 0.09 m/s (200 atoms) and 2.4×10⁻¹⁰ m/s (29
events of one atom). For the hydrogen molecule: 301 levels, E2
A-coefficients agreeing with the 4669 lines of Roueff et al.~(2019) to a
median of 8×10⁻⁴. For electron--electron and proton--proton: no bound
levels, only the Mott cross-section and quadrupole bremsstrahlung. The
check can be done in two steps. Everything in which the Eleventh Thought
Experiment deviated from measurement was the omitted time evolution of
the field, i.e.~unrecorded exchange: the 8170 MHz of 1S and the 1057.85
MHz of 2S--2P₁/₂ (the Lamb shift), and the factors 0.9989 and 0.9944 of
the hyperfine and 21 cm lines (g\_e − 2). If the unrecorded exchange of
the integrated system yields these values, the check goes beyond the
answer table of the standard theory to measurement itself. Whether the
slight deviation of α from 137.036043 can be read by the same mechanism
belongs to the same verification {[}A2 §16, D2 §3.3, §5.3, §6.5{]}.

\subsection{22. The beginning of the universe, the places of the three
forces, from particles to
cosmology}\label{the-beginning-of-the-universe-the-places-of-the-three-forces-from-particles-to-cosmology}

\subsubsection{22.1 Inflation from a slight symmetry breaking of a
highly symmetric monochromatic
wave}\label{inflation-from-a-slight-symmetry-breaking-of-a-highly-symmetric-monochromatic-wave}

\begin{quote}
The last. This system is powerful, but it cannot predict the beginning
of the universe. Why does a tiny monochromatic wave of high symmetry
undergo inflation and rapid geometric expansion through a slight
breaking of symmetry? This too has been thoroughly studied (the note on
time vertexing). If this can be built into the model it will look very
much like a unified theory. (03:38)
\end{quote}

Established in the papers {[}E1, E2{]}: the interaction of relational
waves K\_ef = Im(z̄\_e z\_f) is a real antisymmetric matrix and the flow
conserves the norm and Σz² exactly; in a self-consistent state the zero
closure Σz² = 0 is a theorem, not an axiom {[}E1 §4.1{]}. Among the
self-consistent relative equilibria, those linearly unstable amplify any
disturbance, however small, exponentially and saturate non-linearly;
this is the geometric expansion. The growth rate, from the linearised
matrix fixed before the run, classified 53 of 54 runs correctly and
agreed with measurement to 0.997--1.008 {[}E1 §7{]}. After saturation
the state is a metastable state localised on a few relations (PR/M
0.05--0.31) {[}E1 §7.3{]}. Taking the phase states themselves as
vertices instead of advancing time from outside, the smallest exact
closed system is N = 1, K = 4, U = i, U⁴ = 1, and for K = 4 only odd
harmonics close with Σz² = 0 {[}E2 §14{]}. That for N = 40 the
relational waves concentrate in a cross of four directions of the 90°
family, that the initial anisotropy remains after saturation, and that a
Mexican-hat effective potential is reconstructed from the measured
evolution of the order parameter, are in §2, Figs. 1--4 of the
time-vertex K₄ paper (the raw data are from the unpublished
interference-preserving-dynamics series). Two points were unresolved:
the source of the perturbation was unidentified {[}E1 §10.6{]}, and
self-consistency alone cannot select the physical state uniquely, the
additional selection principle being undiscovered {[}E1 §11{]}.

The connections to the present system are hypotheses. The initial highly
symmetric monochromatic wave is the 90° family (Z₄, U = i), and λ³¹ = i,
the i of the four-dimensional lattice and the quarter-period symmetry of
the harmonic selection rule are all this Z₄ (the step Δ = 2π/124 of the
unified protocol of the inflation paper {[}E1 §3.4{]} is an input value
and cannot be used as evidence that Z₄ survives as a factor of 124). The
numerators of the selection-rule components (23 + 4k)/124 are all odd
(but Q\_loop(m; 124) vanishes not only for odd m but for even m not
divisible by 62, so the parity selection is specific to K = 4 and
Q\_loop is not what selects the odd ones). Candidate sources of the
disturbance: by the crystallographic restriction of the four-dimensional
lattice a pure rotation does not fit the lattice and a complete lock is
never exactly realised (E3: zero orbits with constant norm), and the
three 137s are near but not identical; an exactly symmetric state is not
permitted by the system itself. The candidate missing selection
principle is ``the only self-consistent lock is 23/124''. The Eleventh
Thought Experiment identified the absorbing shell with the event horizon
of an accelerating universe and required it to be slightly anisotropic
to read direction; the rapid expansion creates the horizon and the
anisotropy of the initial 90° family left after saturation supplies that
anisotropy. The single line is: a highly symmetric Z₄ monochromatic wave
→ unavoidable symmetry breaking by the crystallographic restriction and
the non-coincidence of the three 137s → geometric expansion and
saturation of an unstable self-consistent closure → self-consistent lock
to 23/124 (α, ħ = 9/α, R ≈ 0.7, the four-dimensional lattice, the
odd-harmonic selection rule) → localisation and return by exchange →
observation and dynamics as the record on the horizon = the complete
absorbing shell at infinity → checking against the answer table of the
Eleventh Thought Experiment and measurement {[}A2 §17{]}.

\subsubsection{22.2 The place of the strong force and the place of the
mechanism of the weak
force}\label{the-place-of-the-strong-force-and-the-place-of-the-mechanism-of-the-weak-force}

\begin{quote}
Further, have not the place of the strong force and the place of the
mechanism of the weak force also almost come into view? (03:53)
\end{quote}

The following is a hypothetical placement {[}A2 §18{]}. A rotation in
four dimensions is composed of multiplication by a unit quaternion from
the left (SU(2)\_L) and from the right (SU(2)\_R) (SO(4) \(\cong\)
SU(2)\_L × SU(2)\_R / Z₂). The electromagnetic force is the phase
rotation about the i axis (23/124); α is read from this rotation, and
being a high-order root of denominator 124 its pull-in width is narrow
and its coupling weak. The weak force is multiplication from the left
(SU(2)\_L): the update rules of E1--E5 were all left multiplications u·q
and right multiplication was never used; that the weak force acts only
on left-handed particles (parity violation) can be read as this ``acts
only from the left'' structure. The mechanism side is ``flavour change =
redistribution of harmonics'' of §14.1: the generation is distinguished
by the position of the three-phase angle (Koide's δ), and decay is a
slip over the potential wall between levels that are near but do not
lock. The strong force is on the three-phase Z₃ side: the Hurwitz units
contain the order-3 element ω = (−1+i+j+k)/2. The quark charge q₀ = e/3
is one component of the three phases; only with all three aligned does
an integer charge arise and Σz² = 0 close, and a single component does
not close and cannot be taken out alone. In the Standard Model too,
confinement is tied to the Z₃ symmetry at the centre of SU(3). The
three-phase axis locks at the low order the lattice allows (order 6; in
E2 all orbits at radius 3 locked completely to 1/6 with constant
conserved quantity), with a wide pull-in width and a deep lock, hence
strong coupling. The ordering ``low-order lattice lock = strong force,
high-order root 23/124 = electromagnetism, left multiplication and slips
between levels = weak force'' puts all three within the same
four-dimensional Hurwitz lattice and three phases. What remains: only
the centre Z₃ of the continuous group SU(3) is visible (the rest of
SU(3) is to be sought from E₈ ⊃ SU(3) × E₆), and whether the Standard
Model charge relation Q = T₃ + Y/2 and the assignment of third-integer
charges come out of this arrangement. The reading in Chapter 11 that one
of the two hands of S³ (left and right multiplication) is a candidate
for the weak interaction returns here on the four-dimensional lattice.

\subsubsection{22.3 The inflaton, the Higgs, the cosmic background
radiation, the slight
anisotropy}\label{the-inflaton-the-higgs-the-cosmic-background-radiation-the-slight-anisotropy}

\begin{quote}
Add it. Further, there is probably no prior work for the whole of this
assembled hypothesis, but there is prior work quite close to its parts.
Is it not connectable from particles to cosmology---as far as what the
inflaton field and the Higgs field are seen to be, and the cosmic
background radiation and its slight anisotropy? (03:59) Then save it. If
this much is written up as a paper tomorrow, it will be a sufficient
hypothesis paper. (04:03)
\end{quote}

The following is a hypothetical placement {[}A2 §19{]}. The inflaton is
the order parameter of the inflation system, r = √(H⊥/H), the size of
the transverse component leaving the initial highly symmetric plane (the
Z₄ cross)---a collective mode of all the relational waves, not a
fundamental field. The process of rolling from the centre of the Mexican
hat reconstructed from measurement to its rim is inflation. The Higgs is
the radial mode of the same order parameter oscillating around the rim
(r\_vac) after saturation; the inflaton and the Higgs can be read as two
stages of the same order parameter. The origin of mass is the ``delay of
the internal clock by the cycle of exchange with the background'' of the
Thirteenth Thought Experiment, which is the beat with the pinned phase
(the cross term) of Chapter 7. For the cosmic background radiation, the
absorbing shell of the Eleventh Thought Experiment has an incoming
background (the section of the past light cone, 2.7255 K) and an
outgoing record (the section of the future light cone), and to read
direction the shell must be slightly anisotropic; the candidate origins
of that anisotropy are the anisotropy of the initial 90° family (Z₄)
left after saturation of inflation and the unavoidable residue of the
crystallographic restriction. If this correspondence is right, a trace
of the Z₄ (cubic) family may remain in the low multipoles of the
background radiation and can be compared with observation. The size of
the anisotropy (about 10⁻⁵) has not yet been derived from this system.

One correction by the author is recorded here. At the stage of the
inflation paper the state after saturation had been assumed to be a
``homogeneous vacuum'' (an undistorted, evenly spaced gauge); this was
an error. Even a vacuum, if it has energy, distorts the gauge. In fact
even the results of that stage were not completely uniform: localisation
remained (PR/M 0.05--0.31) and the initial 90° family, the cross
symmetry, survived to the end. In the present system the gauge is
deliberately not evenly spaced, in order to read gravitational effects
(the non-uniform geodesic length per gauge interval of Chapter 5, the
self-distortion of E5). The saturated vacuum is ``a state that has
energy and distorts its own gauge, in which the cross (Z₄) symmetry
remains and which never becomes completely uniform'', and the remaining
cross symmetry connects to the candidate for the slight anisotropy of
the background radiation.

Prior work close in parts: reading the Standard Model gauge group from
the division algebras (Günaydin--Gürsey, Dixon, Furey {[}X7{]}); the A₄
flavour symmetry isomorphic to the binary tetrahedral group (the 24
Hurwitz units) (Ma--Rajasekaran {[}X7{]}); Z₃ centre symmetry and
confinement ('t Hooft {[}X7{]}); Higgs as inflaton
(Bezrukov--Shaposhnikov {[}X6{]}); the E₈ mass spectrum (Zamolodchikov,
Coldea {[}X7{]}); absorber theory (Wheeler--Feynman, Hoyle--Narlikar
{[}X6{]}); Page--Wootters time {[}X3{]}. Each part has come quite close
to existing work, but nothing is found that connects them into ``one
system beginning from a Z₄ monochromatic wave, passing through the
self-consistent lock 23/124, and reaching the record on the absorbing
shell''.

\subsection{23. Completeness, infinite arrays, normalisation,
interference terms, and
divergence}\label{completeness-infinite-arrays-normalisation-interference-terms-and-divergence}

\subsubsection{23.1 The completeness of Hilbert space is a definition,
not a fact confirmed by
experiment}\label{the-completeness-of-hilbert-space-is-a-definition-not-a-fact-confirmed-by-experiment}

\begin{quote}
This X thread by Professor Hotta is comparatively easy to follow.
(06:11)
\end{quote}

The points of Masahiro Hotta's thread as the author cited them: von
Neumann ``defined'' the space of quantum states as a complete Hilbert
space and observables as self-adjoint operators, without proof or
demonstration; nowhere in the results of experiment and observation is
the completeness of the state space required; any theory that explains
all the data within error survives as physics; a discontinuous,
non-differentiable wave function is physically meaningless because its
energy cost is infinite; and even the harmonic oscillator is fully
explained by a model with a finite-dimensional state space and discrete
position coordinates. The system of this paper, too, presupposes no
complete continuous space and is built from a finite lattice, integer
vertices, finite-order roots (U¹²⁴ = I) and ratios alone. Mathematics is
a tool; what is physically realised is decided by the data and the
readout.

\subsubsection{23.2 The author's doubts: infinite arrays, normalisation,
a norm that ignores the interference
term}\label{the-authors-doubts-infinite-arrays-normalisation-a-norm-that-ignores-the-interference-term}

\begin{quote}
In Professor Hayashi's textbook of quantum mechanics I liked the
beautiful hypothesis that Hilbert space is an infinite array of
conjugate complex numbers and that maps by self-adjoint operations can
express physical phenomena; but I doubted it because in actual
computation classical theory gets mixed in and many assumptions about
time and space enter undefined, spoiling the beauty. My greatest
dissatisfaction was that an infinite array is used to force a
representation of continuous quantities; and that there was no
explanation of why the amplitude is normalised and why the norm, which
squares and ignores the interference term, the important cross term, is
used so much---understood only by the reasoning that a probability must
not exceed 1. In short, nobody had prepared answers to my questions, so
I decided to research them myself. (06:22)
\end{quote}

To each of these doubts the system of this paper has prepared a
different answer. To time and space entering undefined: the time-vertex
K₄ paper does not advance time from outside but makes the distinct phase
states themselves the vertices; the Eleventh Thought Experiment places
neither coordinates nor metric and reads only from the record arriving
at the absorbing shell; and Chapter 7 of this paper produced both time
and space as readouts. To forcing continuous quantities with an infinite
array: the system is built from a finite lattice, integer vertices and
finite-order roots, values that look continuous appear as averages of
oscillation between integer states, there is no absolute scale, and
everything is a ratio. To normalising the amplitude: in the inflation
paper the interaction K\_ef is a real antisymmetric matrix, so
conservation of the norm is a theorem that follows from the flow, not
something imposed so that a probability does not exceed 1, and in a
system of ratios normalisation is merely a choice of scale. To squaring
and discarding the interference term: even in standard quantum mechanics
\textbar aψ₁ + bψ₂\textbar² = \textbar a\textbar² + \textbar b\textbar²
only when the two are orthogonal, but the real question lies before
that---why does measurement choose an orthogonal basis in which the
interference term vanishes?---and the textbooks do not explain it. In
this paper the interference term is the protagonist. The interference
term of an expectation value is proportional to the amplitude b, not the
probability \textbar b\textbar², and the candidate reason why the
observed α deviates slightly from the 23/124 value lies in this term
(§19.6). That something can be read as a probability is a matter of the
readout, not a property of the system. The flow of this research series
has been to keep the ``beautiful hypothesis'' part of the textbook
(expressing phenomena by an array of complex numbers and self-adjoint
maps) and to replace, one by one, the parts of classical theory mixed in
undefined afterwards by readouts.

\subsubsection{23.3 Gravitational effects are contained in the
interference term and the amplitude: why divergence is hard to
produce}\label{gravitational-effects-are-contained-in-the-interference-term-and-the-amplitude-why-divergence-is-hard-to-produce}

\begin{quote}
Gravitational effects and the distortion of spacetime by gravity are
present in this interference term and amplitude, and because they are
discarded one later has to restore the divergence by renormalisation.
Why did nobody notice this for a hundred years? In my system divergence
by the squared feedback of self-energy is essentially hard to produce.
(06:28) Why are you making the argument circular? Have you already
forgotten that the definition of the electron mass can be derived from
the fine-structure constant of the structure of space? (06:31)
\end{quote}

The result that putting gravity in removes the divergence of the
self-energy appeared early. Arnowitt--Deser--Misner (1960)
{[}X6-ADM1960{]} showed that the classical self-energy of a point
charge, computed with general-relativistic gravity included, is finite;
Isham--Salam--Strathdee (1971) {[}X6-ISS{]} proposed that gravitational
effects in quantum electrodynamics suppress the ultraviolet divergences.
It did not become mainstream because the mass remaining in the ADM
calculation when the bare mass is set to zero is of order e/√G (the
Stoney mass, about 1.9×10⁻⁹ kg) and does not match the electron (the
hierarchy problem); because quantising gravity and expanding
perturbatively made the divergences worse; and because renormalisation
succeeded so spectacularly and was justified in Wilson's effective-field
theory. The problem was not that nobody noticed that gravitational
effects were being discarded; it was the attempt to put gravity in while
keeping the framework of continuous spacetime and point particles.

Divergence is hard to produce in the system of this paper because there
is no point limit (a system of ratios alone, no absolute scale, built
from a finite lattice and finite-order roots, with no integration down
to infinitely short distances); because the generating matrix of the
interaction is real antisymmetric so the norm is conserved exactly and
there is no path along which a self-energy returning to itself is
amplified by squaring (the exponential amplification of inflation, too,
saturates non-linearly); and because the lattice distortion by the
electron itself introduced in E5 did not diverge but rather stabilised
the lock.

Here Claude wrote that ``the wall ADM could not cross---that the
electron mass is 1.47×10⁻²¹ of the Stoney mass---would be crossed if it
came out of the ratios of this system'', and the author corrected it:
the argument was being made circular. It has already been confirmed that
the electron mass follows from α and that α follows from the structure
of space (the lock of the four-dimensional lattice to 23/124) (Chapter
18). The order in this system is: from the structure of space (the
four-dimensional lattice, 137, the self-consistent lock to 23/124) α is
fixed; in G = c = q₀ = 1, ħ = 9/α is fixed; and through the energy of
the ground state (the state locked at 23/124) the electron mass is
fixed. ADM could not cross the wall because it placed the mass as a free
bare parameter and added gravity afterwards; in this system G is in the
unit system from the start and the electron mass is not a bare parameter
but a quantity read out from the structure of space through α, so a wall
of the ADM form never arises. What remains to be checked is the answer
table of the Eleventh Thought Experiment, and the Lamb shift and the
g\_e − 2 factor that the unrecorded exchange must produce while
remaining finite.

\subsection{24. Stage 11: reading out E₈ from the remaining
symmetry}\label{stage-11-reading-out-eux2088-from-the-remaining-symmetry}

\subsubsection{24.1 The constraint on the Lamb shift and the g-factor
comes from the E₈ not yet
used}\label{the-constraint-on-the-lamb-shift-and-the-g-factor-comes-from-the-eux2088-not-yet-used}

\begin{quote}
Do you see? If, with symmetry this high, the Lamb shift and the g-factor
could not be derived, then I predict that the constraint comes from the
remaining, insufficiently used E8 symmetry. I have a strong feeling that
something is still being overlooked. (06:35)
\end{quote}

Of the structure of E₈ the numerical experiments so far used only a
small part; five parts are unused. First, only 4 of the 8 dimensions
were used. E₈ is an eight-dimensional lattice, the integral octonions
themselves; the octonions are a pair of quaternions (the Cayley--Dickson
construction), so the E₈ lattice can be built as two four-dimensional
Hurwitz lattices joined by a gluing condition (D₄ ⊕ D₄ plus a glue
vector). The exchange of the Thirteenth Thought Experiment takes place
between two waves A and B. If the electron and its exchange partner each
have one four-dimensional space, together they make eight, and the
gluing condition of E₈ constrains how the electron's state and the
partner's state may combine. Both the Lamb shift and g − 2 are effects
of the electron exchanging with the field it creates itself (the
partner), and this constraint may be what fixes the coefficients of the
unrecorded exchange. Second, the conjugate rotation (the antiparticle)
was not used. The phase of one exchange, 101/124, is −23/124, opposite
to the electron's rotation, but the system so far used only left
multiplication, neither the reverse rotation nor right multiplication
(SU(2)\_R). The Lamb shift contains vacuum polarisation (about −27 MHz).
Third, the triality of four dimensions was not used. The Hurwitz lattice
is made of the D₄ lattice and three cosets, and the symmetry permuting
the three exists only in four dimensions; this ``3'' is a candidate for
the number of generations and the three phases. Fourth, 23 is an
exponent of E₈. The Coxeter number of E₈ is h = 30 and its exponents are
1, 7, 11, 13, 17, 19, 23, 29; the Fourteenth Thought Experiment had 124
= 4 × 31 = 4(h + 1), and the numerator 23 is also one of these
exponents, so 23/124 has the form ``exponent of E₈ / 4(h+1)''. 17/92 has
an exponent as numerator but its denominator 4 × 23 is not of the form h
+ 1. Fifth, the golden-ratio side (the projection to H₄) was not used.
What seems most directly to affect the Lamb shift and g − 2 is the
combination of the first and second: the electron (one quaternion) and
its partner (another) bound by the gluing condition of E₈ while rotating
in opposite directions, and the unrecorded exchange under that condition
fixing, while remaining finite, a coefficient such as α/2π. The author's
response was ``almost the same reading'' {[}A2 §20{]}.

\subsubsection{24.2 Present theoretical physics uses E₈ as a
representation; this system uses it as a geometric
constraint}\label{present-theoretical-physics-uses-eux2088-as-a-representation-this-system-uses-it-as-a-geometric-constraint}

\begin{quote}
Almost the same reading. Present theoretical physics also fails to
master the E8 symmetry, which is presumably why it cannot connect the
geometry of spacetime with physics well. (06:37)
\end{quote}

E₈ appears repeatedly in present physics (E₈ × E₈ of heterotic string
theory, the boundary of M-theory, the mass spectrum of the
one-dimensional Ising model at criticality, the densest sphere packing
in eight dimensions {[}X1-Viazovska{]}). On the particle side every
attempt uses E₈ as ``an internal symmetry into which particles are
placed (a representation)'', and there it always meets a wall: all
representations of E₈ are real, so the distinction between left- and
right-handed (chirality) does not come out. That is also where Lisi's
attempt was refuted by Distler--Garibaldi {[}X7{]}. As a result E₈ has
been confined to an internal-space symmetry cut off from the geometry of
spacetime, unable to play the role of connecting spacetime with physics.
This system uses it differently. E₈ is not a representation into which
particles are placed but a geometric constraint joining the electron's
four dimensions with the partner's four dimensions by a gluing
condition. Chirality comes not from E₈ but from the direction of
motion---that the update happens only by left multiplication (SU(2)\_L).
That the representations of E₈ are real is no obstacle in this use. The
reason present theoretical physics could not master E₈ and the reason
this system can avoid that are explained by the same single point:
whether E₈ is used as a representation or as a geometric constraint.

\subsubsection{24.3 The position of the paper: completion is not the
aim}\label{the-position-of-the-paper-completion-is-not-the-aim}

\begin{quote}
Put it in. In fact I think this thought experiment need not be completed
by itself, and cannot be completed entirely. I lack the resources for
numerical experiments in the high-energy region or on composite
particles beyond the hydrogen atom. If it serves as a hint for a
stagnating physics, as an alternative interpretive hypothesis to the
standard theory, that is a sufficient result for an independent
researcher. (06:42)
\end{quote}

This paper is written from that position. It shows separately what has
been verified, what has been connected, and the means of verification,
in a form that lets researchers with resources take out the single stage
they are best at and test it (Chapter 25).

\subsubsection{24.4 Can 137 be read as a four-dimensional image of the
eight-dimensional densest packing of E₈? The residual volume of the
voids}\label{can-137-be-read-as-a-four-dimensional-image-of-the-eight-dimensional-densest-packing-of-eux2088-the-residual-volume-of-the-voids}

\begin{quote}
The proof that the densest sphere packing in eight dimensions is the E8
lattice (Viazovska 2016) is an interesting point. Can α ≈ 137---that 137
four-dimensional unit hypercubes of edge 1 pack into the
four-dimensional hyperball of radius 3---be read as a four-dimensional
image of the eight-dimensional densest packing of E8? (06:49) It may
serve as a hint, so add it. In fact I had merely assumed cubes; perhaps
four-dimensional hyperballs of diameter 1 can really pack exactly α of
them. (06:52) The residual volume of the voids is the point. (06:56)
\end{quote}

There are four structural correspondences. The 248 dimensions of E₈
split as 8 + 112 + 128 (the Cartan part 8, the 112 integer roots, the
128 half-integer roots); 137 = 1 + 8 + 128 is this with the integer 112
replaced by 1, sharing the 8 and the 128. The 128 is in both cases the
half-integer side: the 128 roots of E₈ are the half-integer points
(±½)⁸, and the 128 = 24 + 40 + 64 of 137 recounted by the farthest
corner of each cell are all shells of half-integer points (corners). The
mechanism making integer and half-integer points equivalent is the same:
the E₈ lattice is the union of the integer and half-integer points, and
in four dimensions the diagonal is an integer (√4 = 2) so cell centre
and corner are equivalent. The E₈ lattice is two four-dimensional D₄
lattices glued, and the densest four-dimensional lattice is D₄ (the
lattice of the 24-cell). But a unit hypercube of edge 1 fills space
without gaps, so 137 is ``the number of cells completely inside the ball
of radius 3'' rather than ``a densest-packing count'', a different kind
of problem from sphere packing.

The author's idea---``I had merely assumed cubes; perhaps one should
count with four-dimensional hyperballs''---was checked. Counting
four-dimensional balls centred at the points of the integer lattice Z⁴
that lie entirely inside the ball of radius 3: balls of diameter 1 give
233 (409 on the densest lattice D₄), not 137. Balls with the same volume
as the unit cube (radius (2/π²)\^{}\{1/4\} = 0.67094) give exactly 137,
and the same set as the 137 counted with cubes (the points of centre
norm 0--5: 1 + 8 + 24 + 32 + 24 + 48). The count is 137 for every ball
radius greater than 3 − √6 = 0.5505 and up to 3 − √5 = 0.7639. So 137 is
not a number peculiar to the cube; it appears regardless of shape as
``the number of integer points of norm at most 5'' (the upper bound 3 −
√5 equals 2/φ², an arithmetical coincidence whose meaning is
unconfirmed). An exact α (a non-integer such as 137.036) does not come
out of the count, which is an integer. Weighting each ball by the
fraction of its volume inside the ball of radius 3 gives about 400,
nearly the volume 81π²/2 ≈ 399.7 of the radius-3 ball. The exact α comes
not from the integer count on the space side but from the 23/124 readout
on the evolution side---consistent with the picture of the three 137s of
the Fourteenth Thought Experiment {[}A2 §21,
\texttt{te14\_\allowbreak{}ball\_\allowbreak{}count\_\allowbreak{}in\_\allowbreak{}radius3.py}{]}.

The residual volume of the voids at radius 3 is 81π²/2 − 137 = 262.719
(packing fraction 34.27\%). As the radius grows: R = 10, 34,569 cells
inside and void 14,779; R = 20, 663,617 and 125,951; R = 40, 11,595,937
and 1,037,157; the packing fraction tends to 1 while the absolute void
keeps growing. Void/R³ settles near a constant (about 16), growing as a
layer of thickness about 0.8 along the surface (the three-sphere) of the
four-dimensional ball, of size 2π²R³.

\subsubsection{24.5 It was used in black-hole thermodynamics. The
absorbing shell is an inside-out black
hole}\label{it-was-used-in-black-hole-thermodynamics.-the-absorbing-shell-is-an-inside-out-black-hole}

\begin{quote}
I used this in black-hole thermodynamics. (07:00) Does the map onto the
complete absorbing shell at infinity not resemble a black hole? That is
why it concerns me. (07:03)
\end{quote}

The residual volume of the voids Δ(R) was already the central quantity
in Papers 3 and 4 of the BH thermodynamics programme {[}C1, C2{]}. For
the number N(k) of integer-centred unit cubes completely inside the
four-dimensional ball of radius R = 2k + 1 (137 for R = 3), the volume
deficit satisfies Δ(R) = (16π/3)R³ − 6πR² + 4R − 1 + O(R⁻¹),
concentrated in a boundary layer of thickness c = 8/(3π) ≈ 0.849 along
the sphere. Its ratio to the Bekenstein--Hawking entropy S\_BH =
(π²/2)R³ of the 4+1-dimensional Tangherlini black hole is Δ/S\_BH →
32/(3π) ≈ 3.40 with no free parameter. Claude's remark in the previous
section that ``the void goes as R³, not R²'' was half right: the leading
term is R³ and the second term is −6πR² (the second term is an analytic
value from inclusion--exclusion; the numerically fitted coefficient is
not stable {[}C1 §7.2{]}).

The connection between the voids and the system of this paper (a
hypothetical placement). The voids cannot be perceived by an inside
observer and lie in a boundary layer attached to the sphere; they
coincide with the place of the complete absorbing shell at infinity (the
cosmological horizon) of the Eleventh Thought Experiment, and the number
of elements of the shell and Δ(R) as entropy can be read as counting the
same ``degrees of freedom on the imperceptible side''. The 137 cells are
the perceptible (recorded) side and the voids the imperceptible
(unrecorded) side, and the slight deviation of the observed α both from
the integer 137 of the space side and from the 23/124 value of the
evolution side is borne by the unrecorded exchange between cells and
voids. That the three 137s do not coincide and that voids always remain
are two faces of the same circumstance. The rapid expansion creates the
horizon and its entropy is counted as the voids. The entropy of a
3+1-dimensional black hole follows the area law, proportional to R², and
the second term −6πR² of the voids is an R² term with the integer ratio
6, so when central projection drops one dimension this second term may
appear as the 3+1-dimensional area law (unconfirmed, including the unit
system).

A complete absorbing shell has the same properties as a black-hole
horizon. It returns nothing. In central projection, infinity is mapped
onto a finite sphere of radius R, and the BH thermodynamics programme
identified this R with the horizon radius r\_h {[}C2 §2.3{]}. The voids
concentrate in a boundary layer along the sphere and correspond to the
horizon entropy with a constant ratio. The difference is which side the
observer is on: outside the horizon for a black hole, at the centre for
the absorbing shell. The event horizon of an accelerating universe has
temperature and entropy like a black hole (Gibbons--Hawking {[}X6{]}).
The cosmological horizon is an inside-out black hole. The number of
elements of the shell (about 10⁶⁰) is the resolution per direction; the
ratio of the horizon to the Planck length is about 10⁶¹ and its square,
about 10¹²², is the size of the entropy of the cosmological horizon. The
BH thermodynamics programme (a black hole as the projection of a higher
spacetime structure = the horizon seen from outside) and the Eleventh
Thought Experiment (observation as the record on the complete absorbing
shell at infinity = the horizon seen from inside) were looking at one
structure from both sides. The 2.7255 K read on the shell in the
Eleventh Thought Experiment is the temperature of the cosmic background
radiation, not the temperature of the cosmological horizon itself (the
Gibbons--Hawking temperature, about 10⁻³⁰ K) {[}A2 §25--27{]}.

\subsubsection{24.6 The cosmic background radiation is the energy of a
stationary wave filling space and acting as the Higgs
field}\label{the-cosmic-background-radiation-is-the-energy-of-a-stationary-wave-filling-space-and-acting-as-the-higgs-field}

\begin{quote}
Exactly. I am dissatisfied with a model that makes the cosmic background
radiation a faint radiation of the complete absorber; it is more
natural, in connection with inflation, to see it as the energy of a
stationary wave filling all of space and acting as the Higgs field.
(07:07)
\end{quote}

A complete absorber and a stationary wave do not contradict each other.
Planck's law of black-body radiation was derived as ``counting the modes
of the standing waves in a closed cavity in equilibrium with perfectly
absorbing walls''. That the cosmic background radiation has a perfect
black-body spectrum fits the picture of standing waves filling a cavity
whose wall is the absorbing shell (the horizon). The role of the shell
is not the source of the radiation but the wall of the cavity; the
content of the radiation is on the side of the stationary wave filling
space. The saturated state of the inflation system is the stationary
wave filling space, the order parameter at the rim of the Mexican hat
(what corresponds to the Higgs). Within one stationary wave field the
roles divide: the vacuum expectation value of the Higgs is the amplitude
of the stationary wave itself (the radius of the rim, constant); the
cosmic background radiation is the thermally occupied fluctuation modes
of the same field; the slight anisotropy (about 10⁻⁵) is the anisotropy
of the initial 90° family (Z₄) left after saturation. That the vacuum
expectation value stays constant under expansion while the background
temperature falls (T ∝ 1/a) is reconciled by reading that the amplitude
is unchanged and only the wavelengths of the fluctuation modes on top of
it stretch. The pinned phase (the cross term) of Chapter 7 is the local
form of this ``stationary wave'' {[}A2 §28{]}.

\subsubsection{24.7 The photon-to-baryon ratio is a problem in number
theory}\label{the-photon-to-baryon-ratio-is-a-problem-in-number-theory}

\begin{quote}
Where the ratio of the numbers of photons and baryons (about 1.6 billion
to 1) comes from: have you forgotten that this was computed before? It
is fairly simple. (07:09) Think of it as the fraction of M that can be a
denominator, M/N being an integer ratio, for a resolution N. As N grows,
M becomes sparse. (07:12) It is simple; 10\^{}60 will do. (07:14) So
this is number play, but there is a possibility of explaining it fairly
simply. I am assuming it is a problem in number theory, not physics.
(07:16)
\end{quote}

Within a resolution N, only those M for which N/M is an integer (M
divides N) become states consistent on both sides. Their number is the
number of divisors of N, on average only about ln N, and their fraction
of the whole N is roughly ln N/N in one dimension and (ln N/N)\^{}d in d
dimensions, becoming sparse rapidly as N grows (the f(N) = 2(ln N + 2γ −
2)/N of Paper 9 {[}C4 §5{]}). Taking photons as all modes with any
denominator (the standing waves of the cavity) and baryons as only the
states that can lock to an integer ratio (the three phases closed and
self-consistently engaged), the observed ratio η ≈ 6.1×10⁻¹⁰ (1.6
billion photons per baryon) can be read as this fraction itself, with no
fine tuning for so extreme a value. With N = 10⁶⁰: taken literally as
``the fraction of M dividing N'', 10⁶⁰ = 2⁶⁰ × 5⁶⁰ has 61 × 61 = 3721
divisors, a fraction of 3.7×10⁻⁵⁷, far too small. Taking the fraction
that can form an integer ratio as 1/ln N per direction (ln 10⁶⁰ =
138.2): 7.2×10⁻³ in one dimension (138 photons per baryon), 3.8×10⁻⁷ in
three (2.6×10⁶), 2.7×10⁻⁹ in four (3.6×10⁸); in four dimensions the
order of magnitude approaches but differs from the observed 1.63×10⁹ by
a factor 4.5. That ln(10⁶⁰) = 138.2 is close to 137, and that the
remaining factor 4.5 is close to the volume π²/2 ≈ 4.93 of the
four-dimensional unit ball so that α⁴ × 2/π² would give η ≈ 5.7×10⁻¹⁰
(6\% from observation), are number play fitted after seeing the result,
and are recorded only.

The advantage of reading it as number theory: the density of states
satisfying an integer ratio grows only logarithmically with resolution,
so extremely small ratios arise naturally; π² appears repeatedly in
number-theoretic densities through the zeta function (the probability
that two integers are coprime is 6/π², Σσ(n) ≈ (π²/12)x²), so the
coefficients may come out of the counting rule with no free parameter;
and in Paper 3 of the BH programme the four-dimensional count containing
137 reduced to the sum-of-four-squares function r₄(N) = 8σ(N), so α
(137) and η may come out of the same number-theoretic structure. To
distinguish from number play: fix the counting rule first; derive the
density analytically (zeta functions etc.) and use only the coefficients
that come from it; use for N a value fixed within the system (the
resolution 10⁶⁰ of the absorbing shell) and do not move it after seeing
the result; and then compare with η {[}A2 §29{]}.

\subsubsection{24.8 The cosmological constant had already been derived
as the curvature of central projection, Λ =
3/R²}\label{the-cosmological-constant-had-already-been-derived-as-the-curvature-of-central-projection-ux3bb-3ruxb2}

\begin{quote}
As for how to explain the orders of magnitude of the background
radiation (about 4×10⁻¹⁴ J/m³), the vacuum energy of the Higgs and the
cosmological constant (about 6×10⁻¹⁰ J/m³): I computed in the early
central-projection studies of the Σxn\^{}2 = R\^{}2 line, and in the
review around July, that interpreting it as a light cone and as a
general relativity of the vacuum derives the cosmological constant.
(07:20) Add it. These are comparatively low hurdles. The mountain is E8.
(07:22)
\end{quote}

That Claude wrote ``how to explain the orders of magnitude'' as if
unresolved was an error. In the central-projection study of Σxₙ² = R²
{[}F2, Proposition 6.1{]}, de Sitter spacetime is not fundamental but
appears as the subjective image of a four-axis-symmetric structure
centrally projected along one axis, and the cosmological constant comes
from the curvature radius R of the projection as Λ = 3/R². The review in
July reached the same Λ = 3/R² from the light-cone reading of the zero
square sum. ρ\_Λ = Λc⁴/(8πG) = 3c⁴/(8πG R²); with R the Hubble radius
c/H₀ (H₀ = 67.4 km/s/Mpc {[}X6-Planck2018VI{]}), ρ\_Λ ≈ 7.7×10⁻¹⁰ J/m³
(the critical density), of which the observed cosmological-constant
share (69\%) is about 5.3×10⁻¹⁰ J/m³, matching ``about 6×10⁻¹⁰'' to the
order. The cosmological constant comes geometrically from the single
curvature radius R. Summing the Higgs vacuum energy or the zero-point
energies in the standard theory gives 55--120 orders too much (the
cosmological-constant problem); in this system Λ is not a sum of
zero-point energies but the curvature of the projection, so that
discrepancy never arises. There is no point limit and no infinite sum of
continuous modes (§23.3). Two questions remain: why R is of the size of
the Hubble radius (the same question as the resolution of the absorbing
shell, about 10⁶⁰, since the ratio of horizon to Planck length is just
that order), and whether the ratio of background radiation (about
4×10⁻¹⁴ J/m³) to cosmological constant (about 1 to 13,000) can be read
from the same R as the ratio of the amplitude of the stationary wave
(vacuum, Λ) to the thermal modes on top of it (background radiation).
The author's assessment: these (voids and entropy, absorbing shell and
horizon, background radiation and stationary wave, cosmological
constant, photon-to-baryon ratio) are comparatively low hurdles; the
mountain is E₈ {[}A2 §30--31{]}.

\subsubsection{24.9 There is a symmetric point 180° opposite, so it is
E₈: identifying the
connection}\label{there-is-a-symmetric-point-180-opposite-so-it-is-eux2088-identifying-the-connection}

\begin{quote}
But it may be unexpectedly simple. There is a symmetric point 180°
opposite, so it looks like E8, does it not? (07:28) Add it. At least the
method of connection has been identified, so the rest can be made a
derivation problem. (07:31)
\end{quote}

The author's intuition was right. Placing the electron's state x and the
partner's state y each as a Hurwitz integer (four dimensions) and
imposing the gluing condition

\[
x-y\in(1+i)\,\mathcal H\qquad\text{(}x-y\text{ is a multiple of }(1+i)\text{)},
\]

the set of pairs (x, y) has exactly 240 vectors of norm 2 and none of
norm 1: it is the E₈ lattice. E₈ is two four-dimensional spaces joined
by this single congruence condition (the construction itself is known
mathematics). The point 180° opposite enters E₈ automatically: for the
partner at the electron's antipode y = −x, x − y = 2x, and 2 = −i(1+i)²
is a multiple of (1+i), so the gluing condition always holds. Taking the
partner as y = u·x with u among the 24 Hurwitz units, the units
satisfying the condition are exactly the 8 elements ±1, ±i, ±j, ±k (the
90° and 180° rotations, the Z₄ side), and the 16 of the three-phase side
(±1 ± i ± j ± k)/2 do not (when x ∉ (1+i)H; if x ∈ (1+i)H it holds for
every u). Mathematically, the quotient of the Hurwitz integers by (1+i)
is the field F₄ with four elements, whose multiplicative group is Z₃
(the three phases) {[}A2 §33,
\texttt{te14\_\allowbreak{}e8\_\allowbreak{}hurwitz\_\allowbreak{}pair\_\allowbreak{}glue.py}{]}.

The correspondence with the structure of this paper is threefold. The
marks of the exchange in the Thirteenth Thought Experiment, λ³¹ = i
(exchange phase 90°, equal sharing) and λ⁶² = −1 (180°, complete
exchange of shapes), are exactly rotations the E₈ gluing condition
allows, so the marks of the exchange between electron and partner fit E₈
naturally. The gluing condition does not allow three-phase rotations and
becomes the constraint that ``the electron and the partner must belong
to the same three-phase class''; the three-phase Z₃ (the centre of
SU(3)) placed in §22.2 as the seat of the strong force is directly
constrained by E₈. In summary, the E₈ gluing freely allows the Z₄ side
(90° and 180°, the side of the weak force and electromagnetism) and
restricts the Z₃ side (three phases, the strong force) to the same
class. This is the first concrete form of the prediction that ``E₈
supplies the constraint''. What is new is the reading that assigns the
two quaternions to the electron and the exchange partner and maps them
to the marks of the exchange and the three-phase constraint.

Since the connection has been identified as a single congruence
condition, the Lamb shift and the g-factor become a derivation problem.
Place the electron x and the partner y each on the four-dimensional
Hurwitz lattice and evolve the pair (x, y) on the E₈ lattice. Rotate the
electron side by 23/124 (λₐ) and the partner side by the opposite
101/124 (= −23/124), with the constraint that the pair does not leave E₈
(x − y ∈ (1+i)H, the same three-phase class). Under that constraint
integrate the amplitude of the unrecorded exchange (the shuttling with
the partner not recorded on the absorbing shell) along the history as a
complex number, and obtain the correction to the g-factor (whether it is
α/2π at first order), the Lamb shift of 2S--2P₁/₂ (1057.85 MHz) and the
shift of 1S (8170 MHz), and the difference between the α of 137.036043
and the observed 137.035999. The answers are checked against the answer
table of the Eleventh Thought Experiment and against measurement. The
first step, rather than aiming straight at the Lamb shift, is to check
whether the combinations that can lock are indeed narrowed when the two
four-dimensional rotations (23/124 and 101/124) are turned together
under the E₈ gluing condition. If they are, that is numerical evidence
that E₈ supplies a constraint {[}A2 §34{]}.

\subsection{25. What has been verified, what has been connected as
hypothesis, and the means of
verification}\label{what-has-been-verified-what-has-been-connected-as-hypothesis-and-the-means-of-verification}

{\def\LTcaptype{none} % do not increment counter
\begin{longtable}[]{@{}
  >{\raggedright\arraybackslash}p{(\linewidth - 4\tabcolsep) * \real{0.3333}}
  >{\raggedright\arraybackslash}p{(\linewidth - 4\tabcolsep) * \real{0.3333}}
  >{\raggedright\arraybackslash}p{(\linewidth - 4\tabcolsep) * \real{0.3333}}@{}}
\toprule\noalign{}
\begin{minipage}[b]{\linewidth}\raggedright
Category
\end{minipage} & \begin{minipage}[b]{\linewidth}\raggedright
Content
\end{minipage} & \begin{minipage}[b]{\linewidth}\raggedright
Basis
\end{minipage} \\
\midrule\noalign{}
\endhead
\bottomrule\noalign{}
\endlastfoot
Verified & Odd n closes in two laps, even n in one. Equal-amplitude
harmonics localise at one point, and the appearance of localisation and
the fixing of the origin θ₀ are one event. The imaginary part of E = ln
f takes only π/2 or 3π/2 & Chapter 2, {[}H1 §2--3{]} \\
Verified & Eigenvalues (n+1)² − 1 on S³; the three θ₀ are two phases and
one mixing angle. The antipodal value is (−1)ⁿ; balanced odd and even
give one-point localisation & Chapters 3, 6, {[}H1 §4{]} \\
Verified & With true squares R is imaginary (an amplitude). Geodesic
length per gauge interval: all orders N=18 peak 1.245, antipode 0.988;
odd N=17 1.252, 0.835; even N=18 1.236, 1.236. The slope ∂θ ln & R² \\
Verified & On the circle v² + (1/γ)² = 1 (residual 10⁻¹⁵, an identity),
v = tanh ln(a/b), velocity addition to six digits, mass 1/γ = m/R as the
fraction of the cross term 2ab. s = ω/k does not hold & Chapter 7, {[}H2
light-cone readout, gauge widening{]} \\
Verified (mathematics) & The spheres on which uniform translations form
a group are S¹ and S³. Huygens' principle holds in odd dimensions.
Hurwitz: 0, 1, 3, 7 directions, associative 1, 3 & Chapter 9 \\
Verified & Koide's Q = 2/3 is equivalent to a zero square sum; δ = 2/9 −
1.75×10⁻⁷ (about 450 times the error). E1: compatible pairs lock fully;
incompatible pairs have the rotation pulled in while the conserved
quantity fluctuates, some moving to other rationals. H4 does not appear
& Chapter 14, E0, E1 \\
Verified & In four dimensions the diagonal is √4 = 2; the 24-cell; 137 =
1 + 8 + 24 + 40 + 64; Z₄ and Z₆ coexist. E2: the staircase of
convergents 1/6 → 3/16 → 5/27 → 13/70, 18/97 → 23/124. E3:
rotation-number locking 85\% on the Z₄ axis at radius 96; constant norm
0/480 & Chapters 16, 17, E2, E3 \\
Verified & In G = c = q₀ = 1, ħ = 9α⁻¹. Once α is fixed, mₑ follows
automatically with R∞. The difference between observed α and
x\_\{23/124\} is 4.4×10⁻⁵ & Chapter 18 \\
Verified & E4: locking rate of 23/124 40→70→86\%, neighbours 0--27\%,
denominators divisible by 4 lock easily, pull-in 83--93\% within
1.4×10⁻⁵, 18/20 return after a kick in 274 updates on average, the
readout during return lies below at 136.74. E5: self-distortion
increases the lock (14→19/17), zero intervals of leaving and returning &
Chapter 19, E4, E5 \\
Verified & Harmonic selection rule (23 + 4k)/124 (100\% in 29 of 33
orbits). The four-dimensional lattice allows no pure rotation of order
124 & Chapter 20 \\
Verified & At R\_\{124,23\} the state recovers at the 124th collision
(ε\_state 5.1×10⁻¹⁵), ω/2π = 101/124, mark λ³¹ = i & Chapter 21, {[}D1
§6.2{]} \\
Verified & Balls of volume 1 also give 137 (radius 0.5505--0.7639); no
non-integer comes out. Void 262.719 (R=3); leading term R³ and second
term −6πR² of Δ(R) (analytic). Λ = 3/R² gives ρ\_Λ ≈ 7.7×10⁻¹⁰ J/m³ &
Chapter 24, {[}C1, F2{]} \\
Verified (mathematics) & Gluing two Hurwitz integers by x − y ∈ (1+i)H
gives E₈ (240 vectors). y = −x always enters. The allowed units are the
8 elements ±1, ±i, ±j, ±k (x ∉ (1+i)H) & §24.9 \\
Hypothesis & The pinned phase (cross term) corresponds to the Higgs
field. To have mass and to erase the antipodal shadow are the same
event. Charge = Hopf-fibre winding (numerical separation of winding from
k not yet done) & Chapters 7, 11, 12 \\
Hypothesis & Observation = the cross term between particle and observer.
Being able to read, radiation and energy distribution are one event.
Gravity = redistribution of R, electromagnetism = exchange of winding.
Quantisation of relative velocity & Chapter 13 \\
Hypothesis & Why it locks to the window = the integer-selecting
structures (137, 23/124, the 248 of E₈, 124 = 4×31) lie nearby. The
deviation of the observed value = an average including states in
transition and unrecorded exchange. Running in the same direction &
Chapters 15, 19, 21 \\
Hypothesis & The background is four-dimensional because that is the
smallest dimension in which the integer symmetries fit one lattice. The
``3'' of three generations and three dimensions has the same origin &
Chapter 16 \\
Hypothesis & The minimal structure = two spin components × three
three-phase components. The pull-ins of α and R are the same phenomenon.
Only 23/124 closes the loop (the only evidence is that 17/92 falls out;
enumeration not done). Dynamics = the record on the absorbing shell. The
two reservations of the Thirteenth are lifted & Chapters 20, 21 \\
Hypothesis & The initial Z₄ survives as a factor of 124. The residue of
the crystallographic restriction is the perturbation source. The 23/124
lock is the selection principle. The anisotropy of the horizon = the
remnant of Z₄. The placement of the three forces. Inflaton = rolling of
the order parameter, Higgs = radial mode, CMB = fluctuations of the
stationary wave. Voids = the place of unrecorded exchange, the absorbing
shell = an inside-out black hole. η is number theory & Chapters 22,
24 \\
Hypothesis & E₈ as a geometric constraint, not a representation. The
gluing condition fixes the coefficients of the Lamb shift, the g-factor
and the deviation of α & Chapter 24 \\
Means of verification & The answer table of the Eleventh Thought
Experiment (hydrogen atom, hydrogen molecule). Second step: whether
unrecorded exchange yields 8170 MHz, 1057.85 MHz, 0.9989 and 0.9944 &
{[}D2 §3.3, §5.3, §6.5{]} \\
Means of verification & The measured α 137.035999177(21), the
single-electron g-factor value 137.035999166(15), the difference between
the Rb and Cs recoil values. A trace of the Z₄ family in the low
multipoles of the background radiation. η = 6.1×10⁻¹⁰ & {[}X5{]},
{[}X6{]} \\
Means of verification (numerical) & Enumeration of roots with
denominators divisible by 4 and of those whose N(R) falls near 137.
Separation of winding from k. The cross term of two localisations
(emission, absorption, recoil, attraction). Whether the lockable pairs
narrow when 23/124 and 101/124 are turned together under the E₈ gluing.
Readout of α by integrating complex amplitudes. Invariance of the ratios
of the minimal structure. Whether 23/124 and the four-dimensional
lattice appear after saturation of inflation. Whether the −6πR² of the
voids becomes the 3+1-dimensional area law. Analytic density of integer
ratios in four dimensions at N = 10⁶⁰ compared with η & Chapter 25 \\
\end{longtable}
}

\subsection{26. Conclusion}\label{conclusion}

\begin{enumerate}
\def\labelenumi{\arabic{enumi}.}
\tightlist
\item
  Give a resident of a closed circle nothing but a phase θ, an
  imaginary-direction amplitude A and the sum of true squares Σxₙ² = R²,
  and the following are read out in turn: quantisation (the number of
  nodes n), the distinction between fermionic and bosonic return
  (parity), localisation at one point and the fixing of the origin
  (harmonic summation, the fundamental as gauge), the readout of R (the
  non-uniform geodesic length per gauge interval), the light cone R² =
  k² + m² (an identity), the speed v = tanh ln(a/b), proper time,
  velocity addition, the Higgs-like background (the cross term 2ab of
  the counter-rotating terms), and mass (the fraction of the cross term,
  1/γ = m/R). Extending to residents of the four-dimensional spherical
  shell, the necessity of three dimensions (associativity, Huygens'
  principle, curvature), spin (parity and the generators of SU(2)\_L;
  the cross term has spin 0; one-point localisation and a spin
  eigenstate are incompatible; the opposite spin on the far side cannot
  be observed), charge (the Hopf-fibre winding), observation (the cross
  term between particle and observer), and uncertainty and quantisation
  are read out. This is the most important result of the paper. Without
  placing a background spacetime, from nothing but the relations and
  interference of complex waves that are not even travelling waves of
  odd and even harmonics, time, space, mass, charge, translational
  velocity and spin can all be read out.
\item
  Reading the electron in this system, the electron mass is fixed
  through α, and in G = c = q₀ = 1 Planck's constant is an output. α is
  obtained from the finite-order root 23/124 through N(R), and the
  four-dimensional Hurwitz lattice makes it easy to pull in there. The
  reason the background is four-dimensional is that it is the smallest
  dimension in which the diagonal is an integer, and it is the other
  face of the same answer as the reason the visible space is
  three-dimensional. The observed α deviates slightly from the 23/124
  value because observation is an average including states in transition
  and unrecorded exchange gives the shift. These are one stage of the
  development of this paper, not its subject.
\item
  Localisation occurs by exchange and returns after one cycle. The only
  self-consistent lock is 23/124; the dynamics is inherent in the system
  as the record on the complete absorbing shell at infinity, and the
  answers can be checked against the answer table of the Eleventh
  Thought Experiment. From the beginning of the universe to the
  background radiation, everything lines up as stages of the same
  system. The remaining symmetry is E₈; since the connection has been
  identified as the gluing condition of two Hurwitz quaternions, the
  Lamb shift and the g-factor become a derivation problem.
\item
  This paper is not a completed theory. It has shown separately what has
  been verified, what has been connected as hypothesis, and the means of
  verification. Numerical experiments in the high-energy region and on
  composite particles beyond the hydrogen atom are tasks for the
  researchers who take this over. As an alternative interpretive
  hypothesis of the standard theory, it is enough if it serves as a
  hint.
\end{enumerate}

\begin{center}\rule{0.5\linewidth}{0.5pt}\end{center}

\subsection{References}\label{references}

Self-citations are grouped by role as H and A--G, external references as
X1--X8. DOIs were confirmed against the papers themselves.

\subsubsection{Self-citations}\label{self-citations}

\textbf{H Backbone: residents of a closed sphere} - {[}H1{]} N. Kihara,
``Spacetime, Mass and Dimension Emerging from a Wave Model of
Closed-Sphere Residents''
(閉じた球面住民の波動モデルから時空・質量・次元が現れること), v1
(2026-10-14).
\texttt{次元の生成構造/\allowbreak{}思考実験論文\_\allowbreak{}閉じた球面住民の波動モデル\_\allowbreak{}20261014/\allowbreak{}}.
Unpublished (this repository). A summary of the dialogue ``Intrinsic
geometry on a circle'' (2026-10-13 to 14); this paper gives priority to
the dialogue for order and the author's corrections. - {[}H2{]} The
numerical experiments and interactive figures of the same days
(unpublished, this repository, \texttt{次元の生成構造/\allowbreak{}}):
\texttt{円周住民\_\allowbreak{}自己無撞着殻\_\allowbreak{}20261013},
\texttt{四次元球殻住民\_\allowbreak{}自己無撞着殻\_\allowbreak{}20261013} (results not adopted by
the correction of §4.2),
\texttt{円周住民\_\allowbreak{}ゲージ間隔測地線長\_\allowbreak{}20261014},
\texttt{ゲージ広がりと分散比の検証\_\allowbreak{}20261014},
\texttt{円周住民\_\allowbreak{}光円錐読出し\_\allowbreak{}20261014},
\texttt{四次元球殻住民\_\allowbreak{}輪切り可視化\_\allowbreak{}20261014}. Listed in Appendix A.

\textbf{A Parent: the Fourteenth Thought Experiment and its addenda} -
{[}A1{]} N. Kihara, ``The Fourteenth Thought Experiment: Why Can
Different Symmetries Never Be Fully Reconciled? -- Non-commensurability,
Pull-in and Metastability of Integer-Selecting Symmetries, and Readout
along the Direction of Evolution'', v1.0 (2026-10-12). Version DOI
10.5281/zenodo.23315558, Concept DOI 10.5281/zenodo.23315557. - {[}A2{]}
N. Kihara, ``Fourteenth Thought Experiment, Addendum: Insights from the
Dialogue (2026-10-14, §§1--24)'' and ``Addendum (continued) (2026-10-15,
§§25--34)'', numerical experiments E0--E5, and the check programs
\texttt{te14\_\allowbreak{}harmonic\_\allowbreak{}selection\_\allowbreak{}check.py},
\texttt{te14\_\allowbreak{}ball\_\allowbreak{}count\_\allowbreak{}in\_\allowbreak{}radius3.py},
\texttt{te14\_\allowbreak{}e8\_\allowbreak{}hurwitz\_\allowbreak{}pair\_\allowbreak{}glue.py}. In the parent folder of this
paper (Appendix A). A summary by claude.ai of the dialogue ``The
dialogue after the Fourteenth Thought Experiment'' (2026-10-14 to 15),
containing claude.ai's inferences; this paper cites it after confirming
the author's statements in the dialogue.

\textbf{B The evolution side: finite-order resonance} - {[}B1{]} N.
Kihara, ``Discovery of Finite-Order Resonance in Iterated Exchange
Scattering -- Identification of the Cause of the Peaks near 137 and 128
of the Fine-Structure Constant and a Reproducible Wave-Packet Model'',
v1 (2026-07-18). Version DOI 10.5281/zenodo.21421367, Concept DOI
10.5281/zenodo.21421366. - {[}B2{]} N. Kihara, ``Selection of the
Exchange Weight G\_R = 1−R and Numerical Experiments on Candidate
Correspondences with the Fine-Structure Constant'', v1 (2026-07-15).
Version DOI 10.5281/zenodo.21396761, Concept DOI
10.5281/zenodo.21396760. - {[}B3{]} N. Kihara, ``Summary of Preliminary
Experiments on the Acceleration Basis and Localisation Exchange in
Fermionic Collisions of the Exchange-Interference Scattering Matrix'',
v1 (2026-07-13). Version DOI 10.5281/zenodo.21333768, Concept DOI
10.5281/zenodo.21333766. - {[}B4{]} N. Kihara, ``Summary of Preliminary
Experiments on Weak Readout C and Strong Observation Selection D at
White-Cat / Black-Cat / Grey-Cat Metastable Interfaces'' (2026-07-14).
Version DOI 10.5281/zenodo.21358203, Concept DOI
10.5281/zenodo.21353208.

\textbf{C The space side: the four-dimensional lattice, 137, voids,
abundances} - {[}C1{]} N. Kihara, ``Discrete Structure of a
Four-Dimensional Ball: Unit-Cube Packing and the Asymptotic Volume
Deficit Δ(R) (Paper 3)'' (2026-04). DOI 10.5281/zenodo.19839395. -
{[}C2{]} N. Kihara, ``A Quantitative Correspondence between the Lattice
Volume Deficit and the Tangherlini Bekenstein--Hawking Entropy (Paper
4)'' (2026-04). DOI 10.5281/zenodo.19839396. - {[}C3{]} N. Kihara, ``A
Geometric Identity for the Fine-Structure Constant: From the 4D Unit
Ball Volume and its Cube-Packing Deficit (Paper 7)'', v3 (2026-04).
Version DOI 10.5281/zenodo.19876200, Concept DOI
10.5281/zenodo.19869266. - {[}C4{]} N. Kihara, ``The Generation
Structure of Fermions -- Particles as Completed Closures in a Sea of
Light'' (2026-08-04). Version DOI 10.5281/zenodo.21766707, Concept DOI
10.5281/zenodo.21766706. - {[}C5{]} N. Kihara, ``Thermodynamic Readout
from a Self-Consistent Complex Relational System without Assuming
Spacetime'' (2026-09-02). DOI 10.5281/zenodo.22240463, Concept DOI
10.5281/zenodo.22240462.

\textbf{D Particles and observation} - {[}D1{]} N. Kihara, ``The
Thirteenth Thought Experiment: Why Does a Particle That Has Acquired
Mass Not Dissipate and Vanish?'', v2.0 (2026-10-10). Version DOI
10.5281/zenodo.23280070, Concept DOI 10.5281/zenodo.23266250. - {[}D2{]}
N. Kihara, ``The Eleventh Thought Experiment: Emission Records of
Two-Body Systems without a Background Spacetime'', v1.0 (2026-10-07).
Version DOI 10.5281/zenodo.23199300, Concept DOI
10.5281/zenodo.23199299.

\textbf{E The beginning and closed relational systems} - {[}E1{]} N.
Kihara, ``The Mechanism of Inflationary Rapid Expansion in
Self-Consistent Closed Systems of Relational Waves -- Second Edition''
(2026-08-30). Version DOI 10.5281/zenodo.22176949, Concept DOI
10.5281/zenodo.22112008. - {[}E2{]} N. Kihara, ``A Finite Complete
Relational Closed System with Time Vertexed -- the Minimal Exact Closed
System N=1, K=4 by U\^{}K=I and a Complete-Network Analysis of Odd and
Even Harmonics'', v1.2 (2026-09-15). Version DOI
10.5281/zenodo.22763330, Concept DOI 10.5281/zenodo.22763329.

\textbf{F Projection and curvature} - {[}F1{]} N. Kihara, ``The Twelfth
Thought Experiment: From the Fine-Structure Constant to the Observation
Map and Gauge Symmetry'', v1.0 (2026-10-08). Version DOI
10.5281/zenodo.23226259, Concept DOI 10.5281/zenodo.23226258. - {[}F2{]}
N. Kihara, ``Geometric Formulation of 4-Dimensional Space via Central
Projection'', v1.0 (2026-04-05). DOI 10.5281/zenodo.19427780, Concept
DOI 10.5281/zenodo.19427779. - {[}F3{]} N. Kihara, ``Schwarzschild--de
Sitter Exact Solution in the Central Projection Framework''
(2026-04-12). DOI 10.5281/zenodo.19538098, Concept DOI
10.5281/zenodo.19538097. - {[}F4{]} N. Kihara, ``Radial Projection:
Definition and Relation to Central Projection'', v3.3 (2026-06-06).
Version DOI 10.5281/zenodo.20567347, Concept DOI
10.5281/zenodo.20462569.

\textbf{G The research framework and the First to Sixth Thought
Experiments} - {[}G1{]} N. Kihara, ``Design of a Research Project
Exploring the Emergence of Spacetime, Particles and Interactions from
States and Relations'', v1.0 (2026-09-20). Version DOI
10.5281/zenodo.22851945, Concept DOI 10.5281/zenodo.22851944. - {[}G2{]}
N. Kihara, ``The First Thought Experiment: Tracing the Equivalence
Principle back to Two Nameless Degrees of Freedom'', v1.0 (2026-09-20).
Version DOI 10.5281/zenodo.22857953, Concept DOI
10.5281/zenodo.22857952. - {[}G3{]} N. Kihara, ``The Second Thought
Experiment: Deriving a Coulomb-Type Inverse-Square Force from Product
Readout and an Area Clock'', v1.2 (2026-09-21). Version DOI
10.5281/zenodo.22876596, Concept DOI 10.5281/zenodo.22867336. - {[}G4{]}
N. Kihara, ``The Third Thought Experiment: Reading Out the State of a--b
Two Bodies from the Orbits of a--b Two Celestial Bodies'', v1.5
(2026-09-23). Version DOI 10.5281/zenodo.22909661, Concept DOI
10.5281/zenodo.22909660. - {[}G5{]} N. Kihara, ``The Fourth Thought
Experiment: From a Fixed Action S to a State-Dependent Action S\_n'',
v1.0 (2026-09-23). Version DOI 10.5281/zenodo.22919788, Concept DOI
10.5281/zenodo.22919787. - {[}G6{]} N. Kihara, ``The Fifth Thought
Experiment: Was There No Hidden State?'', v2.1 (2026-09-26). Version DOI
10.5281/zenodo.22973087, Concept DOI 10.5281/zenodo.22973086. - {[}G7{]}
N. Kihara, ``The Sixth Thought Experiment: Can Two Kinds of Long-Range
Interaction Be Internalised into One Closed State Map?'', v1.1
(2026-09-26). Version DOI 10.5281/zenodo.22974632, Concept DOI
10.5281/zenodo.22974631. - {[}G8{]} N. Kihara, ``Supplement to the Sixth
Thought Experiment: Does the State Generation Change When the
Normalisation Is Switched from Mass to the Charge Unit?'', v1.0
(2026-09-27). Version DOI 10.5281/zenodo.22985300, Concept DOI
10.5281/zenodo.22985299.

\subsubsection{External references}\label{external-references}

\textbf{X1 Mathematical structures that select integers} -
{[}X1-Hurwitz{]} A. Hurwitz, ``Über die Composition der quadratischen
Formen von beliebig vielen Variablen,'' \emph{Nachr. Ges. Wiss.
Göttingen, Math.-Phys. Kl.}, 309--316 (1898). - {[}X1-Coxeter{]} H. S.
M. Coxeter, ``Integral Cayley numbers,'' \emph{Duke Math. J.}
\textbf{13}, 561--578 (1946). DOI 10.1215/S0012-7094-46-01347-6. -
{[}X1-Korkin--Zolotarev{]} A. Korkine and G. Zolotareff, ``Sur les
formes quadratiques,'' \emph{Math. Ann.} \textbf{6}, 366--389 (1873). -
{[}X1-Conway--Smith{]} J. H. Conway and D. A. Smith, \emph{On
Quaternions and Octonions} (A K Peters, 2003). - {[}X1-Viazovska{]} M.
S. Viazovska, ``The sphere packing problem in dimension 8,'' \emph{Ann.
of Math.} \textbf{185}, 991--1015 (2017). DOI
10.4007/annals.2017.185.3.7. - {[}X1-Bamberg{]} J. Bamberg, G. Cairns
and D. Kilminster, ``The crystallographic restriction, permutations, and
Goldbach's conjecture,'' \emph{Amer. Math. Monthly} \textbf{110},
202--209 (2003). DOI 10.1080/00029890.2003.11919956. -
{[}X1-Hardy--Wright{]} G. H. Hardy and E. M. Wright, \emph{An
Introduction to the Theory of Numbers}, 6th ed.~(Oxford University
Press, 2008). - {[}X1-Koide{]} Y. Koide, ``A new view of quark and
lepton mass hierarchy,'' \emph{Phys. Rev.~D} \textbf{28}, 252 (1983);
first in \emph{Lett. Nuovo Cimento} \textbf{34}, 201 (1982). -
{[}X1-Żenczykowski{]} P. Żenczykowski, ``A note on Koide's doubly
special parametrization of quark masses,'' \emph{Open Phys.}
\textbf{16}, 427--429 (2018). DOI 10.1515/phys-2018-0058.

\textbf{X2 Dynamics of pull-in and incommensurability} - {[}X2-Adler{]}
R. Adler, ``A study of locking phenomena in oscillators,'' \emph{Proc.
IRE} \textbf{34}, 351--357 (1946). - {[}X2-Bak{]} P. Bak, ``Commensurate
phases, incommensurate phases and the devil's staircase,''
\emph{Rep.~Prog. Phys.} \textbf{45}, 587 (1982). DOI
10.1088/0034-4885/45/6/001. - {[}X2-Jensen{]} M. H. Jensen, P. Bak and
T. Bohr, ``Transition to chaos by interaction of resonances in
dissipative systems. I. Circle maps,'' \emph{Phys. Rev.~A} \textbf{30},
1960--1969 (1984). DOI 10.1103/PhysRevA.30.1960. - {[}X2-Kuramoto{]} Y.
Kuramoto, ``Self-entrainment of a population of coupled non-linear
oscillators,'' \emph{Lecture Notes in Physics} \textbf{39}, 420--422
(1975). DOI 10.1007/BFb0013365.

\textbf{X3 Finite recurrence and the readout of time} -
{[}X3-Page--Wootters{]} D. N. Page and W. K. Wootters, ``Evolution
without evolution: Dynamics described by stationary observables,''
\emph{Phys. Rev.~D} \textbf{27}, 2885 (1983). DOI
10.1103/PhysRevD.27.2885. - {[}X3-Schwinger1960{]} J. Schwinger,
``Unitary operator bases,'' \emph{Proc. Natl. Acad. Sci. USA}
\textbf{46}, 570--579 (1960). - {[}X3-Gabor{]} D. Gabor, ``Theory of
communication. Part 1,'' \emph{J. IEE III} \textbf{93}, 429--441 (1946).
DOI 10.1049/ji-3-2.1946.0074.

\textbf{X4 Localisation, non-dissipation, mass} - {[}X4-deBroglie{]} L.
de Broglie, ``Waves and quanta,'' \emph{Nature} \textbf{112}, 540
(1923). DOI 10.1038/112540a0. - {[}X4-Haus{]} H. A. Haus, ``Mode-locking
of lasers,'' \emph{IEEE J. Sel. Top. Quantum Electron.} \textbf{6},
1173--1185 (2000). DOI 10.1109/2944.902165.

\textbf{X5 Measured values and the answer table of the standard theory}
- {[}X5-CODATA2022{]} P. J. Mohr, D. B. Newell, B. N. Taylor and E.
Tiesinga, ``CODATA recommended values of the fundamental physical
constants: 2022,'' \emph{Rev.~Mod. Phys.} \textbf{97}, 025002 (2025). -
{[}X5-Fan{]} X. Fan, T. G. Myers, B. A. D. Sukra and G. Gabrielse,
``Measurement of the electron magnetic moment,'' \emph{Phys. Rev.~Lett.}
\textbf{130}, 071801 (2023). - {[}X5-Morel{]} L. Morel, Z. Yao, P. Cladé
and S. Guellati-Khélifa, ``Determination of the fine-structure constant
with an accuracy of 81 parts per trillion,'' \emph{Nature} \textbf{588},
61--65 (2020). DOI 10.1038/s41586-020-2964-7. - {[}X5-Parker{]} R. H.
Parker, C. Yu, W. Zhong, B. Estey and H. Müller, ``Measurement of the
fine-structure constant as a test of the Standard Model,''
\emph{Science} \textbf{360}, 191--195 (2018). DOI
10.1126/science.aap7706. - {[}X5-Peil{]} S. Peil and G. Gabrielse,
``Observing the quantum limit of an electron cyclotron: QND measurements
of quantum jumps between Fock states,'' \emph{Phys. Rev.~Lett.}
\textbf{83}, 1287 (1999). DOI 10.1103/PhysRevLett.83.1287. -
{[}X5-Lamb--Retherford{]} W. E. Lamb and R. C. Retherford, ``Fine
structure of the hydrogen atom by a microwave method,'' \emph{Phys.
Rev.} \textbf{72}, 241--243 (1947). DOI 10.1103/PhysRev.72.241. -
{[}X5-Bethe{]} H. A. Bethe, ``The electromagnetic shift of energy
levels,'' \emph{Phys. Rev.} \textbf{72}, 339--341 (1947). DOI
10.1103/PhysRev.72.339. - {[}X5-Schwinger1948{]} J. Schwinger, ``On
quantum-electrodynamics and the magnetic moment of the electron,''
\emph{Phys. Rev.} \textbf{73}, 416--417 (1948). DOI
10.1103/PhysRev.73.416. - {[}X5-Eides{]} M. I. Eides, H. Grotch and V.
A. Shelyuto, ``Theory of light hydrogenlike atoms,'' \emph{Phys. Rep.}
\textbf{342}, 63 (2001). - {[}X5-Keshavarzi{]} A. Keshavarzi, D. Nomura
and T. Teubner, ``The muon g−2 and α(M\_Z²): a new data-based
analysis,'' \emph{Phys. Rev.~D} \textbf{101}, 014029 (2020). DOI
10.1103/PhysRevD.101.014029. - {[}X5-Roueff{]} E. Roueff et al., ``The
full infrared spectrum of molecular hydrogen,'' \emph{Astron.
Astrophys.} \textbf{630}, A58 (2019).

\textbf{X6 Absorbers, horizons, cosmology, self-energy} -
{[}X6-Wheeler--Feynman{]} J. A. Wheeler and R. P. Feynman, ``Interaction
with the absorber as the mechanism of radiation,'' \emph{Rev.~Mod.
Phys.} \textbf{17}, 157 (1945). - {[}X6-Hoyle--Narlikar{]} F. Hoyle and
J. V. Narlikar, ``Time symmetric electrodynamics and the arrow of time
in cosmology,'' \emph{Proc. R. Soc. Lond. A} \textbf{277}, 1 (1964). -
{[}X6-Gibbons--Hawking{]} G. W. Gibbons and S. W. Hawking,
``Cosmological event horizons, thermodynamics, and particle creation,''
\emph{Phys. Rev.~D} \textbf{15}, 2738--2751 (1977). DOI
10.1103/PhysRevD.15.2738. - {[}X6-Bekenstein{]} J. D. Bekenstein,
``Black holes and entropy,'' \emph{Phys. Rev.~D} \textbf{7}, 2333
(1973). - {[}X6-Tangherlini{]} F. R. Tangherlini, ``Schwarzschild field
in n dimensions and the dimensionality of space problem,'' \emph{Nuovo
Cimento} \textbf{27}, 636 (1963). - {[}X6-ADM1960{]} R. Arnowitt, S.
Deser and C. W. Misner, ``Gravitational-electromagnetic coupling and the
classical self-energy problem,'' \emph{Phys. Rev.} \textbf{120},
313--320 (1960). - {[}X6-ISS{]} C. J. Isham, A. Salam and J. Strathdee,
``Infinity suppression in gravity-modified quantum electrodynamics,''
\emph{Phys. Rev.~D} \textbf{3}, 1805--1817 (1971). DOI
10.1103/PhysRevD.3.1805. - {[}X6-Planck2018VI{]} Planck Collaboration,
``Planck 2018 results. VI. Cosmological parameters,'' \emph{Astron.
Astrophys.} \textbf{641}, A6 (2020). - {[}X6-PDG-BBN{]} B. D. Fields, P.
Molaro and S. Sarkar, ``Big-Bang nucleosynthesis,'' in Particle Data
Group, \emph{Review of Particle Physics}, \emph{Phys. Rev.~D}
\textbf{110}, 030001 (2024). - {[}X6-Schwarz{]} D. J. Schwarz, C. J.
Copi, D. Huterer and G. D. Starkman, ``CMB anomalies after Planck,''
\emph{Class. Quantum Grav.} \textbf{33}, 184001 (2016). -
{[}X6-Bezrukov--Shaposhnikov{]} F. Bezrukov and M. Shaposhnikov, ``The
Standard Model Higgs boson as the inflaton,'' \emph{Phys. Lett. B}
\textbf{659}, 703 (2008).

\textbf{X7 Particle symmetries and the physics of E₈} - {[}X7-Furey{]}
C. Furey, ``SU(3)\_C × SU(2)\_L × U(1)\_Y (× U(1)\_X) as a symmetry of
division algebraic ladder operators,'' \emph{Eur. Phys. J. C}
\textbf{78}, 375 (2018). - {[}X7-Günaydin--Gürsey{]} M. Günaydin and F.
Gürsey, ``Quark structure and octonions,'' \emph{J. Math. Phys.}
\textbf{14}, 1651--1667 (1973). - {[}X7-Pati--Salam{]} J. C. Pati and A.
Salam, ``Lepton number as the fourth `color','' \emph{Phys. Rev.~D}
\textbf{10}, 275--289 (1974). DOI 10.1103/PhysRevD.10.275. -
{[}X7-Ma--Rajasekaran{]} E. Ma and G. Rajasekaran, ``Softly broken A4
symmetry for nearly degenerate neutrino masses,'' \emph{Phys. Rev.~D}
\textbf{64}, 113012 (2001). DOI 10.1103/PhysRevD.64.113012. -
{[}X7-Lisi{]} A. G. Lisi, ``An exceptionally simple theory of
everything,'' arXiv:0711.0770 (2007). - {[}X7-Distler--Garibaldi{]} J.
Distler and S. Garibaldi, ``There is no `Theory of Everything' inside
E8,'' \emph{Commun. Math. Phys.} \textbf{298}, 419--436 (2010). DOI
10.1007/s00220-010-1006-y. - {[}X7-'t Hooft{]} G. 't Hooft, ``On the
phase transition towards permanent quark confinement,'' \emph{Nucl.
Phys. B} \textbf{138}, 1 (1978). - {[}X7-Zamolodchikov{]} A. B.
Zamolodchikov, ``Integrals of motion and S-matrix of the (scaled) T=T\_c
Ising model with magnetic field,'' \emph{Int. J. Mod. Phys. A}
\textbf{4}, 4235--4248 (1989). - {[}X7-Coldea{]} R. Coldea et al.,
``Quantum criticality in an Ising chain: Experimental evidence for
emergent E8 symmetry,'' \emph{Science} \textbf{327}, 177--180 (2010).
DOI 10.1126/science.1180085.

\textbf{X8 Close prior work raised in Chapter 8 (surveyed within the
dialogue; volume and page details are to be confirmed by the author
before citation elsewhere)} - {[}X8-Shanahan{]} D. Shanahan, ``A Case
for Lorentzian Relativity,'' \emph{Found. Phys.} \textbf{44} (2014). -
{[}X8-Feynman--Hibbs{]} R. P. Feynman and A. R. Hibbs, \emph{Quantum
Mechanics and Path Integrals} (McGraw-Hill, 1965): the checkerboard
model. - {[}X8-Jacobson--Schulman{]} T. Jacobson and L. S. Schulman,
\emph{J. Phys. A} \textbf{17}, 375 (1984): mass as the frequency of path
reversals. - {[}X8-Bondi{]} H. Bondi, the k-calculus (\emph{Relativity
and Common Sense}, 1964). - {[}X8-Adams{]} J. F. Adams, ``On the
non-existence of elements of Hopf invariant one,'' \emph{Ann. of Math.}
\textbf{72}, 20--104 (1960): the parallelisable spheres are S¹, S³, S⁷.
- {[}X8-Hestenes{]} D. Hestenes, ``The zitterbewegung interpretation of
quantum mechanics,'' \emph{Found. Phys.} \textbf{20}, 1213--1232 (1990).

\begin{center}\rule{0.5\linewidth}{0.5pt}\end{center}

\subsection{Appendix A　Reproduction and the interactive
figures}\label{appendix-a-reproduction-and-the-interactive-figures}

Everything is in GitHub \texttt{WurabeSeiji/\allowbreak{}ai-chat-logs-open} (the
Zenodo record \url{https://doi.org/10.5281/zenodo.23384929} holds the Japanese
and English md, tex and PDF, the seven interactive HTML figures,
\texttt{build\_\allowbreak{}tex.py} and \texttt{SHA256SUMS.txt}). The folder of this
paper is
\texttt{匿名頂点状態生成幾何-匿名内部観測者から不変な関係量の体系/\allowbreak{}第十四思考実験\_\allowbreak{}整数を選ぶ対称性の非整合と発展方向の読出し\_\allowbreak{}20261012/\allowbreak{}統合仮説論文\_\allowbreak{}20261015/\allowbreak{}}.
The full texts of the primary-source dialogues are
\texttt{chat\_\allowbreak{}logs/\allowbreak{}01\_\allowbreak{}円周上の内在的な幾何学\_\allowbreak{}20261013-14\_\allowbreak{}full.md} and
\texttt{chat\_\allowbreak{}logs/\allowbreak{}02\_\allowbreak{}中性子の寿命と弱い力によるフレーバー変化\_\allowbreak{}20261014-15\_\allowbreak{}full.md}.
The correspondence table between each stage and the passages of the
dialogue, the sections of the paper and the experiment folders is §4 of
\texttt{working\_\allowbreak{}notes.md}.

\textbf{Numerical experiments of the sphere residents
(\texttt{次元の生成構造/\allowbreak{}})}

{\def\LTcaptype{none} % do not increment counter
\begin{longtable}[]{@{}
  >{\raggedright\arraybackslash}p{(\linewidth - 8\tabcolsep) * \real{0.2000}}
  >{\raggedright\arraybackslash}p{(\linewidth - 8\tabcolsep) * \real{0.2000}}
  >{\raggedright\arraybackslash}p{(\linewidth - 8\tabcolsep) * \real{0.2000}}
  >{\raggedright\arraybackslash}p{(\linewidth - 8\tabcolsep) * \real{0.2000}}
  >{\raggedright\arraybackslash}p{(\linewidth - 8\tabcolsep) * \real{0.2000}}@{}}
\toprule\noalign{}
\begin{minipage}[b]{\linewidth}\raggedright
Content
\end{minipage} & \begin{minipage}[b]{\linewidth}\raggedright
Folder
\end{minipage} & \begin{minipage}[b]{\linewidth}\raggedright
Computation and plotting
\end{minipage} & \begin{minipage}[b]{\linewidth}\raggedright
Results
\end{minipage} & \begin{minipage}[b]{\linewidth}\raggedright
Interactive figures (HTML)
\end{minipage} \\
\midrule\noalign{}
\endhead
\bottomrule\noalign{}
\endlastfoot
Self-consistent shell (§4.1, results not adopted) &
\texttt{円周住民\_\allowbreak{}自己無撞着殻\_\allowbreak{}20261013},
\texttt{四次元球殻住民\_\allowbreak{}自己無撞着殻\_\allowbreak{}20261013} &
\texttt{selfconsistent\_\allowbreak{}shell\_\allowbreak{}s1.py}, \texttt{plot\_\allowbreak{}…\_\allowbreak{}s1.py},
\texttt{…\_\allowbreak{}s3.py} &
\texttt{results\_\allowbreak{}summary\_\allowbreak{}selfconsistent\_\allowbreak{}shell\_\allowbreak{}s1.json},
\texttt{…\_\allowbreak{}s3.json} & --- \\
Gauge interval and geodesic length (Chapter 5) &
\texttt{円周住民\_\allowbreak{}ゲージ間隔測地線長\_\allowbreak{}20261014} &
\texttt{gauge\_\allowbreak{}geodesic\_\allowbreak{}s1.py}, \texttt{plot\_\allowbreak{}gauge\_\allowbreak{}geodesic\_\allowbreak{}s1.py}
(generates PNG and SVG of \texttt{fig1\_\allowbreak{}R\_\allowbreak{}vs\_\allowbreak{}phase},
\texttt{fig2\_\allowbreak{}gauge\_\allowbreak{}interval\_\allowbreak{}lengths}, \texttt{fig3\_\allowbreak{}resident\_\allowbreak{}ruler})
& \texttt{results\_\allowbreak{}summary\_\allowbreak{}gauge\_\allowbreak{}geodesic\_\allowbreak{}s1.json} &
\texttt{circle\_\allowbreak{}imag\_\allowbreak{}amplitude\_\allowbreak{}radial\_\allowbreak{}log\_\allowbreak{}gradient.html} \\
Gauge widening versus dispersion ratio (§7.6, does not hold) &
\texttt{ゲージ広がりと分散比の検証\_\allowbreak{}20261014} &
\texttt{gauge\_\allowbreak{}stretch\_\allowbreak{}vs\_\allowbreak{}dispersion.py},
\texttt{plot\_\allowbreak{}gauge\_\allowbreak{}stretch\_\allowbreak{}vs\_\allowbreak{}dispersion.py}
(\texttt{fig\_\allowbreak{}gauge\_\allowbreak{}stretch\_\allowbreak{}vs\_\allowbreak{}dispersion}) &
\texttt{results\_\allowbreak{}summary\_\allowbreak{}gauge\_\allowbreak{}stretch\_\allowbreak{}vs\_\allowbreak{}dispersion.json} & --- \\
Light-cone readout (§7.9), action of phase rotation (§8.1), observer
(Chapter 10) & \texttt{円周住民\_\allowbreak{}光円錐読出し\_\allowbreak{}20261014} &
\texttt{lightcone\_\allowbreak{}readout\_\allowbreak{}s1.py},
\texttt{plot\_\allowbreak{}lightcone\_\allowbreak{}readout\_\allowbreak{}s1.py}
(\texttt{fig\_\allowbreak{}lightcone\_\allowbreak{}readout\_\allowbreak{}s1}) &
\texttt{results\_\allowbreak{}summary\_\allowbreak{}lightcone\_\allowbreak{}readout\_\allowbreak{}s1.json} &
\texttt{circle\_\allowbreak{}phase\_\allowbreak{}rotation\_\allowbreak{}action\_\allowbreak{}lightcone.html},
\texttt{circle\_\allowbreak{}phase\_\allowbreak{}rotation\_\allowbreak{}action\_\allowbreak{}lightcone\_\allowbreak{}v2.html},
\texttt{sr\_\allowbreak{}observer\_\allowbreak{}lightcone\_\allowbreak{}s1.html} \\
Slices of S³ (Chapter 6), action on S³ (§8.2), observer on S³ (Chapter
10) & \texttt{四次元球殻住民\_\allowbreak{}輪切り可視化\_\allowbreak{}20261014} & --- & --- &
\texttt{s3\_\allowbreak{}shell\_\allowbreak{}slice\_\allowbreak{}amplitude\_\allowbreak{}gradient.html},
\texttt{s3\_\allowbreak{}phase\_\allowbreak{}rotation\_\allowbreak{}action\_\allowbreak{}lightcone.html},
\texttt{sr\_\allowbreak{}s3\_\allowbreak{}observer\_\allowbreak{}lightcone.html} \\
\end{longtable}
}

The HTML files are interactive figures that open in a browser on their
own; they are included in this Zenodo record (Version DOI
10.5281/zenodo.23384929). Download and open in a browser. The
direct-download URLs are as follows (the GitHub raw links display the
source; save and open).

{\def\LTcaptype{none} % do not increment counter
\begin{longtable}[]{@{}
  >{\raggedright\arraybackslash}p{(\linewidth - 6\tabcolsep) * \real{0.2500}}
  >{\raggedright\arraybackslash}p{(\linewidth - 6\tabcolsep) * \real{0.2500}}
  >{\raggedright\arraybackslash}p{(\linewidth - 6\tabcolsep) * \real{0.2500}}
  >{\raggedright\arraybackslash}p{(\linewidth - 6\tabcolsep) * \real{0.2500}}@{}}
\toprule\noalign{}
\begin{minipage}[b]{\linewidth}\raggedright
Interactive figure
\end{minipage} & \begin{minipage}[b]{\linewidth}\raggedright
Content
\end{minipage} & \begin{minipage}[b]{\linewidth}\raggedright
Zenodo (direct download)
\end{minipage} & \begin{minipage}[b]{\linewidth}\raggedright
GitHub (raw)
\end{minipage} \\
\midrule\noalign{}
\endhead
\bottomrule\noalign{}
\endlastfoot
\texttt{circle\_\allowbreak{}imag\_\allowbreak{}amplitude\_\allowbreak{}radial\_\allowbreak{}log\_\allowbreak{}gradient.html} & Circle:
imaginary amplitude (vertical) and slope ∂θ ln & R² & (radial) (§5.3) \\
\texttt{circle\_\allowbreak{}phase\_\allowbreak{}rotation\_\allowbreak{}action\_\allowbreak{}lightcone.html} & Circle:
action of phase rotation and light-cone distribution (first version,
§8.1) &
\href{https://zenodo.org/records/23384929/files/circle_phase_rotation_action_lightcone.html?download=1}{download}
&
\href{https://raw.githubusercontent.com/WurabeSeiji/ai-chat-logs-open/main/\%E6\%AC\%A1\%E5\%85\%83\%E3\%81\%AE\%E7\%94\%9F\%E6\%88\%90\%E6\%A7\%8B\%E9\%80\%A0/\%E5\%86\%86\%E5\%91\%A8\%E4\%BD\%8F\%E6\%B0\%91_\%E5\%85\%89\%E5\%86\%86\%E9\%8C\%90\%E8\%AA\%AD\%E5\%87\%BA\%E3\%81\%97_20261014/circle_phase_rotation_action_lightcone.html}{raw} \\
\texttt{circle\_\allowbreak{}phase\_\allowbreak{}rotation\_\allowbreak{}action\_\allowbreak{}lightcone\_\allowbreak{}v2.html} & Circle:
action of phase rotation and light-cone distribution
(translated-distance slider, separate odd/even N, §8.1) &
\href{https://zenodo.org/records/23384929/files/circle_phase_rotation_action_lightcone_v2.html?download=1}{download}
&
\href{https://raw.githubusercontent.com/WurabeSeiji/ai-chat-logs-open/main/\%E6\%AC\%A1\%E5\%85\%83\%E3\%81\%AE\%E7\%94\%9F\%E6\%88\%90\%E6\%A7\%8B\%E9\%80\%A0/\%E5\%86\%86\%E5\%91\%A8\%E4\%BD\%8F\%E6\%B0\%91_\%E5\%85\%89\%E5\%86\%86\%E9\%8C\%90\%E8\%AA\%AD\%E5\%87\%BA\%E3\%81\%97_20261014/circle_phase_rotation_action_lightcone_v2.html}{raw} \\
\texttt{sr\_\allowbreak{}observer\_\allowbreak{}lightcone\_\allowbreak{}s1.html} & Circle residents: the
special-relativity thought experiment (observer and particle, Chapter
10) &
\href{https://zenodo.org/records/23384929/files/sr_observer_lightcone_s1.html?download=1}{download}
&
\href{https://raw.githubusercontent.com/WurabeSeiji/ai-chat-logs-open/main/\%E6\%AC\%A1\%E5\%85\%83\%E3\%81\%AE\%E7\%94\%9F\%E6\%88\%90\%E6\%A7\%8B\%E9\%80\%A0/\%E5\%86\%86\%E5\%91\%A8\%E4\%BD\%8F\%E6\%B0\%91_\%E5\%85\%89\%E5\%86\%86\%E9\%8C\%90\%E8\%AA\%AD\%E5\%87\%BA\%E3\%81\%97_20261014/sr_observer_lightcone_s1.html}{raw} \\
\texttt{s3\_\allowbreak{}shell\_\allowbreak{}slice\_\allowbreak{}amplitude\_\allowbreak{}gradient.html} & Slices of the
four-dimensional shell S³: amplitude (shading, opacity) and slope ∂ψ ln
& R² & (Chapter 6) \\
\texttt{s3\_\allowbreak{}phase\_\allowbreak{}rotation\_\allowbreak{}action\_\allowbreak{}lightcone.html} & Four-dimensional
shell S³: action of phase rotation and light-cone distribution (§8.2) &
\href{https://zenodo.org/records/23384929/files/s3_phase_rotation_action_lightcone.html?download=1}{download}
&
\href{https://raw.githubusercontent.com/WurabeSeiji/ai-chat-logs-open/main/\%E6\%AC\%A1\%E5\%85\%83\%E3\%81\%AE\%E7\%94\%9F\%E6\%88\%90\%E6\%A7\%8B\%E9\%80\%A0/\%E5\%9B\%9B\%E6\%AC\%A1\%E5\%85\%83\%E7\%90\%83\%E6\%AE\%BB\%E4\%BD\%8F\%E6\%B0\%91_\%E8\%BC\%AA\%E5\%88\%87\%E3\%82\%8A\%E5\%8F\%AF\%E8\%A6\%96\%E5\%8C\%96_20261014/s3_phase_rotation_action_lightcone.html}{raw} \\
\texttt{sr\_\allowbreak{}s3\_\allowbreak{}observer\_\allowbreak{}lightcone.html} & Four-dimensional shell S³:
the special-relativity thought experiment with an observer (Chapter 10)
&
\href{https://zenodo.org/records/23384929/files/sr_s3_observer_lightcone.html?download=1}{download}
&
\href{https://raw.githubusercontent.com/WurabeSeiji/ai-chat-logs-open/main/\%E6\%AC\%A1\%E5\%85\%83\%E3\%81\%AE\%E7\%94\%9F\%E6\%88\%90\%E6\%A7\%8B\%E9\%80\%A0/\%E5\%9B\%9B\%E6\%AC\%A1\%E5\%85\%83\%E7\%90\%83\%E6\%AE\%BB\%E4\%BD\%8F\%E6\%B0\%91_\%E8\%BC\%AA\%E5\%88\%87\%E3\%82\%8A\%E5\%8F\%AF\%E8\%A6\%96\%E5\%8C\%96_20261014/sr_s3_observer_lightcone.html}{raw} \\
\end{longtable}
}

The PNG and SVG files of the plotting programs were generated on
2026-10-15 by re-running the computation programs unchanged in each
folder and then the plotting programs (Figs. 1--5). The re-run results
agree with the saved summary JSON (the gauge-interval ratios agree to
four decimals in every system, the s versus ω/k table agrees in every
row, the light-cone residual is 1.75×10⁻¹⁵, the v − tanh difference
7.77×10⁻¹⁶, and the five velocity-addition cases agree to six digits).
The font Noto Sans CJK JP specified by the plotting programs was
installed on this Mac by the author's decision
(\texttt{\textasciitilde{}/Library/Fonts/NotoSansCJKjp-Regular.otf});
the programs were not modified. Copies of the PNG files of Figs. 1--10
are in \texttt{figures/\allowbreak{}} of this paper.

\textbf{Numerical experiments of the Fourteenth Thought Experiment
(parent folder of this paper)}

{\def\LTcaptype{none} % do not increment counter
\begin{longtable}[]{@{}
  >{\raggedright\arraybackslash}p{(\linewidth - 6\tabcolsep) * \real{0.2500}}
  >{\raggedright\arraybackslash}p{(\linewidth - 6\tabcolsep) * \real{0.2500}}
  >{\raggedright\arraybackslash}p{(\linewidth - 6\tabcolsep) * \real{0.2500}}
  >{\raggedright\arraybackslash}p{(\linewidth - 6\tabcolsep) * \real{0.2500}}@{}}
\toprule\noalign{}
\begin{minipage}[b]{\linewidth}\raggedright
Experiment
\end{minipage} & \begin{minipage}[b]{\linewidth}\raggedright
Folder
\end{minipage} & \begin{minipage}[b]{\linewidth}\raggedright
Run
\end{minipage} & \begin{minipage}[b]{\linewidth}\raggedright
Figure
\end{minipage} \\
\midrule\noalign{}
\endhead
\bottomrule\noalign{}
\endlastfoot
E0, E1 & \texttt{数値実験E1\_\allowbreak{}レプトン三相位相の格子引き込み\_\allowbreak{}20261014/\allowbreak{}}
& \texttt{te14\_\allowbreak{}E0\_\allowbreak{}koide\_\allowbreak{}phase\_\allowbreak{}check.py},
\texttt{te14\_\allowbreak{}E1\_\allowbreak{}lepton\_\allowbreak{}phase\_\allowbreak{}lattice.py}, \texttt{te14\_\allowbreak{}E1\_\allowbreak{}plot.py}
& \texttt{te14\_\allowbreak{}E1\_\allowbreak{}fig.png}, \texttt{.svg} \\
E2 & \texttt{数値実験E2\_\allowbreak{}四次元フルヴィッツ格子と23\_\allowbreak{}124\_\allowbreak{}20261014/\allowbreak{}} &
\texttt{te14\_\allowbreak{}E2\_\allowbreak{}hurwitz\_\allowbreak{}124\_\allowbreak{}23.py},
\texttt{te14\_\allowbreak{}E2\_\allowbreak{}radius\_\allowbreak{}scan.py}, \texttt{te14\_\allowbreak{}E2\_\allowbreak{}plot.py} &
\texttt{te14\_\allowbreak{}E2\_\allowbreak{}fig.png}, \texttt{.svg} (the readout in panel (b) uses
the withdrawn measurement and is not used here) \\
E3 & \texttt{数値実験E3\_\allowbreak{}基底状態と準位系列\_\allowbreak{}20261014/\allowbreak{}} &
\texttt{te14\_\allowbreak{}E3\_\allowbreak{}ground\_\allowbreak{}and\_\allowbreak{}levels.py}, \texttt{te14\_\allowbreak{}E3\_\allowbreak{}plot.py}
(with the E2 programs in the same folder) & \texttt{te14\_\allowbreak{}E3\_\allowbreak{}fig.png},
\texttt{.svg} \\
E4 & \texttt{数値実験E4\_\allowbreak{}選択・引き込み・回復\_\allowbreak{}20261014/\allowbreak{}} &
\texttt{te14\_\allowbreak{}E4\_\allowbreak{}selection\_\allowbreak{}pullin\_\allowbreak{}recovery.py},
\texttt{te14\_\allowbreak{}E4\_\allowbreak{}mod4\_\allowbreak{}check.py}, \texttt{te14\_\allowbreak{}E4\_\allowbreak{}plot.py} &
\texttt{te14\_\allowbreak{}E4\_\allowbreak{}fig.png}, \texttt{.svg} \\
E5 & \texttt{数値実験E5\_\allowbreak{}ゲージの自己歪み\_\allowbreak{}20261014/\allowbreak{}} &
\texttt{te14\_\allowbreak{}E5\_\allowbreak{}gauge\_\allowbreak{}backreaction.py},
\texttt{te14\_\allowbreak{}E5\_\allowbreak{}control.py}, \texttt{te14\_\allowbreak{}E5\_\allowbreak{}plot.py} &
\texttt{te14\_\allowbreak{}E5\_\allowbreak{}fig.png}, \texttt{.svg} \\
Harmonic selection rule & parent folder &
\texttt{te14\_\allowbreak{}harmonic\_\allowbreak{}selection\_\allowbreak{}check.py} & --- \\
137 with balls & parent folder &
\texttt{te14\_\allowbreak{}ball\_\allowbreak{}count\_\allowbreak{}in\_\allowbreak{}radius3.py} (argument r) & --- \\
E₈ gluing & parent folder & \texttt{te14\_\allowbreak{}e8\_\allowbreak{}hurwitz\_\allowbreak{}pair\_\allowbreak{}glue.py} &
--- \\
\end{longtable}
}

On 2026-10-15 E0--E5 were re-run with the original programs unchanged
and checked against the original summaries (the summaries of E1--E3
agree completely; in E4 the re-lock targets and transient lengths agree
completely and only the last digits of the mean rotation number during
transients differ; E5 reproduces the values in its README). Environment:
Python 3.9, numpy 2.0.2, matplotlib 3.9.4.

\subsection{Appendix B　Record of
corrections}\label{appendix-b-record-of-corrections}

Claude's errors and the author's corrections that this paper kept in the
text as part of the derivation.

{\def\LTcaptype{none} % do not increment counter
\begin{longtable}[]{@{}
  >{\raggedright\arraybackslash}p{(\linewidth - 4\tabcolsep) * \real{0.3333}}
  >{\raggedright\arraybackslash}p{(\linewidth - 4\tabcolsep) * \real{0.3333}}
  >{\raggedright\arraybackslash}p{(\linewidth - 4\tabcolsep) * \real{0.3333}}@{}}
\toprule\noalign{}
\begin{minipage}[b]{\linewidth}\raggedright
Place
\end{minipage} & \begin{minipage}[b]{\linewidth}\raggedright
Claude's error
\end{minipage} & \begin{minipage}[b]{\linewidth}\raggedright
Author's correction
\end{minipage} \\
\midrule\noalign{}
\endhead
\bottomrule\noalign{}
\endlastfoot
§2.3 & Took the points of aligned phase as nodes. Restricted to
e\^{}\{iks\}. Introduced C cos(mθ+δ) & Nodes are the phases where the
amplitude vanishes. There is only θ and the imaginary amplitude A \\
§2.4 & n=0 is a state in which nothing appears & The two states ±iA \\
§2.5 & Drew the imaginary axis radially & The radial direction is that
of a background space; not allowed \\
§2.8 & Left the amplitude as a coefficient. Mixed the imaginary
amplitude into the phase. Skipped the Σ expansion & Amplitude into the
exponent too; the imaginary amplitude as an independent factor;
expand \\
§2.10 & θ₀ needs a second wave & The arrangement of the component phases
is θ₀ itself. The fundamental is the gauge \\
§4.1--4.2 & \textbar xₙ\textbar², a real radius, a centroid,
\textbar Σxₙ\textbar² & True squares. Not a norm. There is no outside of
the circle \\
§7.2 & The dispersion √(k²+m²) is an assumption not derived from the
axiom & It is the light cone Σxₙ² = 0 itself (§7.8) \\
§7.4 & The internal clock oscillates & A phase pinned to space that does
not oscillate. Exactly the Higgs field \\
§7.6 & The phase velocity exceeds the speed of light & The spatial gauge
is widened \\
§8.1 & A naive phase advance is a problem. The distribution is not a
physical proof & Shifting θ₀ is enough. Action and distribution close
within the axiom. The rotation speed is unnecessary \\
§8.5 & What if a, b are given phases & a, b are real. z = a + ib \\
Chapter 10 & β ticked from outside. Several localisations needed & β is
the observer's wave. The observer is a resident of the circle; no
god's-eye view \\
§11.4 & Odd alone is not localised & Odd alone is localised at two
points \\
§13.1 & The observer accelerated by a slider, R computed separately &
Observation is the cross term. No acceleration without radiation. The
apparatus too changes R \\
Chapter 15 & Electron = complete lock; the argument that metastability
needs 23/124 & No ``why'' needed. The problem is why it locks to that
window \\
Chapters 15--16 & The design of E2 was facile & The diagonal of the
four-dimensional lattice is an integer. Why the background is
four-dimensional \\
§17.3 & h fixed at its SI value, ``2000 times off'' & In G = c = q₀ = 1,
ħ = 9α⁻¹ is an output. Withdrawn \\
§18.2 & Made the electron mass a separate test & If α is 137 it
necessarily agrees. Only α is tested \\
§19.2 & Did not notice the reason for the deviation & Observation is an
average over many samples; transitions are observed too \\
§19.4 & The gauge was fixed & The electron itself distorts the gauge \\
§19.6 & A classical-probability average & The interference term of the
expectation value is proportional to the amplitude \\
§20.1 & At what radius does the electron live & There is no absolute
scale. Everything is a ratio \\
§20.1 & The selection rule in all 33 orbits & 29 (4 lack the
symmetry) \\
§21.1--21.2 & The gauge side needs a delay. R is an input and a dynamics
is needed & The partner wave holds and returns. The pull-in of α is the
mechanism that makes R 0.7. Dynamics is inherent as the absorbing
shell \\
§23.3 & Made the electron mass a separate test as a ratio to the Stoney
mass & A circular argument. The electron mass from α, α from the
structure of space \\
§24.4 & The void goes as R³, not R² & Leading term R³, second term
−6πR². Used in BH thermodynamics \\
§24.7 & Could not find the earlier computation of the photon--baryon
ratio & The fraction of integer ratios for resolution N, N = 10⁶⁰ \\
§24.8 & Left the orders of magnitude of the cosmological constant
unresolved & Already derived as Λ = 3/R² \\
E2→E3 & Measured the rotation number as the mean per-step angle &
Compute from the (integer) winding numbers \\
Sphere-resident paper v1 & 0.972 in the table of §6; conflation of
Adams' theorem with group composition & 0.988; distinguished \\
\end{longtable}
}

\subsection{Version history}\label{version-history}

\begin{itemize}
\tightlist
\item
  \textbf{v0.1 --- 2026-10-15}: Written following claude.ai's
  instruction v1, in a 25-chapter structure with the self-consistent
  lock α = 23/124 as the central claim. Rejected by the author (α is one
  stage of the development, not the subject; the development from
  residents of the circle to residents of the four-dimensional shell,
  and the most important result---that time, space, mass, charge,
  translational velocity and spin can be read out from the relations and
  interference of complex waves alone---had been dropped). Kept as a
  record in \texttt{unified\_\allowbreak{}hypothesis\_\allowbreak{}paper\_\allowbreak{}ja\_\allowbreak{}v0.1.md}.
\item
  \textbf{v0.2 --- 2026-10-15}: Following instruction v3, rewritten in
  full with the full texts of the two dialogues as primary sources, in
  the exact order of the author's questions and corrections. The chapter
  order follows the correspondence table in §4 of
  \texttt{working\_\allowbreak{}notes.md}.
\item
  \textbf{v1.0 --- 2026-10-15}: Published version. By the author's
  decision the quotations were tidied (typographical errors, colloquial
  interjections and expressions of reproach removed; meaning unchanged),
  Figs. 1--10 embedded, and the figures of the sphere-resident
  experiments generated by re-running. Registered on Zenodo with
  Japanese and English md, tex and PDF and the seven interactive HTML
  files (Version DOI 10.5281/zenodo.23384929, Concept DOI
  10.5281/zenodo.23384928). The claims of this version are not rewritten
  retroactively when later numerical experiments appear; differences are
  recorded from v1.1 onwards.
\end{itemize}

\subsection{Record of AI involvement}\label{record-of-ai-involvement}

The primary sources of this paper are two dialogues between the author
and Claude (claude.ai). Claude's errors in the dialogues and the
author's corrections are kept in the text as part of the derivation. The
writing was done by Claude Code; the structure follows the order of the
author's eleven stages and of the dialogues, and Claude Code did not
create an order by its own inference. The re-running and checking of the
numerical experiments E0--E5 were done by Claude Code. Every statement
in the text rests on the primary sources and the texts of the earlier
papers. The English version is a translation of the Japanese v1.0 by
Claude Code.

\end{document}
